% ==============================================================================
%  Apuntes de Regresión y ANOVA (RANO) — bloque de REGRESIÓN (Temas 1–6)
%  Archivo principal: compila este archivo con latexmk.
%
%  Estructura prevista (ampliable):
%    - bloque_regresion.tex : lista de temas de regresión (Temas 1–6).
%    - más adelante, bloque_anova.tex (Temas 7…) con su propio documento
%      (índice «Índice — ANOVA») y un documento global que incluya ambos bloques.
% ==============================================================================
\documentclass[11pt,a4paper,oneside]{report}

% ==============================================================================
%  Paquetes y configuración general
% ==============================================================================

% --- Idioma y fuentes ----------------------------------------------------------
\usepackage[T1]{fontenc}
\usepackage[utf8]{inputenc}
\usepackage{lmodern}
\usepackage[spanish,es-nodecimaldot,es-tabla,es-noshorthands]{babel}
\usepackage{microtype}

% --- Matemáticas ---------------------------------------------------------------
\usepackage{amsmath, amssymb, amsthm, mathtools}
\usepackage{bm}

% --- Página --------------------------------------------------------------------
\usepackage[a4paper, top=2.5cm, bottom=2.5cm, left=2.7cm, right=2.7cm,
            headheight=14pt]{geometry}
\usepackage{setspace}
\setstretch{1.15}
\setlength{\parindent}{0pt}
\setlength{\parskip}{0.5em}

% --- Tablas, listas, figuras ---------------------------------------------------
\usepackage{booktabs}
\usepackage{array}
\usepackage{multirow}
\usepackage{enumitem}
\setlist{itemsep=0.2em, topsep=0.3em}
\usepackage{graphicx}
\graphicspath{{figuras/}}
\usepackage{float}
\usepackage{caption}
\usepackage{subcaption}

% --- Gráficos ------------------------------------------------------------------
\usepackage{tikz}
\usetikzlibrary{arrows.meta, positioning, shapes, calc, decorations.pathreplacing, babel}
\usepackage{pgfplots}
\pgfplotsset{compat=1.18}
\usepackage{pgfplotstable}

% --- Colores y cajas -----------------------------------------------------------
\usepackage[dvipsnames]{xcolor}
\usepackage[most]{tcolorbox}

% --- Encabezados y títulos ----------------------------------------------------
\usepackage{titlesec}
\usepackage{fancyhdr}

% --- Enlaces (cargar casi al final) ---------------------------------------------
\definecolor{enlace}{RGB}{50,75,120}
\usepackage[colorlinks=true, linkcolor=enlace, urlcolor=enlace,
            citecolor=enlace, bookmarksnumbered=true]{hyperref}
\usepackage{bookmark}
\usepackage[spanish, nameinlink, noabbrev]{cleveref}

% --- Estilo de capítulos (un capítulo = un tema) -------------------------------
\titleformat{\chapter}[display]
  {\normalfont\bfseries}
  {\LARGE Tema \thechapter}
  {0.4em}
  {\Huge}
  [{\vspace{0.3em}\titlerule[1pt]}]
\titlespacing*{\chapter}{0pt}{-20pt}{25pt}

\titleformat{\section}
  {\normalfont\Large\bfseries}{\thesection}{0.8em}{}
\titleformat{\subsection}
  {\normalfont\large\bfseries}{\thesubsection}{0.8em}{}

% --- Encabezado / pie ----------------------------------------------------------
\pagestyle{fancy}
\fancyhf{}
\renewcommand{\chaptermark}[1]{\markboth{Tema \thechapter. #1}{}\gdef\tituloTemaActual{#1}}
\fancyhead[L]{\small\nouppercase{\leftmark}}
\fancyhead[R]{\small Regresión y ANOVA}
\fancyfoot[C]{\thepage}
\renewcommand{\headrulewidth}{0.4pt}
\fancypagestyle{plain}{\fancyhf{}\fancyfoot[C]{\thepage}\renewcommand{\headrulewidth}{0pt}}

% --- Estilos de gráficos comunes -----------------------------------------------
\pgfplotsset{
  mini/.style={width=4.4cm, height=4.2cm, ticks=none, axis lines=box,
               xlabel={$X$}, ylabel={$Y$}, xlabel style={yshift=2pt},
               ylabel style={rotate=-90, yshift=-2pt}, every axis plot/.append style={thin}},
  puntos/.style={only marks, mark=*, mark size=0.9pt, black},
  recta/.style={NavyBlue, thick, no marks},
  % gráficos de residuos (Tema 4): como mini, con r_i frente a los valores ajustados
  residuos/.style={mini, xlabel={$\est{y}_i$}, ylabel={$r_i$}, xmin=0, xmax=10},
  cero/.style={gray, dashed, no marks, domain=0:10},
}

% ==============================================================================
%  Entornos (cajas). Estilo sobrio, inspirado en el repositorio de referencia:
%  fondo tenue, marco fino y título en negrita dentro de la caja.
%
%  Uso:
%    \begin{definicion}[Título][etiqueta] ... \end{definicion}  (ambos opcionales)
%    \begin{teorema}[...][...]      ... \end{teorema}      (también: proposicion)
%    \begin{ejemplo}[...][...]      ... \end{ejemplo}
%    \begin{demostracion}[...]      ... \end{demostracion}
%    \begin{nota}[...]              ... \end{nota}
%    \begin{resumen}[Título]        ... \end{resumen}
% ==============================================================================

% --- Colores de las cajas --------------------------------------------------------
\definecolor{fondoDef}{RGB}{237,241,247}   % azul grisáceo muy claro
\definecolor{fondoTeo}{RGB}{236,236,238}   % gris claro
\colorlet{fondoDem}{Green!15}              % verde claro (como el repositorio)
\definecolor{fondoEje}{RGB}{246,245,241}   % gris cálido muy claro
\definecolor{fondoNot}{RGB}{249,249,249}   % casi blanco
\definecolor{marco}{RGB}{140,140,145}      % marco gris

% Colores usados en figuras
\colorlet{colRes}{Maroon}

% --- Estilo base ---------------------------------------------------------------
\tcbset{
  baseapuntes/.style={
    enhanced, breakable,
    colframe=marco, boxrule=0.4pt, arc=1pt,
    left=6pt, right=6pt, top=5pt, bottom=5pt,
    fonttitle=\bfseries, coltitle=black,
    detach title, before upper={\tcbtitle.\ },
    before skip=0.8em, after skip=0.8em,
  }
}

% --- Entornos numerados (numeración por tema: 1.1, 1.2, ...; contador común) ------
% Sintaxis: \begin{definicion}[Título opcional][etiqueta opcional]
%   p.ej.   \begin{definicion}[Correlación de Pearson][def:pearson]
%   y luego \cref{def:pearson}  ->  "definición 1.1"
\newcommand{\titulocaja}[2]{#1~\thetcbcounter\ifx\relax#2\relax\else\ (#2)\fi}
\tcbset{etiqueta/.code={\ifx\relax#1\relax\else\pgfkeysalso{label={#1}}\fi}}

\NewTColorBox[auto counter, number within=chapter, crefname={definición}{definiciones}]
  {definicion}{ O{} O{} }{baseapuntes, colback=fondoDef,
  title={\titulocaja{Definición}{#1}}, etiqueta={#2}}

\NewTColorBox[use counter from=definicion, crefname={teorema}{teoremas}]
  {teorema}{ O{} O{} }{baseapuntes, colback=fondoTeo,
  title={\titulocaja{Teorema}{#1}}, etiqueta={#2}}

\NewTColorBox[use counter from=definicion, crefname={proposición}{proposiciones}]
  {proposicion}{ O{} O{} }{baseapuntes, colback=fondoTeo,
  title={\titulocaja{Proposición}{#1}}, etiqueta={#2}}

\NewTColorBox[use counter from=definicion, crefname={ejemplo}{ejemplos}]
  {ejemplo}{ O{} O{} }{baseapuntes, colback=fondoEje,
  title={\titulocaja{Ejemplo}{#1}}, etiqueta={#2}}

% --- Entornos sin numerar -------------------------------------------------------
\newtcolorbox{demostracion}[1][]{
  baseapuntes, colback=fondoDem,
  fontupper=\setlength{\jot}{6pt},   % más separación entre líneas de align* (pasos en vertical)
  title={Demostración\ifx\relax#1\relax\else\ (#1)\fi},
  after upper={\par\hfill$\blacksquare$},
}

\newtcolorbox{nota}[1][]{
  baseapuntes, colback=fondoNot,
  title={Nota\ifx\relax#1\relax\else\ (#1)\fi},
}

\newtcolorbox{resumen}[1][Resumen]{
  baseapuntes, unbreakable,   % los resúmenes son cortos: nunca se parten entre páginas
  colback=white, boxrule=0.8pt,
  detach title, before upper={{\bfseries #1}\par\smallskip},
  title={#1},
}

% ==============================================================================
%  Comandos / notación común a todos los temas
%  (usar SIEMPRE estos comandos para que la notación sea homogénea)
% ==============================================================================

% --- Conjuntos y probabilidad ---------------------------------------------------
\newcommand{\R}{\mathbb{R}}
\DeclareMathOperator{\E}{E}                  % esperanza      \E[Y]
\DeclareMathOperator{\Var}{Var}              % varianza       \Var(Y)
\DeclareMathOperator{\Cov}{Cov}              % covarianza     \Cov(X,Y)
\DeclareMathOperator{\Cor}{Cor}              % correlación
\newcommand{\Normal}{\mathcal{N}}            % \Normal(0,\sigma^2)
\newcommand{\iid}{\overset{\text{i.i.d.}}{\sim}}
\newcommand{\bajoHo}{\overset{H_0}{\sim}}    % distribución bajo H0

% --- Vectores y matrices (negrita) ----------------------------------------------
\newcommand{\vect}[1]{\bm{#1}}               % \vect{Y}, \vect{\beta}
\newcommand{\T}{^{\mathsf{T}}}               % traspuesta: X\T

% --- Sumas de cuadrados y medias cuadráticas ------------------------------------
\newcommand{\SC}[1]{SS_{\text{#1}}}          % \SC{hip}   -> SS_hip
\newcommand{\MC}[1]{MS_{\text{#1}}}          % \MC{res}   -> MS_res

% --- Estadística descriptiva --------------------------------------------------
\newcommand{\media}[1]{\overline{#1}}        % \media{y}
\newcommand{\est}[1]{\widehat{#1}}           % estimador: \est{\mu}

% --- Distribuciones -----------------------------------------------------------
\newcommand{\Fdist}[2]{F_{#1,\,#2}}          % \Fdist{1}{n-1}
\newcommand{\tdist}[1]{t_{#1}}
\newcommand{\chisq}[1]{\chi^2_{#1}}

% --- Texto --------------------------------------------------------------------
\newcommand{\Rsoft}{\textsf{R}}              % el software R
\newcommand{\codigo}[1]{\texttt{#1}}

% --- Otros ----------------------------------------------------------------------
\newcommand{\cond}{\mid}                     % condicionado: \E(Y \cond X=x)
\newcommand{\gl}{\text{g.l.}}                % grados de libertad
\newcommand{\sumi}{\sum_{i=1}^{n}}           % suma habitual de 1 a n

% --- Regresión (Tema 3 en adelante) ----------------------------------------------
\newcommand{\Ynew}{Y_{h(\mathrm{new})}}      % nueva observación en x_h
\newcommand{\Xnew}{X_{h(\mathrm{new})}}      % predicción inversa
\newcommand{\Ymnew}{\media{Y}_{h(\mathrm{new})}} % media de m nuevas observaciones
\newcommand{\tq}[2]{t_{#1;\,#2}}             % cuantil: \tq{n-2}{1-\alpha/2}
\newcommand{\Fq}[3]{F_{#1,\,#2;\,#3}}        % cuantil: \Fq{1}{n-2}{1-\alpha}
\newcommand{\ee}{\operatorname{ee}}          % error estándar
\newcommand{\SSX}{SS_X}
\newcommand{\raizMSE}[1]{\sqrt{MSE\left(#1\right)}}

% --- Validación (Tema 4 en adelante) -------------------------------------------
\newcommand{\DFFits}{\mathrm{DFFits}}        % \DFFits_i
\newcommand{\DFBeta}{\mathrm{DFBeta}}        % \DFBeta_{ji}
\newcommand{\sumk}{\sum_{k=1}^{n}}           % suma en k de 1 a n

% --- Notación matricial (Tema 5 en adelante) ------------------------------------
\newcommand{\XtX}{\vect{X}\T\vect{X}}                 % X'X
\newcommand{\XtXi}{(\vect{X}\T\vect{X})^{-1}}         % (X'X)^{-1}
\newcommand{\bhat}{\est{\vect{\beta}}}                % beta estimado (vector)


% --- Compilar solo algunos temas (más rápido mientras se trabaja) -------------
% Descomenta y deja solo los temas que quieras. Las referencias cruzadas
% a los temas excluidos se conservan gracias a los .aux.
% \includeonly{temas/tema1/tema1}

\begin{document}

\pagenumbering{roman}
% ==============================================================================
%  Portada
% ==============================================================================
\thispagestyle{empty}
\begin{center}
  \vspace*{3cm}
  \rule{\linewidth}{1.2pt}\par
  \vspace{0.6cm}
  {\Huge\bfseries Regresión y ANOVA\par}
  \vspace{0.4cm}
  {\Large Apuntes de la asignatura\par}
  \vspace{0.4cm}
  \rule{\linewidth}{1.2pt}\par
  \vspace{2cm}
  {\large \textbf{Autor:} David Vaquerizo Medina\par}
  \vspace{0.6cm}
  {\large Grado en Estadística / INdat\par}
  {\large Universidad de Valladolid\par}
  {\large Curso 2026--2027\par}
\end{center}
\clearpage

% ==============================================================================
%  Aviso (página tras la portada)
% ==============================================================================
\thispagestyle{empty}
\begin{center}
  {\Huge\bfseries AVISO\par}
\end{center}
\vspace{1cm}

\setlength{\parindent}{1.5em}
Estos apuntes han sido elaborados con la ayuda de herramientas de inteligencia artificial, a
partir de las diapositivas de la asignatura \emph{Regresión y ANOVA} de la profesora Itziar
Fernández Martínez (Universidad de Valladolid). Su contenido se basa en dicho material, que ha
sido reorganizado y completado con las demostraciones y aclaraciones que se han considerado
necesarias.

Aunque el contenido ha sido revisado, los apuntes pueden contener errores o imprecisiones. No
sustituyen al material oficial de la asignatura, que debe consultarse ante cualquier duda.

Si detectas algún error, se agradece que se comunique para poder corregirlo:

\begin{center}
  \href{mailto:david.vaquerizo24@estudiantes.uva.es}{\texttt{david.vaquerizo24@estudiantes.uva.es}}
\end{center}
\setlength{\parindent}{0pt}

\vfill
\begin{center}
  {\small Versión compilada el \today\par}
\end{center}
\clearpage


{\hypersetup{linkcolor=black}%
 \renewcommand{\contentsname}{Índice --- Regresión}%
 \tableofcontents}
\cleardoublepage
\pagenumbering{arabic}

% ------------------------------------------------------------------------------
%  TEMAS DE REGRESIÓN (la lista está en bloque_regresion.tex)
% ------------------------------------------------------------------------------
% ==============================================================================
%  Bloque de REGRESIÓN: Temas 1–6
%  Se usa en main.tex (apuntes de regresión) y se reutilizará en el documento
%  global junto con el bloque de ANOVA. Añadir aquí una línea \include por tema.
% ==============================================================================
% ==============================================================================
%  TEMA 1 — Introducción a los modelos lineales
%  Fuente: diapositivas T1_IntroduccionModelosLineales (.pdf / .md)
% ==============================================================================
\chapter{Introducción a los modelos lineales}
\label{cap:tema1}

% ==============================================================================
%  1.1  Modelos estadísticos: utilidad y fundamentos
% ==============================================================================
\section{Modelos estadísticos: utilidad y fundamentos}
\label{sec:t1-modelos}

\subsection{La tarea del estadístico}

La tarea de un estadístico es \textbf{extraer información de un conjunto de datos} para:
\begin{itemize}
  \item describir un fenómeno,
  \item controlar o predecir el fenómeno,
  \item confirmar alguna hipótesis sobre el fenómeno,
  \item informar para la toma de decisiones.
\end{itemize}

\paragraph{El proceso estadístico.} De una \emph{población} se extrae una \emph{muestra}
(muestreo), que se resume mediante la \emph{estadística descriptiva}. Para extraer
conclusiones sobre la población a partir de la muestra (\emph{inferencia}) se necesita la
\emph{probabilidad}, es decir, describir cómo se han generado los datos de la muestra. Esta es
la tarea de la \textbf{modelización estadística}.

\begin{figure}[H]
  \centering
  \begin{tikzpicture}[
      caja/.style={draw, rounded corners, fill=fondoDef, minimum width=2.6cm,
                   minimum height=0.9cm, align=center, font=\small},
      flecha/.style={-{Stealth[length=2.5mm]}, thick}]
    \node[caja] (pob) at (0,0) {Población};
    \node[caja] (mue) at (7,0) {Muestra};
    \draw[flecha] (pob.north east) to[bend left=25]
      node[above, font=\small] {Muestreo} (mue.north west);
    \node[font=\small, right=0.3cm of mue] {$\to$ Estadística descriptiva};
    \draw[flecha, colRes!80!black] (mue.south west) to[bend left=25]
      node[below, font=\small, align=center]
      {\textbf{Probabilidad} + \textbf{Inferencia}\\[-1pt]
       \footnotesize (modelización estadística)} (pob.south east);
  \end{tikzpicture}
  \caption{El proceso estadístico. La modelización estadística permite pasar de la muestra
    a la población.}
\end{figure}

\subsection{Componente sistemática y componente aleatoria}

Los datos suelen combinar dos aspectos:
\begin{itemize}
  \item cierta regularidad en algún sentido: \textbf{componente sistemática};
  \item una variación irregular: \textbf{componente aleatoria}.
\end{itemize}

\begin{ejemplo}[Descomposición sistemática + aleatoria][ej:t1-componentes]
  \begin{enumerate}[label=(\alph*)]
    \item \emph{Consumo de cigarrillos de un fumador.} Varía irregularmente de semana en
      semana en torno a un nivel ``medio'': $C_i = m + e_i$.
    \item \emph{Ventas de patatas en un establecimiento.} Oscilan a lo largo del año con un
      patrón estacional, pero no todas las semanas son iguales: $V_t = m(t) + e_t$.
    \item \emph{Elección entre dos rutas, A y B, para cruzar la ciudad.} Los datos son los
      tiempos de viaje durante varios días usando las dos rutas. El tiempo de cada viaje se
      descompone en una componente fija, la ``norma'', y otra aleatoria (tráfico,
      semáforos\ldots): $t_r = T + u_r$, $r=1,\ldots,n$.
  \end{enumerate}
  En los tres casos, $m$, $m(t)$ y $T$ son la componente sistemática y $e_i$, $e_t$, $u_r$ la
  aleatoria.

  En (c), para saber qué ruta es más rápida no basta con comparar tiempos individuales, ya que
  cada uno incluye la componente aleatoria $u_r$. Hay que comparar las componentes fijas $T_A$
  y $T_B$ de cada ruta, lo que supone hacer inferencia sobre la diferencia de dos medias a
  partir de las dos muestras de tiempos (\cref{sec:t1-dos-medias}).
\end{ejemplo}

\subsection{Modelo estadístico}

\begin{definicion}[Modelo estadístico]
  La característica principal de un modelo estadístico es que \textbf{la variabilidad se
  representa por medio de distribuciones de probabilidad}. Por ello, la formulación del
  modelo incluye alguna especificación sobre la componente aleatoria (o \emph{perturbación}),
  como que sigue cierta distribución de probabilidad, por ejemplo $\Normal(0,\sigma^2)$.
\end{definicion}

\paragraph{Etapas del ajuste de un modelo estadístico.}
\begin{enumerate}
  \item \textbf{Especificación}: seleccionar una función apropiada para el modelo y las
    variables a incluir. Puede ser muy útil el análisis exploratorio de los datos
    (\cref{sec:t1-descriptivo}). Para un mismo conjunto de datos pueden ser apropiados
    diferentes modelos y tipos de análisis.
  \item \textbf{Estimación--ajuste}: seleccionar un método apropiado de estimación del
    modelo y estimar los parámetros.
  \item \textbf{Diagnóstico}: comprobar si el modelo propuesto es correcto. A veces esta etapa sirve
    para seleccionar ``el mejor modelo''.
\end{enumerate}

\paragraph{Tipos de modelos estadísticos.}
\begin{itemize}
  \item \textbf{Paramétricos}: el modelo queda perfectamente especificado una vez conocido
    el conjunto de parámetros que, a su vez, especifican la distribución de probabilidad
    subyacente. \emph{Ejemplo}: si la componente aleatoria tiene distribución
    $\Normal(0,\sigma^2)$, para estimar el modelo hay que estimar $\sigma$.
  \item \textbf{No paramétricos}: la distribución de los datos no puede definirse en
    términos de un conjunto de parámetros. \emph{Ejemplo}: se suponen las perturbaciones
    independientes y con media 0, sin asumir ninguna distribución concreta.
\end{itemize}

\begin{nota}[Notación]
  En estos apuntes $\Normal(\mu,\sigma^2)$ indica siempre una normal de media $\mu$ y
  \emph{varianza} $\sigma^2$ (desviación típica $\sigma$).
\end{nota}

\subsection{El modelo de regresión}

En muchos problemas existe una relación entre dos o más variables cuya naturaleza interesa
estudiar. El \textbf{análisis de regresión} es la técnica estadística que se utiliza para
modelizar e investigar este tipo de relaciones.

\begin{definicion}[Modelo y función de regresión][def:funcion-regresion]
  Los \textbf{modelos de regresión} describen cómo una \emph{variable respuesta} $Y$, tratada
  como \emph{aleatoria}, depende de un conjunto de \emph{variables explicativas}
  $X_1,\ldots,X_k$, tratadas como \emph{fijas}. Se supone que $Y$ tiene una distribución de
  probabilidad para cada nivel de las $X$, cuya media varía con ellas de forma sistemática.
  Esta relación sistemática es la \textbf{función de regresión de $Y$ sobre $X$}:
  \[
    m(x_1,\ldots,x_k) = \E\left(Y \cond X_1=x_1,\ldots,X_k=x_k\right).
  \]
\end{definicion}

Cada forma de la función de regresión y de la distribución de $Y$ da lugar a un modelo de
regresión distinto, que se estima a partir de los datos. Cuando no se hace ninguna suposición
sobre la forma de $m$ se obtienen los \textbf{modelos de regresión no paramétricos}, que se
ajustan estimando $m(x_1,\ldots,x_k)$ punto a punto.

\subsection{El modelo lineal}

\begin{definicion}[Modelo lineal][def:modelo-lineal]
  Un modelo es \textbf{lineal} si su componente sistemática es una \textbf{función lineal de
  sus parámetros}. Los ejemplos más sencillos son la constante $\alpha$ y la recta
  $\alpha+\beta x$. En ningún caso aparecen expresiones del tipo $\alpha^2$, $\alpha\beta$,
  $\log(\beta)$\ldots
\end{definicion}

\begin{nota}
  ``Lineal'' se refiere a los \emph{parámetros}, no a las variables: $Y=\beta_0+\beta_1X^2+U$
  o $Y=\beta_0+\beta_1\log X+U$ son modelos lineales; $Y=\beta_0 e^{\beta_1X}+U$ no lo es.
\end{nota}

Formalmente, el \textbf{modelo lineal general} se expresa como
\[
  \vect{Y} = Z\vect{\theta} + \vect{U},
\]
donde $\vect{Y}$ es el vector de la variable respuesta (univariante o multivariante), $Z$ una
matriz conocida que contiene los valores de las variables explicativas, $\vect{\theta}$ un
vector de parámetros y $\vect{U}$ un vector aleatorio. Puede acompañarse de un conjunto de
restricciones lineales sobre los parámetros, $R\vect{\theta}=\vect{c}$.

\paragraph{Modelos de regresión lineal.} Son aquellos en los que una respuesta continua depende
linealmente de un conjunto de variables explicativas,
\[
  Y = \beta_0 + \beta_1X_1 + \cdots + \beta_kX_k + U .
\]
Constituyen la base de la mayoría de los análisis estadísticos, aunque no son adecuados en
todos los casos: en muchas aplicaciones la respuesta $Y$ es \emph{discreta}, o los datos
sugieren la presencia de \emph{componentes no lineales}. En esas situaciones se recurre a:
\begin{itemize}
  \item \textbf{modelos no lineales} $Y=g(X_1,\ldots,X_k)+U$ con $g$ no lineal, que requieren
    nuevas herramientas de ajuste (por ejemplo, mínimos cuadrados ponderados);
  \item \textbf{modelos lineales generalizados}, con distribuciones de la familia
    exponencial para la respuesta, de la que la binomial y la Poisson son casos particulares
    importantes.
\end{itemize}

% ==============================================================================
%  1.2  Relación entre variables. Relaciones lineales y linealizables.
%       Asociación y causalidad.
% ==============================================================================
\section{Relación entre variables}
\label{sec:t1-relacion}

\subsection{Relación funcional y relación estadística}

\begin{definicion}[Relación funcional y relación estadística]
  \begin{itemize}
    \item Una \textbf{relación funcional}, expresada mediante una fórmula matemática, es una
      relación perfecta entre variables: el valor de las variables explicativas determina
      exactamente el valor de la respuesta.
    \item Una \textbf{relación estadística} es algo intermedio entre una relación funcional y
      la ausencia de relación: conocidas las variables explicativas, la respuesta no puede
      predecirse exactamente, pero sí de forma aproximada.
  \end{itemize}
\end{definicion}

El interés está en modelizar relaciones estadísticas, mediante modelos del tipo
\begin{equation}\label{eq:t1-modelo-general}
  Y = m(X_1,\ldots,X_k) + \varepsilon ,
\end{equation}
donde $m(X_1,\ldots,X_k)$ es el \emph{patrón sistemático} de la relación, en torno al cual se
distribuyen los datos con cierta variabilidad que recoge el \emph{error aleatorio}
$\varepsilon$. A este error se le atribuyen una serie de propiedades (distribución, media,
varianza\ldots).

\begin{figure}[H]
  \centering
  \begin{subfigure}[t]{0.48\linewidth}
    \centering
    \begin{tikzpicture}
      \begin{axis}[width=6cm, height=5cm, xlabel={Unidades vendidas ($X$)},
          ylabel={Ventas en \$ ($Y$)}, xmin=0, xmax=160, ymin=0, ymax=320]
        \addplot[recta, domain=0:150] {2*x};
        \addplot[puntos, mark size=1.5pt] coordinates {(25,50) (75,150) (130,260)};
        \node[font=\footnotesize, anchor=west] at (axis cs:95,120) {$Y=2X$};
      \end{axis}
    \end{tikzpicture}
    \caption{Relación funcional: los puntos están exactamente sobre la recta.}
  \end{subfigure}\hfill
  \begin{subfigure}[t]{0.48\linewidth}
    \centering
    \begin{tikzpicture}
      \begin{axis}[width=6cm, height=5cm, xlabel={Edad ($X$)},
          ylabel={Nivel de esteroides ($Y$)}, xmin=7, xmax=26, ymin=0, ymax=30]
        \addplot[recta, domain=8:25, samples=50] {-26.3+4.84*x-0.1165*x^2};
        \addplot[puntos, mark size=1.3pt] coordinates {
          (8,2) (8,6) (9,9.5) (9,12) (9,4) (10,11) (11,15) (12,11) (12,19) (13,13)
          (13,21) (14,21) (14,24) (16,22) (17,19) (17,24) (17,25) (18,21) (18,20.5)
          (19,22.5) (21,20) (21,24) (23,23) (23,27) (24.5,26) (25,22) (25,20)};
      \end{axis}
    \end{tikzpicture}
    \caption{Relación estadística: los puntos se distribuyen alrededor de un patrón
      sistemático (en este caso, curvo).}
  \end{subfigure}
  \caption{Relación funcional frente a relación estadística.}
\end{figure}

\begin{definicion}[Variables relacionadas][def:variables-relacionadas]
  Dos variables $X$ e $Y$ están \textbf{relacionadas} si la distribución condicional de
  $Y \cond X=x$ cambia al cambiar $x$. En caso contrario, las variables son
  \textbf{independientes}.
\end{definicion}

Uno de los objetivos de la modelización estadística es describir de forma sencilla cómo
cambian las distribuciones condicionales de $Y$ al cambiar $X$. La misma idea se extiende a la
relación entre $Y$ y un conjunto de variables explicativas $X_1,\ldots,X_k$.

\subsection{Selección de las variables explicativas}

Una parte esencial del ajuste de un modelo es la \textbf{selección de las variables
explicativas} $X_1,\ldots,X_k$, especialmente cuando $k$ es muy grande en relación con el
tamaño muestral $n$. Los criterios de selección son:
\begin{itemize}
  \item \textbf{Criterio principal}: en qué medida contribuyen a explicar la variabilidad de la
    respuesta. Son preferibles las que explican más.
  \item \textbf{Conocimiento del fenómeno}: por ejemplo, su importancia como agente causal.
  \item \textbf{Cuestiones prácticas de la toma de datos}: la precisión con la que puede
    observarse la variable y la facilidad, rapidez o coste de cada observación.
\end{itemize}

\subsection{Forma de la relación: relaciones lineales y linealizables}

Otra decisión importante es la \textbf{forma de la relación}. Normalmente no se conoce a priori
y debe decidirse empíricamente a partir de los datos. Las funciones lineales suelen ser una
buena elección por varias razones:
\begin{itemize}
  \item su \emph{simplicidad}, que es la más evidente, aunque no la única;
  \item muchas \emph{relaciones no lineales} pueden \textbf{linealizarse} mediante
    transformaciones de $X$, de $Y$ o de ambas (inversa, logaritmo, cuadrado, etc.);
  \item las relaciones de forma muy compleja pueden aproximarse mediante \textbf{funciones
    lineales ``a trozos''}.
\end{itemize}
Todo ello amplía el campo de aplicación de los modelos lineales y los convierte en una
herramienta fundamental para el ajuste de modelos estadísticos.

\begin{figure}[H]
  \centering
  \begin{tikzpicture}
    \begin{axis}[width=6.5cm, height=4.2cm, xmin=0, xmax=1.05, ymin=0, ymax=1.1,
        xtick={0,0.2,0.4,0.6,0.8,1},
        xticklabels={$x_0{=}0$,$x_1$,$x_2$,$x_3$,$x_4$,$x_5{=}1$},
        ytick=\empty, axis lines=left, tick label style={font=\scriptsize}]
      \addplot[NavyBlue, thick, domain=0:1, samples=60] {4*x*(1-x)};
      \addplot[red, thick, mark=none] coordinates
        {(0,0) (0.2,0.64) (0.4,0.96) (0.6,0.96) (0.8,0.64) (1,0)};
      \foreach \x in {0.2,0.4,0.6,0.8} {
        \edef\tmp{\noexpand\draw[dashed, gray] (axis cs:\x,0) -- (axis cs:\x,{4*\x*(1-\x)});}
        \tmp }
    \end{axis}
  \end{tikzpicture}
  \caption{Aproximación de una relación no lineal (azul) mediante una función lineal a
    trozos (rojo).}
\end{figure}

\subsection{Asociación y causalidad}
\label{sec:t1-causalidad}

Una relación puede ser \emph{débil}, en cuyo caso puede carecer de importancia práctica. Pero,
\textbf{aunque exista una relación estadística fuerte entre $X$ e $Y$, no puede asumirse que
$X$ sea la causa de $Y$} (relación causa--efecto):
\begin{itemize}
  \item la relación observada puede deberse a la influencia de terceras variables
    (\textbf{variables confusoras}), cuyos efectos solo pueden neutralizarse mediante el
    \emph{diseño de experimentos} (\cref{sec:t1-datos});
  \item puede existir una relación causal en sentido contrario, de $Y$ hacia $X$;
  \item hay incluso situaciones en las que la relación no es interpretable.
\end{itemize}

\begin{ejemplo}[Chocolate y premios Nobel][ej:t1-chocolate]
  En 2012 F.~H.~Messerli publicó en el \emph{New England Journal of Medicine} el estudio
  ``El consumo de chocolate, la función cognitiva y los Premios Nobel''. Con datos por
  países, encontró una relación muy fuerte ($r=0.791$, $p<0.0001$) entre el consumo anual de
  chocolate per cápita y el número de premios Nobel por cada 10 millones de habitantes.

  Poco después aparecieron otros trabajos (Maurage, Heeren y Pesenti, 2013, \emph{The Journal
  of Nutrition}) que señalaban otros índices relacionados con el número de premiados:
  \begin{center}
    \small
    \begin{tabular}{lc}
      \toprule
      Relación entre países & $r$ \\
      \midrule
      Nobel -- consumo de té            & $0.03$ \\
      Nobel -- consumo de vino          & $0.16$ \\
      Nobel -- tiendas IKEA por 10 M hab. & $0.82$ \\
      Nobel -- PIB per cápita           & $0.73$ \\
      Consumo de chocolate -- PIB per cápita & $0.66$ \\
      \bottomrule
    \end{tabular}
  \end{center}
  El número de tiendas de IKEA está más relacionado con el número de premios Nobel que el
  consumo de chocolate.

  Esta relación no permite concluir que comer chocolate aumente la probabilidad de obtener un
  premio Nobel. La tabla apunta a una \textbf{variable confusora}, la riqueza del país (PIB per
  cápita), relacionada tanto con el consumo de chocolate ($r=0.66$) como con el número de
  premios Nobel ($r=0.73$). Además, se trata de datos observacionales agregados por países,
  con los que no pueden establecerse relaciones causa--efecto (\cref{sec:t1-datos}).
\end{ejemplo}

% ==============================================================================
%  1.3  Tipos de datos según su procedencia o presentación
% ==============================================================================
\section{Tipos de datos según su procedencia}
\label{sec:t1-datos}

\subsection{Datos observacionales y experimentales}

Según procedan de estudios no experimentales o experimentales, los datos se clasifican en
\emph{observacionales} y \emph{experimentales}.

\begin{definicion}[Datos observacionales y experimentales]
  \begin{itemize}
    \item \textbf{Datos observacionales}: se obtienen de estudios no experimentales, en los
      que \emph{no se controlan} las variables explicativas de interés. Esto limita las
      conclusiones, ya que es muy difícil establecer relaciones causa--efecto.
    \item \textbf{Datos experimentales}: se obtienen de experimentos controlados y permiten
      extraer conclusiones más sólidas sobre posibles relaciones causa--efecto. La
      \emph{aleatorización} es un mecanismo importante para controlar las variables
      explicativas.
  \end{itemize}
\end{definicion}

Un experimento controlado tiene como objetivo comparar el efecto sobre la respuesta de
distintos \textbf{tratamientos} (combinaciones de niveles de las variables explicativas),
aplicándolos a diferentes \textbf{unidades experimentales} (individuos u objetos sobre los que
se realiza la investigación).

\begin{definicion}[Aleatorización]
  La \textbf{aleatorización} consiste en asignar de forma aleatoria cada unidad
  experimental a un tratamiento. Es útil porque tiende a equilibrar los posibles efectos de
  otras variables (no controladas) que podrían afectar a la respuesta.
\end{definicion}

\begin{ejemplo}[Efecto de fumar en la capacidad pulmonar]
  \begin{center}
    \small
    \begin{tabular}{p{0.45\linewidth}p{0.45\linewidth}}
      \toprule
      \textbf{Diseño experimental} & \textbf{Diseño observacional} \\
      \midrule
      Elegir 100 personas que actualmente no fumen. & --- \\[2pt]
      Asignar \textbf{aleatoriamente} 50 de ellas al grupo ``fumadores'' y las otras 50 al
      grupo ``no fumadores''. & --- \\[2pt]
      Cada persona del grupo ``fumadores'' fumará un paquete al día durante 10 años. &
      Encontrar a 50 personas fumadoras de un paquete al día durante 10 años. \\[2pt]
      Cada persona del grupo ``no fumadores'' permanecerá sin fumar durante 10 años. &
      Encontrar a 50 personas que no hayan fumado durante 10 años. \\[2pt]
      Transcurridos los 10 años, medir la capacidad pulmonar de las 100 personas. &
      Medir la capacidad pulmonar de las 100 personas. \\
      \bottomrule
    \end{tabular}
  \end{center}
  En el diseño observacional, los fumadores y los no fumadores pueden diferir en otras
  variables (edad, trabajo, deporte\ldots) que afectan a la capacidad pulmonar; en el
  experimental, la aleatorización equilibra esos efectos entre los dos grupos.
\end{ejemplo}

\paragraph{Estudios muestrales.} Entre ambos tipos de estudio se sitúan los \textbf{estudios
muestrales}: las inferencias se realizan sobre la población adecuada, pero la presencia de
variables influyentes no controladas o desconocidas puede dificultar la interpretación de las
relaciones entre las variables observadas.

\paragraph{Precauciones en el análisis.} Al analizar los datos hay que tener en cuenta:
\begin{itemize}
  \item posibles \textbf{cambios en las condiciones de recogida} de los datos (de personal, de
    máquinas, envejecimiento\ldots) que puedan influir en los resultados;
  \item la presencia de variables explicativas \textbf{fuertemente correlacionadas} entre sí,
    cuyos efectos se confunden;
  \item un \textbf{rango pequeño} de alguna variable importante, que impide establecer su
    influencia fuera de ese rango.
\end{itemize}

\subsection{Efectos fijos, aleatorios y mixtos}

Según cómo se recojan las variables explicativas se distingue entre:

\begin{definicion}[Modelos de efectos fijos, aleatorios y mixtos]
  \begin{itemize}
    \item \textbf{Modelo de efectos fijos.} Una \emph{variable fija} es aquella que se supone
      medida sin error, de modo que los valores observados en un estudio serían los mismos en
      otro. Las variables explicativas toman valores predeterminados, que se controlan en el
      diseño y la realización del experimento.
    \item \textbf{Modelo de efectos aleatorios.} Las variables explicativas son \emph{variables
      aleatorias}, cuyos valores son una muestra aleatoria de todos los valores posibles, y se
      espera poder generalizar los resultados al resto de valores. Esta generalización debe
      hacerse con cautela, ya que los resultados solo son válidos en el rango de variación
      conjunta de las variables implicadas.
    \item \textbf{Modelo mixto}: combina efectos fijos y aleatorios. En muchas aplicaciones los
      modelos son mixtos.
  \end{itemize}
\end{definicion}

\begin{ejemplo}[Dosis de un fármaco]
  Se mide la efectividad de un fármaco con dosis de 5, 10 y 15\,mg.
  \begin{itemize}
    \item Si esas tres dosis se han elegido intencionadamente y los resultados solo se van a
      aplicar a ellas, ``dosis'' es un efecto \textbf{fijo}.
    \item Si se consideran una muestra aleatoria de todas las dosis posibles entre 5 y 15\,mg
      y los resultados se quieren aplicar a cualquier dosis de ese rango, ``dosis'' es un
      efecto \textbf{aleatorio}.
  \end{itemize}
\end{ejemplo}

% ==============================================================================
%  1.4  Análisis descriptivo previo a la formulación del modelo.
%       Gráficos. Correlación.
% ==============================================================================
\section{Análisis descriptivo previo al modelo}
\label{sec:t1-descriptivo}

El análisis descriptivo de los datos ayuda a formular el modelo (etapa de
\emph{especificación}). Conviene tener en cuenta las escalas elegidas y el rango de variación de
$X$ e $Y$, y completarlo con otros métodos descriptivos, gráficos o numéricos.

\subsection{El diagrama de dispersión}

Para estudiar gráficamente la dependencia entre \textbf{dos variables} resulta muy útil el
\textbf{diagrama (o gráfico) de dispersión}. Los datos son pares $(x_i,y_i)$, $i=1,\ldots,n$, y
el objetivo es ver cómo cambian los valores de $Y$ al cambiar los de $X$. Cada par se
representa como un punto, con $x_i$ en el eje horizontal (\textbf{variable explicativa}) e
$y_i$ en el vertical (\textbf{variable respuesta}). La forma de la nube de puntos da una idea
de si existe relación entre $X$ e $Y$ y de qué forma tiene.

\begin{ejemplo}[Edad y distancia de lectura de una señal][ej:conductores]
  En una muestra de 30 conductores de entre 18 y 82 años se mide la distancia máxima (en pies)
  a la que pueden leer una señal de tráfico. Se quiere estudiar la relación entre la edad y esa
  distancia, con el objetivo de mejorar la seguridad de los conductores mayores. Como interesa
  el efecto de la edad sobre la distancia:
  \begin{itemize}
    \item la \textbf{edad} es la variable explicativa $X$;
    \item la \textbf{distancia} es la variable respuesta $Y$.
  \end{itemize}
  Cada conductor corresponde a un punto: el conductor 1 (18 años, 510 pies) es el punto
  $(18,510)$, el conductor 2 (32 años, 410 pies) es el punto $(32,410)$, etc.

  \begin{center}
    \begin{tikzpicture}
      \begin{axis}[width=8cm, height=5.5cm, xlabel={Edad del conductor (años)},
          ylabel={Distancia (pies)}, xmin=15, xmax=85, ymin=260, ymax=610,
          label style={font=\small}, tick label style={font=\footnotesize}]
        \addplot[puntos, mark size=1.3pt] table[x=edad, y=distancia]
          {datos/conductores.dat};
        \addplot[only marks, mark=*, mark size=2pt, red] coordinates {(18,510)};
        \addplot[only marks, mark=*, mark size=2pt, blue] coordinates {(32,410)};
        \draw[red, dashed] (axis cs:18,260) -- (axis cs:18,510) -- (axis cs:15,510);
        \draw[blue, dashed] (axis cs:32,260) -- (axis cs:32,410) -- (axis cs:15,410);
      \end{axis}
    \end{tikzpicture}\\
    {\footnotesize Datos reconstruidos a partir del gráfico de las diapositivas.}
  \end{center}
\end{ejemplo}

\subsection{Interpretación de un diagrama de dispersión}

Al interpretar un diagrama de dispersión se analizan:
\begin{itemize}
  \item el \textbf{patrón general} de la relación: \emph{dirección}, \emph{forma} y
    \emph{fuerza};
  \item las \textbf{desviaciones} respecto de ese patrón: los \emph{outliers} (datos atípicos).
\end{itemize}

\paragraph{Dirección.} Puede ser positiva, negativa o no estar definida.
\begin{itemize}
  \item \textbf{Relación positiva}: los incrementos de una variable van asociados a
    incrementos de la otra.
  \item \textbf{Relación negativa}: los incrementos de una variable van asociados a
    disminuciones de la otra.
  \item No todas las relaciones pueden clasificarse como positivas o negativas.
\end{itemize}

\begin{figure}[H]
  \centering
  \begin{subfigure}[t]{0.32\linewidth}
    \centering
    \begin{tikzpicture}
      \begin{axis}[mini]
        \addplot[puntos] table {datos/tema1/dir_pos.dat};
      \end{axis}
    \end{tikzpicture}
    \caption{Relación positiva.}
  \end{subfigure}\hfill
  \begin{subfigure}[t]{0.32\linewidth}
    \centering
    \begin{tikzpicture}
      \begin{axis}[mini]
        \addplot[puntos] table {datos/tema1/dir_neg.dat};
      \end{axis}
    \end{tikzpicture}
    \caption{Relación negativa.}
  \end{subfigure}\hfill
  \begin{subfigure}[t]{0.32\linewidth}
    \centering
    \begin{tikzpicture}
      \begin{axis}[mini]
        \addplot[puntos] table {datos/tema1/dir_nula.dat};
      \end{axis}
    \end{tikzpicture}
    \caption{Sin dirección clara.}
  \end{subfigure}
  \caption{Dirección de una relación.}
\end{figure}

\paragraph{Forma.} Lineal, curva (cuadrática, exponencial\ldots), por grupos, etc.

\paragraph{Fuerza.} Depende de lo cerca que se sitúen los puntos del patrón que describe la
relación. En una relación lineal \textbf{fuerte} los puntos están muy próximos a la recta; en
una relación lineal más \textbf{débil} siguen el mismo patrón lineal, pero mucho más dispersos a
su alrededor.

\begin{figure}[H]
  \centering
  \begin{subfigure}[t]{0.45\linewidth}
    \centering
    \begin{tikzpicture}
      \begin{axis}[mini]
        \addplot[puntos] table {datos/tema1/fuerte.dat};
        \addplot[recta] table[y={create col/linear regression={y=y}}] {datos/tema1/fuerte.dat};
      \end{axis}
    \end{tikzpicture}
    \caption{Relación lineal fuerte: puntos muy próximos a la recta.}
  \end{subfigure}\hfill
  \begin{subfigure}[t]{0.45\linewidth}
    \centering
    \begin{tikzpicture}
      \begin{axis}[mini]
        \addplot[puntos] table {datos/tema1/debil.dat};
        \addplot[recta] table[y={create col/linear regression={y=y}}] {datos/tema1/debil.dat};
      \end{axis}
    \end{tikzpicture}
    \caption{Relación lineal débil: puntos más dispersos alrededor de la recta.}
  \end{subfigure}
  \caption{Fuerza de una relación lineal.}
\end{figure}

\paragraph{Outliers.} Son observaciones atípicas, es decir, datos que se apartan del patrón de
la relación. Pueden darse distintos casos:
\begin{itemize}
  \item \textbf{Atípico en $X$}: su valor de $X$ es mucho mayor (o menor) que el del resto,
    pero sigue el patrón de la relación, por lo que \emph{no} es atípico en la relación.
  \item \textbf{Atípico en la relación}: los valores de $X$ e $Y$ no son extremos por separado,
    pero el par no sigue el patrón del resto.
\end{itemize}

\begin{figure}[H]
  \centering
  \begin{subfigure}[t]{0.45\linewidth}
    \centering
    \begin{tikzpicture}
      \begin{axis}[mini, width=5.5cm]
        \addplot[puntos] table {datos/tema1/atipico_x.dat};
        \addplot[only marks, mark=*, mark size=1.8pt, red] coordinates {(40,640)};
        \addplot[recta] table[y={create col/linear regression={y=y}}]
          {datos/tema1/atipico_x.dat};
      \end{axis}
    \end{tikzpicture}
    \caption{Atípico en $X$ (en rojo): valor de $X$ extremo, pero sigue la relación.}
  \end{subfigure}\hfill
  \begin{subfigure}[t]{0.45\linewidth}
    \centering
    \begin{tikzpicture}
      \begin{axis}[mini, width=5.5cm]
        \addplot[puntos] table {datos/tema1/atipico_rel.dat};
        \addplot[only marks, mark=*, mark size=1.8pt, red] coordinates {(0.5,9.2)};
      \end{axis}
    \end{tikzpicture}
    \caption{Atípico en la relación (en rojo): el par no sigue el patrón del resto.}
  \end{subfigure}
  \caption{Tipos de datos atípicos.}
\end{figure}

\begin{ejemplo}[Conductores: interpretación][ej:conductores-interp]
  En el \cref{ej:conductores}:
  \begin{itemize}
    \item \textbf{Dirección}: negativa; a medida que un conductor envejece, disminuye la
      distancia a la que es capaz de identificar la señal.
    \item \textbf{Forma}: aparentemente lineal.
    \item \textbf{Fuerza}: es difícil de valorar sin una referencia; parece
      \emph{moderadamente fuerte}, ya que los datos están bastante dispersos en torno a la
      recta que resume la relación.
    \item \textbf{Outliers}: todas las observaciones siguen razonablemente bien el patrón
      lineal, por lo que no parece haber valores atípicos.
  \end{itemize}
  \begin{center}
    \begin{tikzpicture}
      \begin{axis}[width=7.5cm, height=5cm, xlabel={Edad (años)}, ylabel={Distancia (pies)},
          xmin=15, xmax=85, ymin=260, ymax=610, label style={font=\small},
          tick label style={font=\footnotesize}]
        \addplot[puntos, mark size=1.3pt] table[x=edad, y=distancia] {datos/conductores.dat};
        \addplot[recta, domain=17:83] {576-3*x};
      \end{axis}
    \end{tikzpicture}
  \end{center}
\end{ejemplo}

\paragraph{Varias variables explicativas.} Cuando hay varias variables se utiliza una
\textbf{matriz de gráficos de dispersión}, formada por los diagramas de dispersión de cada par
de variables.

\begin{ejemplo}[Consumo de combustible (Weisberg)]
  Datos de 2001 de los estados de EE.\,UU.:
  \begin{center}
    \small
    \begin{tabular}{ll}
      \toprule
      \textit{Drivers} & Número de conductores con licencia en el estado \\
      \textit{FuelC}   & Gasolina vendida para uso en carretera (miles de galones) \\
      \textit{Income}  & Renta personal per cápita en 2000 (miles de dólares) \\
      \textit{Miles}   & Millas de autopistas con ayuda federal en el estado \\
      \textit{Pop}     & Población de 16 años o más en 2001 \\
      \textit{Tax}     & Impuesto estatal sobre la gasolina (céntimos por galón) \\
      \textit{State}   & Nombre del estado \\
      \midrule
      \textit{Fuel}    & $1000\times \textit{FuelC}/\textit{Pop}$ \\
      \textit{Dlic}    & $1000\times \textit{Drivers}/\textit{Pop}$ \\
      $\log(\textit{Miles})$ & Logaritmo en base 2 de \textit{Miles} \\
      \bottomrule
    \end{tabular}
  \end{center}
  La matriz de dispersión de \textit{Tax}, \textit{Dlic}, \textit{Income},
  $\log(\textit{Miles})$ y \textit{Fuel} es una cuadrícula $5\times5$ en la que la casilla
  $(i,j)$ contiene el diagrama de dispersión de la variable de la fila $i$ frente a la de la
  columna $j$, y la diagonal, los nombres de las variables. Permite ver de un vistazo cómo se
  relaciona la respuesta (\textit{Fuel}) con cada variable explicativa, y estas entre sí.
\end{ejemplo}

\subsection{El coeficiente de correlación lineal de Pearson}
\label{sec:t1-correlacion}

La \textbf{correlación} mide el grado de \textbf{relación lineal} entre $X$ e $Y$. Para dos
variables cuantitativas se utiliza el \textbf{coeficiente de correlación lineal de Pearson}.

\begin{definicion}[Coeficiente de correlación de Pearson][def:pearson]
  \begin{itemize}
    \item \textbf{Poblacional}:
      \[
        \rho_{XY} = \frac{\E\big[(X-\E X)(Y-\E Y)\big]}{\sqrt{\Var(X)}\,\sqrt{\Var(Y)}} .
      \]
    \item \textbf{Muestral}:
      \[
        r_{xy} = \frac{1}{n}\sumi \frac{(x_i-\media{x})(y_i-\media{y})}{s_x\,s_y},
        \qquad
        s_x=\sqrt{\frac1n\sumi(x_i-\media{x})^2},\quad
        s_y=\sqrt{\frac1n\sumi(y_i-\media{y})^2}.
      \]
  \end{itemize}
\end{definicion}

\begin{nota}
  En la fórmula muestral, $s_x$ y $s_y$ son las desviaciones típicas \emph{con divisor $n$},
  coherentes con el factor $\tfrac1n$. Al simplificar, estos factores se cancelan:
  \[
    r_{xy} = \frac{\sumi (x_i-\media{x})(y_i-\media{y})}
                  {\sqrt{\sumi (x_i-\media{x})^2}\,\sqrt{\sumi (y_i-\media{y})^2}},
  \]
  por lo que se obtiene el mismo valor usando $\tfrac{1}{n-1}$ \emph{a la vez} en la
  covarianza y en las dos varianzas. Lo que no es correcto es mezclar ambos divisores.
\end{nota}

\paragraph{Propiedades.}
\begin{itemize}
  \item Es \textbf{adimensional}.
  \item Toma valores entre $-1$ y $1$, y solo vale $-1$ o $1$ si existe una relación lineal
    \emph{exacta} entre las variables.
  \item Es \textbf{invariante ante cambios de escala}: no varía al cambiar las unidades de
    medida de cualquiera de las dos variables (salvo, quizá, el signo).
  \item Es \textbf{invariante ante cambios de origen} y, por tanto, ante transformaciones
    lineales (salvo el signo).
\end{itemize}

\paragraph{Interpretación.}
\begin{itemize}
  \item El \textbf{signo} indica la dirección de la relación lineal: positiva ($r_{xy}>0$) o
    negativa ($r_{xy}<0$).
  \item La \textbf{distancia a 0} indica la fuerza de la relación lineal: los valores próximos
    a $1$ o $-1$ indican una relación lineal fuerte, y los próximos a $0$, ausencia de relación
    lineal.
\end{itemize}

\begin{figure}[H]
  \centering
  \begin{subfigure}[t]{0.32\linewidth}
    \centering
    \begin{tikzpicture}
      \begin{axis}[mini]
        \addplot[puntos, mark size=0.6pt] table {datos/tema1/r_097.dat};
        \addplot[recta] table[y={create col/linear regression={y=y}}] {datos/tema1/r_097.dat};
      \end{axis}
    \end{tikzpicture}
    \caption{$r=0.97$: relación lineal positiva fuerte.}
  \end{subfigure}\hfill
  \begin{subfigure}[t]{0.32\linewidth}
    \centering
    \begin{tikzpicture}
      \begin{axis}[mini]
        \addplot[puntos, mark size=0.6pt] table {datos/tema1/r_m079.dat};
        \addplot[recta] table[y={create col/linear regression={y=y}}] {datos/tema1/r_m079.dat};
      \end{axis}
    \end{tikzpicture}
    \caption{$r=-0.79$: relación lineal negativa.}
  \end{subfigure}\hfill
  \begin{subfigure}[t]{0.32\linewidth}
    \centering
    \begin{tikzpicture}
      \begin{axis}[mini]
        \addplot[puntos, mark size=0.6pt] table {datos/tema1/r_nulo.dat};
        \addplot[recta] table[y={create col/linear regression={y=y}}] {datos/tema1/r_nulo.dat};
      \end{axis}
    \end{tikzpicture}
    \caption{$r=-0.02$: ausencia de relación lineal (plot nulo).}
  \end{subfigure}
  \caption{Nubes de puntos y su coeficiente de correlación.}
\end{figure}

\paragraph{Precaución.} $r_{xy}$ \textbf{por sí solo no basta} para determinar si una relación
es lineal o no: \textbf{es necesario el diagrama de dispersión}. La siguiente figura muestra
tres situaciones en las que el valor de $r$ resulta engañoso.

\begin{figure}[H]
  \centering
  \begin{subfigure}[t]{0.32\linewidth}
    \centering
    \begin{tikzpicture}
      \begin{axis}[mini]
        \addplot[puntos, mark size=0.6pt] table {datos/tema1/no_lineal.dat};
        \addplot[recta] table[y={create col/linear regression={y=y}}]
          {datos/tema1/no_lineal.dat};
        \node[font=\scriptsize, anchor=north] at (rel axis cs:0.5,0.97) {$r=-0.06$};
      \end{axis}
    \end{tikzpicture}
    \caption{Relación no lineal: existe una relación clara (parabólica), pero $r\approx0$
      porque no es lineal.}
  \end{subfigure}\hfill
  \begin{subfigure}[t]{0.32\linewidth}
    \centering
    \begin{tikzpicture}
      \begin{axis}[mini]
        \addplot[only marks, mark=*, mark size=0.7pt, NavyBlue] table {datos/tema1/grupo1.dat};
        \addplot[only marks, mark=*, mark size=0.7pt, ForestGreen] table {datos/tema1/grupo2.dat};
        \addplot[black, thin, no marks] table[y={create col/linear regression={y=y}}]
          {datos/tema1/grupos.dat};
        \node[font=\scriptsize, anchor=north west] at (rel axis cs:0.02,0.97) {$r=0.98$};
        \node[font=\scriptsize, NavyBlue] at (rel axis cs:0.25,0.35) {$r=0.21$};
        \node[font=\scriptsize, ForestGreen] at (rel axis cs:0.75,0.55) {$r=0.05$};
      \end{axis}
    \end{tikzpicture}
    \caption{Grupos de individuos: apenas hay relación dentro de cada grupo, pero al
      juntarlos $r=0.98$.}
  \end{subfigure}\hfill
  \begin{subfigure}[t]{0.32\linewidth}
    \centering
    \begin{tikzpicture}
      \begin{axis}[mini]
        \addplot[puntos, mark size=0.7pt] table {datos/tema1/sin_atipico.dat};
        \addplot[only marks, mark=o, mark size=2.5pt, red] coordinates {(10,10)};
        \addplot[only marks, mark=*, mark size=0.8pt, black] coordinates {(10,10)};
        \addplot[black, thin, no marks] table[y={create col/linear regression={y=y}}]
          {datos/tema1/con_atipico.dat};
        \addplot[red, thin, no marks] table[y={create col/linear regression={y=y}}]
          {datos/tema1/sin_atipico.dat};
        \node[font=\scriptsize, anchor=north west] at (rel axis cs:0.02,0.97) {$r=0.70$};
        \node[font=\scriptsize, red, anchor=north west] at (rel axis cs:0.02,0.82) {$r=-0.13$};
      \end{axis}
    \end{tikzpicture}
    \caption{Dato atípico (en rojo): sin él $r=-0.13$; al incluirlo, $r=0.70$.}
  \end{subfigure}
  \caption{Situaciones en las que el coeficiente de correlación es engañoso.}
\end{figure}

% ==============================================================================
%  1.5  El modelo lineal. Inferencias sobre una, dos y varias medias.
%       Introducción al ANOVA.
% ==============================================================================
\section{El modelo lineal para inferencias sobre medias. Introducción al ANOVA}
\label{sec:t1-medias}

Esta sección trata los métodos de \textbf{inferencia sobre las medias} de una, dos y varias
poblaciones. Estas inferencias pueden basarse en los principios de las distribuciones
muestrales o en \textbf{modelos lineales}. En esta asignatura interesan los modelos lineales,
que también se apoyan en las \textbf{distribuciones muestrales}, aunque con un enfoque algo
distinto.

\paragraph{Proceso de ``hacer inferencia''.}
\begin{enumerate}
  \item Se identifican uno o más \textbf{parámetros} que caracterizan la población de interés.
  \item Se toma una \textbf{muestra representativa} de la población.
  \item Para estimar el parámetro se elige un \textbf{estadístico} que pueda calcularse con la
    muestra. Al calcularse a partir de una muestra aleatoria, es una variable aleatoria.
  \item Se determina su \textbf{distribución en el muestreo}, es decir, la distribución
    teórica del estadístico en las distintas muestras posibles, que describe el
    comportamiento de todos sus valores.
  \item Con esa distribución y el valor observado en la muestra, se \textbf{hace inferencia}
    sobre el parámetro.
\end{enumerate}

\begin{definicion}[Suma de cuadrados y media cuadrática][def:SS-MS]
  \begin{itemize}
    \item \textbf{Suma de cuadrados} ($SS$): suma de las desviaciones cuadráticas de los valores
      observados respecto de la media muestral,
      \[ SS_{xx} = \sumi (x_i-\media{x})^2 . \]
    \item \textbf{Media cuadrática} ($MS$): suma de cuadrados dividida por sus grados de
      libertad,
      \[ MS_x = \frac{SS_{xx}}{n-1} = \frac{\sumi (x_i-\media{x})^2}{n-1} = S^2 , \]
      es decir, la varianza muestral corregida $S^2$ es una media cuadrática.
  \end{itemize}
\end{definicion}

\begin{nota}
  Análogamente, $SS_{xy}=\sumi(x_i-\media{x})(y_i-\media{y})$. Con esta notación, el
  coeficiente de correlación de la \cref{def:pearson} es
  $r_{xy} = SS_{xy}/\sqrt{SS_{xx}\,SS_{yy}}$.
\end{nota}

Los tres contrastes que se estudian a continuación siguen el mismo esquema:
\begin{resumen}[Esquema común: modelo restringido frente a no restringido]
  \begin{enumerate}
    \item Se plantea un modelo lineal \textbf{no restringido} (compatible con $H_1$) y el modelo
      \textbf{restringido} que impone $H_0$.
    \item Se estima cada uno por \textbf{mínimos cuadrados} y se calcula su suma de cuadrados
      de los errores. Como el modelo restringido es un caso particular del no restringido,
      $\SC{restringido} \ge \SC{no restringido}$.
    \item El aumento $\SC{hipótesis}=\SC{restringido}-\SC{no restringido}$ mide cuánto empeora
      el ajuste al imponer $H_0$. Si $H_0$ es cierta, este aumento debería ser pequeño.
    \item Se comparan $\MC{hipótesis}$ y $\MC{no restringido}$ con un estadístico $F$ y se
      rechaza $H_0$ si $F$ es grande.
  \end{enumerate}
\end{resumen}

% ------------------------------------------------------------------------------
\subsection{Inferencias para una media}
\label{sec:t1-una-media}

Se parte de una muestra aleatoria de tamaño $n$ de una población normal con media $\mu$ y
desviación típica $\sigma$, y el objetivo es hacer inferencias sobre $\mu$. Para ello se
plantea el \textbf{modelo lineal}
\[
  y_i = \mu + \varepsilon_i, \qquad i=1,\ldots,n,
\]
donde $y_i$ es el $i$-ésimo valor observado de la respuesta $Y$, $\mu$ es la media poblacional
y $\varepsilon_1,\ldots,\varepsilon_n$ son variables aleatorias independientes
$\Normal(0,\sigma^2)$.
\begin{itemize}
  \item \textbf{Componente sistemática}: $\mu$. En ausencia de variabilidad ($\sigma=0$),
    $y_i=\mu$ para todo $i$. Como la media de los $\varepsilon_i$ es 0, $\E[Y]=\mu$.
  \item \textbf{Componente aleatoria} (``error''): los $\varepsilon_i=y_i-\mu$, que describen la
    variabilidad de los valores individuales alrededor de $\mu$.
  \item \textbf{Parámetro $\sigma^2$}: es la varianza de $\varepsilon$ y mide la dispersión del
    error, es decir, la variabilidad del modelo. Sirve como \emph{medida del ajuste}: una
    varianza pequeña indica que los valores observados están cerca de la media, y por tanto
    un buen ajuste.
\end{itemize}

\paragraph{Estimación por mínimos cuadrados.} El criterio más utilizado para estimar $\mu$ es
el \textbf{criterio de mínimos cuadrados}, que consiste en buscar el valor $\est{\mu}$ que
minimiza la suma de los cuadrados de las diferencias entre los valores observados y la media
estimada,
\[
  SS_{yy} = \sumi \est{\varepsilon}_i^{\,2} = \sumi (y_i-\est{\mu})^2 .
\]
Derivando respecto de $\mu$ e igualando a cero,
\[
  \frac{\partial}{\partial \mu}\sumi (y_i-\mu)^2 = -2\sumi y_i + 2n\mu = 0
  \quad\Longrightarrow\quad
  \est{\mu} = \frac{\sumi y_i}{n} = \media{y} .
\]
Por tanto, $SS_{yy}=\sumi (y_i-\media{y})^2$, que es el numerador de $s_y^2$, la estimación de
$\sigma^2$. Es decir, $\est{\mu}=\media{y}$ minimiza la variabilidad del modelo y proporciona
el modelo que mejor se ajusta a los datos.

\paragraph{Contraste $H_0\colon \mu=\mu_0$ frente a $H_1\colon \mu\neq\mu_0$.} Hay que elegir
entre dos modelos:
\begin{enumerate}
  \item \textbf{Modelo restringido} (verifica $H_0$): $y_i = \mu_0 + \varepsilon_i$.
  \item \textbf{Modelo no restringido} ($\mu$ libre; apoyaría la alternativa):
    $y_i = \mu + \varepsilon_i$.
\end{enumerate}
Se elige el que minimiza la varianza. Las sumas de cuadrados de ambos modelos son
\[
  \SC{restringido} = \sumi (y_i-\mu_0)^2, \qquad
  \SC{no restringido} = \sumi (y_i-\media{y})^2 .
\]
Como $\media{y}$ es el estimador mínimo cuadrático, minimiza la suma de cuadrados y
$\SC{restringido} \ge \SC{no restringido}$. La distancia entre ambas es la base del contraste:
\begin{align*}
  \SC{restringido} &= \sumi (y_i-\mu_0)^2 = \sumi (y_i-\media{y}+\media{y}-\mu_0)^2 \\
  &= \sumi (y_i-\media{y})^2 + \sumi (\media{y}-\mu_0)^2
     + 2(\media{y}-\mu_0)\underbrace{\sumi (y_i-\media{y})}_{=\,0} \\
  &= \SC{no restringido} + n(\media{y}-\mu_0)^2 .
\end{align*}

\begin{proposicion}[Descomposición de la variabilidad (una media)][prop:desc-una-media]
  \[
    \underbrace{\sumi (y_i-\mu_0)^2}_{\SC{restringido}}
    = \underbrace{\sumi (y_i-\media{y})^2}_{\SC{no restringido}}
    + \underbrace{n(\media{y}-\mu_0)^2}_{\SC{hipótesis}}
  \]
  \begin{itemize}
    \item $\SC{restringido}$ tiene $n$ grados de libertad: es la suma de $n$ desviaciones
      respecto de una cantidad, $\mu_0$, que no se calcula a partir de los datos.
    \item $\SC{no restringido}$ tiene $n-1$ grados de libertad: suma de $n$ desviaciones
      respecto de $\media{y}$, estimada por mínimos cuadrados.
    \item $\SC{hipótesis}$ es lo que aumenta el error al imponer que $H_0$ es cierta; tiene
      1 grado de libertad.
  \end{itemize}
  La igualdad también se cumple para los grados de libertad: $n=(n-1)+1$.
\end{proposicion}

\paragraph{Medias cuadráticas y estadístico $F$.}
\[
  \MC{no restringido} = \frac{\sumi (y_i-\media{y})^2}{n-1}, \qquad
  \MC{hipótesis} = \frac{n(\media{y}-\mu_0)^2}{1}.
\]
Ambas estiman la misma varianza $\sigma^2$ si $H_0$ es cierta, y entonces
\[
  F = \frac{\MC{hipótesis}}{\MC{no restringido}} \bajoHo \Fdist{1}{n-1} .
\]
Si $H_0$ no es cierta, el numerador (y por tanto $F$) tiende a ser mayor. \textbf{Se rechaza
$H_0$ si $F_{\text{obs}} > F_{1,\,n-1;\,1-\alpha}$}.

\begin{nota}
  Como $\MC{no restringido}=s^2$, se tiene $F=\left(\dfrac{\media{y}-\mu_0}{s/\sqrt{n}}\right)^2=t^2$:
  es el cuadrado del estadístico $t$ habitual para una media, y $F_{1,n-1}$ es la distribución
  de $t_{n-1}^2$. Ambos contrastes son equivalentes.
\end{nota}

% ------------------------------------------------------------------------------
\subsection{Inferencias para dos medias}
\label{sec:t1-dos-medias}

Se consideran dos poblaciones de una variable $Y$ normal, con medias $\mu_1$, $\mu_2$ y
varianzas $\sigma_1^2$, $\sigma_2^2$, de las que se obtienen muestras aleatorias de tamaños
$n_1$ y $n_2$, con valores observados $y_{ij}$, $i=1,2$, $j=1,\ldots,n_i$. El interés está en
la diferencia de medias $\delta=\mu_1-\mu_2$. El \textbf{modelo lineal} es
\[
  y_{ij} = \mu_i + \varepsilon_{ij},
\]
donde $y_{ij}$ es el $j$-ésimo valor observado en la población $i$, $\mu_i$ es la media de la
población $i$ y los $\varepsilon_{ij}$ son observaciones de una variable aleatoria
$\Normal(0,\sigma^2)$. Obsérvese que \textbf{se supone que las dos poblaciones tienen la misma
varianza} $\sigma^2$ ($\sigma_1^2=\sigma_2^2=\sigma^2$).

\paragraph{Contraste $H_0\colon \mu_1=\mu_2$ frente a $H_1\colon \mu_1\neq\mu_2$.}
\begin{itemize}
  \item \textbf{No restringido}: $y_{ij}=\mu_i+\varepsilon_{ij}$. Minimizando
    $\sum_i\sum_j (y_{ij}-\est{\mu}_i)^2$ se obtienen $\est{\mu}_1=\media{y}_1$ y
    $\est{\mu}_2=\media{y}_2$ (las medias de cada muestra).
  \item \textbf{Restringido} ($H_0$ cierta): $y_{ij}=\mu+\varepsilon_{ij}$, con estimador
    mínimo cuadrático $\est{\mu}=\media{y}$ (la media de todas las observaciones).
\end{itemize}

\begin{proposicion}[Descomposición de la variabilidad (dos medias)][prop:desc-dos-medias]
  $\SC{restringido} = \SC{no restringido} + \SC{hipótesis}$, donde
  \begin{itemize}
    \item $\SC{restringido} = \sum_i\sum_j (y_{ij}-\media{y})^2$, con $n_1+n_2-1$ g.l.;
    \item $\SC{no restringido} = \sum_j (y_{1j}-\media{y}_1)^2 + \sum_j (y_{2j}-\media{y}_2)^2
      = SS_1+SS_2$, con $n_1+n_2-2$ g.l.;
    \item $\SC{hipótesis} = n_1(\media{y}_1-\media{y})^2 + n_2(\media{y}_2-\media{y})^2$, con
      1 g.l. (modelo sin restricciones: 2 parámetros; con restricciones: 1 parámetro).
  \end{itemize}
\end{proposicion}

\paragraph{Medias cuadráticas y estadístico $F$.}
\[
  \MC{no restringido} = \frac{SS_1+SS_2}{n_1+n_2-2}, \qquad
  \MC{hipótesis} = n_1(\media{y}_1-\media{y})^2 + n_2(\media{y}_2-\media{y})^2 .
\]
Ambas estiman $\sigma^2$ si $H_0$ es cierta, y entonces
\[
  F = \frac{\MC{hipótesis}}{\MC{no restringido}} \bajoHo \Fdist{1}{n_1+n_2-2} .
\]
Si $H_0$ no es cierta, $F$ tiende a ser mayor. \textbf{Se rechaza $H_0$ si
$F_{\text{obs}} > F_{1,\,n_1+n_2-2;\,1-\alpha}$}.

\begin{nota}
  $\MC{no restringido}$ es la varianza combinada $s_p^2$ del contraste $t$ para dos muestras
  con varianzas iguales, y $F=t^2$. De nuevo, el enfoque del modelo lineal conduce al contraste
  clásico.
\end{nota}

% ------------------------------------------------------------------------------
\subsection{Inferencias para varias medias: el ANOVA}
\label{sec:t1-anova}

El modelo anterior se generaliza al caso de $k>2$ medias. Suponiendo que los datos son
muestras aleatorias independientes de tamaños $n_i$ de $k$ poblaciones normales, el modelo
lineal es
\[
  y_{ij} = \mu_i + \varepsilon_{ij}, \qquad i=1,\ldots,k,\quad j=1,\ldots,n_i,
\]
donde $y_{ij}$ es el $j$-ésimo valor observado en la población $i$, $\mu_i$ es la media de la
población $i$ y los $\varepsilon_{ij}$ son observaciones de una variable aleatoria
$\Normal(0,\sigma^2)$ (misma varianza en todas las poblaciones).

\begin{definicion}[ANOVA][def:anova]
  Este modelo se conoce como \textbf{análisis de la varianza (ANOVA)}. Se denota por
  $N=\sum_{i=1}^k n_i$ el número total de observaciones.
\end{definicion}

\paragraph{Contraste.}
\[
  H_0\colon \mu_i=\mu_j \ \ \forall i\neq j
  \qquad\text{frente a}\qquad
  H_1\colon \mu_i\neq\mu_j \ \text{para uno o más pares}.
\]
\begin{itemize}
  \item \textbf{No restringido}: minimizando $\sum_i\sum_j (y_{ij}-\est{\mu}_i)^2$ se obtiene
    $\est{\mu}_i=\media{y}_i$, $i=1,\ldots,k$.
  \item \textbf{Restringido}: $y_{ij}=\mu+\varepsilon_{ij}$, con $\est{\mu}=\media{y}$.
\end{itemize}

\begin{proposicion}[Descomposición de la variabilidad (ANOVA)][prop:desc-anova]
  $\SC{restringido} = \SC{no restringido} + \SC{hipótesis}$, donde
  \begin{itemize}
    \item $\SC{restringido} = \sum_i\sum_j (y_{ij}-\media{y})^2$, con $N-1$ g.l.;
    \item $\SC{no restringido} = \sum_j (y_{1j}-\media{y}_1)^2+\cdots+\sum_j (y_{kj}-\media{y}_k)^2
      = SS_1+\cdots+SS_k$, con $N-k$ g.l.;
    \item $\SC{hipótesis} = \sum_i n_i(\media{y}_i-\media{y})^2$, con $k-1$ g.l. (modelo sin
      restricciones: $k$ parámetros; con restricciones: 1 parámetro).
  \end{itemize}
\end{proposicion}

\paragraph{Medias cuadráticas y estadístico $F$.}
\[
  \MC{no restringido} = \frac{SS_1+\cdots+SS_k}{N-k}, \qquad
  \MC{hipótesis} = \frac{\sum_i n_i(\media{y}_i-\media{y})^2}{k-1}.
\]
Ambas estiman $\sigma^2$ si $H_0$ es cierta, y entonces
\[
  F = \frac{\MC{hipótesis}}{\MC{no restringido}} \bajoHo \Fdist{k-1}{N-k} .
\]
Si $H_0$ no es cierta, $F$ tiende a ser mayor. \textbf{Se rechaza $H_0$ si
$F_{\text{obs}} > F_{k-1,\,N-k;\,1-\alpha}$}.

\begin{nota}
  El caso de dos medias es el ANOVA con $k=2$ y $N=n_1+n_2$.
\end{nota}

\begin{resumen}[Resumen de los tres contrastes]
  \centering
  \small
  \renewcommand{\arraystretch}{1.35}
  \begin{tabular}{lccc}
    \toprule
    & \textbf{Una media} & \textbf{Dos medias} & \textbf{$k$ medias (ANOVA)} \\
    \midrule
    $H_0$ & $\mu=\mu_0$ & $\mu_1=\mu_2$ & $\mu_1=\cdots=\mu_k$ \\
    \midrule
    $\SC{restringido}$ & $\sum_i(y_i-\mu_0)^2$ & $\sum_i\sum_j(y_{ij}-\media{y})^2$
      & $\sum_i\sum_j(y_{ij}-\media{y})^2$ \\
    \quad g.l. & $n$ & $n_1+n_2-1$ & $N-1$ \\
    $\SC{no restringido}$ & $\sum_i(y_i-\media{y})^2$ & $SS_1+SS_2$ & $SS_1+\cdots+SS_k$ \\
    \quad g.l. & $n-1$ & $n_1+n_2-2$ & $N-k$ \\
    $\SC{hipótesis}$ & $n(\media{y}-\mu_0)^2$ & $\sum_{i=1}^{2} n_i(\media{y}_i-\media{y})^2$
      & $\sum_{i=1}^{k} n_i(\media{y}_i-\media{y})^2$ \\
    \quad g.l. & $1$ & $1$ & $k-1$ \\
    \midrule
    $F$ bajo $H_0$ & $\Fdist{1}{n-1}$ & $\Fdist{1}{n_1+n_2-2}$ & $\Fdist{k-1}{N-k}$ \\
    \bottomrule
  \end{tabular}
  \par\medskip\raggedright
  En todos los casos $F=\MC{hipótesis}/\MC{no restringido}$ y se rechaza $H_0$ si
  $F_{\text{obs}}$ supera el cuantil $1-\alpha$ de la distribución $F$ correspondiente.
\end{resumen}


% ==============================================================================
%  TEMA 2 — El modelo de regresión simple
%  Fuente: diapositivas T2_RegresionSimple (.pdf / .md)
% ==============================================================================
\chapter{El modelo de regresión simple}
\label{cap:tema2}

% ==============================================================================
%  2.1  El análisis de regresión
% ==============================================================================
\section{El análisis de regresión}
\label{sec:t2-analisis}

El término \textbf{regresión} se utiliza para describir relaciones estadísticas entre
variables. El \textbf{análisis de regresión} es la técnica estadística que permite modelizar e
investigar la relación de una variable $Y$ con otra u otras variables $X_1,\ldots,X_k$
($Y \leftarrow X_1,\ldots,X_k$), es decir, estudiar cómo afectan a $Y$ los cambios en las $X_i$.
\begin{itemize}
  \item $Y$: \textbf{variable respuesta} o \textbf{dependiente}. Es una variable aleatoria y la
    de mayor interés en el problema.
  \item $X_1,\ldots,X_k$: \textbf{regresores}, \textbf{predictores}, variables
    \textbf{independientes} o \textbf{explicativas}. Aportan información sobre la variabilidad
    de $Y$, y el interés está en evaluar el efecto que sus cambios producen en $Y$.
\end{itemize}

\subsection{Componentes del modelo de regresión}

El modelo de regresión debe expresar formalmente las dos componentes de una relación
estadística (\cref{sec:t1-relacion}):
\begin{enumerate}
  \item la tendencia \textbf{sistemática} con la que cambia $Y$ al cambiar los predictores,
    $\vect{X}=(X_1,\ldots,X_k)$;
  \item la componente \textbf{aleatoria}, que recoge la variabilidad con la que los datos se
    distribuyen alrededor de esa tendencia.
\end{enumerate}
Para ello se supone que:
\begin{enumerate}
  \item existe una \textbf{distribución de $Y$ para cada nivel de $\vect{X}$};
  \item la \textbf{media} de estas distribuciones condicionadas \textbf{varía de forma
    sistemática con $\vect{X}$}.
\end{enumerate}
Los modelos son, por tanto, de la forma
\[
  Y = m(X_1,\ldots,X_k) + \varepsilon,
  \qquad
  m(x_1,\ldots,x_k) = \E\left(Y \cond X_1=x_1,\ldots,X_k=x_k\right),
\]
donde $m$ es la función de regresión (\cref{def:funcion-regresion}).

\begin{ejemplo}[Evaluación de empleados][ej:t2-empleados]
  Se quiere evaluar la relación entre la eficiencia de 10 empleados en dos momentos del año.
  La variable respuesta $Y$ es la evaluación de cada trabajador a final de año, y la variable
  explicativa $X$, la evaluación a mitad de año. Para cada valor de $X$ existe una
  distribución de probabilidad de $Y$, y la \emph{curva de regresión} une las medias de esas
  distribuciones.

  \begin{center}
    \begin{tikzpicture}[x=0.09cm, y=0.06cm]
      % ejes
      \draw[-{Stealth}] (40,40) -- (105,40) node[below] {$X$};
      \draw[-{Stealth}] (40,40) -- (40,100) node[left] {$Y$};
      \foreach \t in {50,70,90} {
        \draw (\t,40) -- (\t,39) node[below, font=\scriptsize] {\t};
        \draw (40,\t) -- (39,\t) node[left, font=\scriptsize] {\t};
      }
      \node[below, font=\small] at (72,32) {Evaluación a mitad de año};
      \node[rotate=90, font=\small] at (29,70) {Evaluación final};
      % curva de regresión m(x)
      \draw[NavyBlue, thick, domain=45:100, samples=40]
        plot (\x, {50+0.9*(\x-45)-0.006*(\x-45)^2});
      \node[NavyBlue, font=\scriptsize, anchor=west] at (100,80.8) {curva de regresión};
      % densidades condicionadas (tumbadas)
      \foreach \xx in {55,72,90} {
        \pgfmathsetmacro{\mm}{50+0.9*(\xx-45)-0.006*(\xx-45)^2}
        \draw[gray] (\xx,{\mm-15}) -- (\xx,{\mm+15});
        \draw[colRes, thick, domain=-15:15, samples=40]
          plot ({\xx+9*exp(-(\x)^2/40)}, {\mm+\x});
        \fill[NavyBlue] (\xx,\mm) circle (0.5pt);
      }
      \node[colRes, font=\scriptsize, align=center] at (82,100) {distribución de $Y$\\ para cada $X$};
    \end{tikzpicture}
  \end{center}
\end{ejemplo}

\paragraph{Filosofía.}
\begin{itemize}
  \item La variabilidad de la respuesta $Y$ depende de muchas causas, quizá infinitas.
  \item De ellas, unas pocas, $X_1,\ldots,X_k$, son \emph{causas importantes} que pueden
    observarse y controlarse.
  \item El resto, $X_{k+1},X_{k+2},\ldots$, son muchas otras causas no observables,
    desconocidas o incontrolables.
\end{itemize}
Por tanto, $Y$ puede modelizarse como
\[
  Y = f(X_1,\ldots,X_k) + g(X_{k+1},X_{k+2},\ldots),
\]
donde $g(X_{k+1},X_{k+2},\ldots)=\varepsilon$ es una \textbf{perturbación aleatoria}. Así:
\[
  Y = f(X_1,\ldots,X_k) + \varepsilon .
\]

\subsection{Objetivos y decisiones del análisis de regresión}

\paragraph{Objetivos.} El análisis de regresión tiene varios objetivos:
\begin{enumerate}
  \item \textbf{Conocer el modelo}, para averiguar el tipo o forma de la relación (lineal,
    polinómica\ldots), medir su fuerza y comprender el papel y la importancia de cada variable
    explicativa.
  \item \textbf{Predecir} el valor de $Y$ en un individuo a partir de sus valores de
    $X_1,\ldots,X_k$.
  \item \textbf{Optimizar procesos}: averiguar los valores de $X_1,\ldots,X_k$ que proporcionan
    el valor óptimo de $Y$.
\end{enumerate}

\begin{ejemplo}[Usos de los modelos de regresión]
  \begin{itemize}
    \item Estimar la ganancia de peso de niños debida a varios suplementos dietéticos.
    \item Predecir el resultado académico a partir de las calificaciones de bachillerato y la
      EBAU.
    \item Estimar la cantidad de ventas en relación con los gastos de publicidad.
    \item Predecir el consumo de calefacción según las temperaturas diarias y otros factores
      climáticos.
  \end{itemize}
\end{ejemplo}

\paragraph{Decisiones importantes.}
\begin{itemize}
  \item Selección de las ``mejores'' variables explicativas.
  \item Forma de la componente sistemática de la relación: lineal, cuadrática\ldots
  \item \textbf{Alcance (\emph{scope}) del modelo}. Uno de los usos principales de la
    regresión es la predicción, por lo que hay que tener en cuenta el \emph{rango y la
    dispersión de los valores de $X$ observados} para determinar el rango de validez de las
    predicciones del modelo ajustado. Conviene tenerlo presente ya en la fase de recogida de
    datos.
\end{itemize}

\subsection{El modelo de regresión lineal}

Si $Y=f(X_1,\ldots,X_k)+\varepsilon$, la regresión busca la función $f$ que mejor se ajusta a
los datos. Cuando $f$ es lineal, el modelo es
\[
  Y = \beta_0 + \beta_1X_1 + \beta_2X_2 + \cdots + \beta_kX_k + \varepsilon,
\]
y el objetivo se reduce a \textbf{estimar los coeficientes} $\beta_i$. La asignatura se centra
en los modelos de regresión lineal por varios motivos:
\begin{itemize}
  \item son simples y fáciles de manejar;
  \item la mayoría de los modelos pueden linealizarse mediante transformaciones;
  \item casi cualquier función se puede aproximar localmente por funciones lineales.
\end{itemize}

\begin{ejemplo}[Tipos de regresión lineal][ej:t2-tipos]
  \begin{itemize}
    \item \textbf{Regresión simple} (un solo predictor):
      \begin{itemize}
        \item Número de accidentes de tráfico mortales frente al número de habitantes de cada
          estado: la nube de puntos es claramente creciente y aproximadamente lineal.
        \item La longitud de un caimán es fácil de medir con fotografías aéreas, pero su peso
          es mucho más difícil de estimar, por lo que se quiere predecir el peso a partir de
          la longitud. El peso crece de forma \emph{no lineal} con la longitud, pero de forma
          aproximadamente lineal con el cubo de la longitud; el modelo
          $Y=\beta_0+\beta_1X^3+\varepsilon$ es lineal en los parámetros
          (\cref{def:modelo-lineal}). Es un ejemplo de relación \emph{linealizable}.
      \end{itemize}
    \item \textbf{Regresión lineal múltiple}: consumo de un vehículo (millas por galón) frente
      a su potencia y su peso. La superficie ajustada es un plano.
    \item \textbf{Regresión lineal a trozos}: descripción de la morfología de las glándulas de
      Meibomio mediante una poligonal.
  \end{itemize}
\end{ejemplo}

% ==============================================================================
%  2.2  El modelo de regresión lineal simple: formulación, hipótesis e
%       interpretación de los parámetros
% ==============================================================================
\section{El modelo de regresión lineal simple}
\label{sec:t2-modelo}

El modelo de regresión más sencillo es el \textbf{modelo de regresión lineal simple}, en el
que intervienen:
\begin{itemize}
  \item una variable respuesta $Y$ cuantitativa;
  \item una única variable explicativa $X$ cuantitativa.
\end{itemize}
Se estudia si los cambios en $X$ afectan \emph{de forma lineal} a $Y$ ($Y\leftarrow X$).

\begin{ejemplo}[Conductores: formulación del modelo][ej:t2-conductores]
  En el \cref{ej:conductores} (Tema 1) la edad es la variable explicativa $X$ y la distancia
  máxima de lectura de la señal es la respuesta $Y$. El diagrama de dispersión sugiere una
  relación lineal negativa, por lo que se modeliza la distancia media como una función lineal
  de la edad más un error aleatorio, que es el modelo que se define a continuación.
\end{ejemplo}

\subsection{Formulación del modelo}

\begin{definicion}[Modelo de regresión lineal simple][def:rls]
  \[
    Y_i = \beta_0 + \beta_1 X_i + \varepsilon_i, \qquad i=1,\ldots,n,
  \]
  donde
  \begin{itemize}
    \item $Y_i$ es la respuesta del $i$-ésimo individuo;
    \item $\beta_0$ y $\beta_1$ son \textbf{parámetros}, llamados \textbf{coeficientes de la
      regresión};
    \item $X_i$ es el valor de la variable explicativa (predictor) del $i$-ésimo individuo.
      \textbf{Se asume conocido} (no aleatorio);
    \item $\varepsilon_i$ representa la componente aleatoria: es una \textbf{variable
      aleatoria} $\Normal(0,\sigma^2)$, y $\varepsilon_1,\ldots,\varepsilon_n$ son
      \textbf{i.i.d.}
  \end{itemize}
\end{definicion}

\paragraph{Parámetros del modelo.}
\begin{itemize}
  \item $\beta_0$: \textbf{\emph{intercept}}, u ordenada en el origen.
  \item $\beta_1$: \textbf{pendiente}. $\beta_1=0$ si no hay relación (lineal) entre $X$ e $Y$.
  \item $\sigma^2$: \textbf{varianza del error}.
\end{itemize}

\subsection{Hipótesis del modelo}
\label{sec:t2-hipotesis}

Al ajustar un modelo de regresión lineal simple se asume que los \textbf{errores} son variables
aleatorias independientes e igualmente distribuidas, \textbf{normales con media 0 y varianza
constante}:
\[
  \varepsilon_i \sim \Normal(0,\sigma^2) \ \text{independientes}, \qquad \forall i=1,\ldots,n.
\]
Esto supone cuatro hipótesis, que deben comprobarse para que el modelo sea válido:
\begin{enumerate}
  \item \textbf{Linealidad}: $\E(\varepsilon_i)=0$ (la media de $Y$ es realmente la recta).
  \item \textbf{Homogeneidad de varianzas} (homocedasticidad): $\Var(\varepsilon_i)=\sigma^2$,
    que no depende de $i$.
  \item \textbf{Normalidad}: $\varepsilon_i \sim \Normal$.
  \item \textbf{Independencia}: $\varepsilon_i$ y $\varepsilon_j$ independientes,
    $\forall i\neq j$.
\end{enumerate}

\begin{proposicion}[Distribución de la respuesta][prop:t2-dist-Y]
  Las hipótesis sobre $\varepsilon$ son equivalentes a
  \[
    Y_i \sim \Normal(\beta_0+\beta_1X_i,\ \sigma^2) \ \text{independientes},
    \qquad \forall i=1,\ldots,n .
  \]
\end{proposicion}

En efecto, la respuesta $Y_i=\beta_0+\beta_1X_i+\varepsilon_i$ es la suma de una constante,
$\beta_0+\beta_1X_i$ (los $X_i$ son conocidos), y de la variable aleatoria $\varepsilon_i$:
\begin{itemize}
  \item $Y_i$ es una variable aleatoria normal, por ser combinación lineal de variables
    aleatorias normales;
  \item $\E(Y_i)=\beta_0+\beta_1X_i$, ya que $\E(\varepsilon_i)=0$;
  \item $\Var(Y_i)=\Var(\varepsilon_i)=\sigma^2$, ya que sumar una constante no altera la
    varianza;
  \item las respuestas de individuos distintos son incorreladas: para $i\neq j$,
    \[
      \Cov(Y_i,Y_j)=\Cov(\beta_0+\beta_1X_i+\varepsilon_i,\ \beta_0+\beta_1X_j+\varepsilon_j)
      =\Cov(\varepsilon_i,\varepsilon_j)=0,
    \]
    por la independencia de los errores. Además, $Y_i$ e $Y_j$ son independientes, ya que cada
    una es función únicamente de su propio error y $\varepsilon_i$ y $\varepsilon_j$ son
    independientes.
\end{itemize}

\begin{nota}
  Como los errores $\varepsilon_i$ no se observan, las hipótesis se comprueban sobre los
  \emph{residuos} $e_i$ (\cref{def:t2-ajuste}), mediante el \textbf{análisis residual}.
\end{nota}

\subsection{Interpretación de los parámetros}

La parte determinista del modelo, $\beta_0+\beta_1x_i$, especifica que la \textbf{media
poblacional} de la respuesta $Y$ cuando $X=x_i$ es una \textbf{línea recta}:
\[
  Y_i = \underbrace{\beta_0+\beta_1x_i}_{\E(Y_i\cond X=x_i)\,=\,\E(Y_i)} + \varepsilon_i .
\]
Por ser una recta:
\begin{itemize}
  \item $\beta_0$ (\textbf{\emph{intercept}}) es el valor esperado de la respuesta cuando $X=0$;
  \item $\beta_1$ (\textbf{pendiente}) es el incremento (o decremento, si es negativa) de la
    media de la respuesta cuando $X$ aumenta 1 unidad.
\end{itemize}

\begin{figure}[H]
  \centering
  \begin{tikzpicture}[x=0.11cm, y=0.13cm]
    \draw[-{Stealth}] (0,0) -- (70,0) node[below] {$X$};
    \draw[-{Stealth}] (0,0) -- (0,40) node[left] {$Y$};
    \foreach \t in {10,20,30,40,50,60} \draw (\t,0) -- (\t,-0.6) node[below, font=\scriptsize] {\t};
    \foreach \t in {10,20,30} \draw (0,\t) -- (-0.8,\t) node[left, font=\scriptsize] {\t};
    \draw[NavyBlue, thick] (0,0) -- (68,34) node[right, font=\scriptsize] {$\E(Y)=0+0.5X$};
    \foreach \xx/\lab/\dy in {20/1/2.2,30/3/7,55/2/2.2} {
      \pgfmathsetmacro{\mm}{0.5*\xx}
      \draw[gray] (\xx,{\mm-8}) -- (\xx,{\mm+8});
      \draw[colRes, thick, domain=-8:8, samples=40] plot ({\xx+4.5*exp(-(\x)^2/10)}, {\mm+\x});
      \fill[NavyBlue] (\xx,\mm) circle (1.3pt);
      \node[font=\scriptsize, anchor=east] at (\xx-0.8,\mm+\dy) {$\E(Y_{\lab})=\pgfmathprintnumber{\mm}$};
    }
    \fill (20,7.3) circle (1pt) node[below right, font=\scriptsize] {$Y_1$};
    \fill (55,29.5) circle (1pt) node[below right, font=\scriptsize] {$Y_2$};
    \fill (30,16.4) circle (1pt) node[right, font=\scriptsize] {$Y_3$};
  \end{tikzpicture}
  \caption{Recta de regresión poblacional $\E(Y)=0+0.5X$ y distribución normal de $Y$ para tres
    valores de $X$, todas con la misma varianza. Cada observación $Y_i$ es un valor de su
    distribución.}
\end{figure}

\subsection{Ajustar el modelo: recta estimada, valores ajustados y residuos}

Ajustar el modelo consiste en estimar sus parámetros a partir de los datos: la \textbf{recta de
regresión} ($\beta_0$ y $\beta_1$, que describen el modelo) y $\sigma^2$, que permite cuantificar
el error que se comete.

\begin{definicion}[Recta estimada, valores ajustados y residuos][def:t2-ajuste]
  Sean $\est{\beta}_0$, $\est{\beta}_1$ y $\est{\sigma}^2$ los estimadores de los parámetros.
  \begin{itemize}
    \item La \textbf{recta de regresión estimada} es $\est{Y}=\est{\beta}_0+\est{\beta}_1X$.
    \item Los \textbf{valores ajustados} son $\est{y}_i=\est{\beta}_0+\est{\beta}_1x_i$,
      $i=1,\ldots,n$: los valores estimados por el modelo.
    \item Las desviaciones respecto de la verdadera recta son los \textbf{errores}
      $\varepsilon_i=y_i-(\beta_0+\beta_1x_i)$, que son desconocidos porque también lo son los
      coeficientes. Se estiman mediante los \textbf{residuos}:
      \[
        e_i = y_i-\est{y}_i = y_i-(\est{\beta}_0+\est{\beta}_1x_i).
      \]
  \end{itemize}
\end{definicion}

\begin{figure}[H]
  \centering
  \begin{tikzpicture}
    \begin{axis}[width=8cm, height=5.5cm, axis lines=left, xtick={4.5}, xticklabels={$x_i$},
        ytick={3.2,4.55}, yticklabels={$y_i$,$\est{y}_i$}, xmin=0, xmax=10, ymin=0, ymax=9.5,
        clip=false]
      \addplot[recta, domain=0.3:9.8] {1.4+0.7*x};
      \node[font=\scriptsize, NavyBlue, anchor=west] at (axis cs:9.9,8.3)
        {$\est{y}=\est{\beta}_0+\est{\beta}_1x$};
      \addplot[puntos, mark size=1.5pt] coordinates {(1.5,3.5) (2,2.7) (2.5,2.8) (3.3,3.4)
        (3.7,3.3) (5.3,5.4) (6.2,6.9) (6.8,5.8) (7.8,6.7) (8.6,8.0) (9.3,7.7)};
      \addplot[only marks, mark=*, mark size=1.8pt, red] coordinates {(4.5,3.2)};
      \addplot[only marks, mark=square, mark size=1.8pt, NavyBlue] coordinates {(4.5,4.55)};
      \draw[dashed, gray] (axis cs:0,3.2) -- (axis cs:4.5,3.2);
      \draw[dashed, gray] (axis cs:0,4.55) -- (axis cs:4.5,4.55);
      \draw[dashed, gray] (axis cs:4.5,0) -- (axis cs:4.5,3.2);
      \draw[red, thick] (axis cs:4.5,3.2) -- (axis cs:4.5,4.55)
        node[midway, right, font=\scriptsize] {$e_i$ (residuo)};
      \node[font=\scriptsize, red, anchor=north west] at (axis cs:4.6,3.1)
        {$(x_i,y_i)$ observación};
      \node[font=\scriptsize, NavyBlue, anchor=south east] at (axis cs:4.4,4.6)
        {$(x_i,\est{y}_i)$ valor ajustado};
    \end{axis}
  \end{tikzpicture}
  \caption{Observación, valor ajustado y residuo.}
\end{figure}

Los estimadores tienen una distribución concreta, lo que permite utilizar \textbf{métodos de
inferencia} para evaluar la significación y la precisión de las estimaciones (Tema 3).

\begin{ejemplo}[Conductores: recta ajustada e interpretación][ej:t2-conductores-ajuste]
  En el ejemplo de los conductores, ajustar el modelo es estimar la recta que representa el
  patrón general de los datos. Se obtiene
  \[
    \widehat{\text{distancia}} = 576 - 3\cdot\text{edad},
    \qquad \est{\beta}_0=576,\quad \est{\beta}_1=-3 .
  \]
  \begin{center}
    \begin{tikzpicture}
      \begin{axis}[width=7.5cm, height=5cm, xlabel={Edad (años)}, ylabel={Distancia (pies)},
          xmin=15, xmax=85, ymin=260, ymax=610, label style={font=\small},
          tick label style={font=\footnotesize}]
        \addplot[puntos, mark size=1.3pt] table[x=edad, y=distancia] {datos/conductores.dat};
        \addplot[red, thick, no marks, domain=17:83] {576-3*x};
        \draw[NavyBlue, -{Stealth}] (axis cs:60,260) -- (axis cs:60,396) -- (axis cs:15,396);
        \node[font=\scriptsize, red, anchor=south west] at (axis cs:40,540)
          {$\widehat{\text{dist.}}=576-3\,\text{edad}$};
      \end{axis}
    \end{tikzpicture}
  \end{center}
  \begin{itemize}
    \item \textbf{Predicción}: para un conductor de 60 años, $576-3\cdot60=396\approx400$ pies.
    \item \textbf{\emph{Intercept}} $\est{\beta}_0=576$: la función de regresión vale 576 en $X=0$,
      lo que equivaldría a decir que un conductor de 0 años vería la señal a 576 pies. En este
      caso \textbf{no es interpretable}, ya que la edad no puede ser 0 (y además queda muy
      fuera del rango observado, 18--82 años).
    \item \textbf{Pendiente} $\est{\beta}_1=-3$: por cada año más de edad, la distancia máxima a
      la que se lee la señal \emph{disminuye, en media}, 3 pies.
  \end{itemize}
\end{ejemplo}

\subsection{Versión centrada del modelo}

El modelo también puede formularse utilizando como predictor $X_i-\media{X}$ (desviación del
valor de $X$ en la observación $i$ respecto de su media) en lugar de $X_i$:
\[
  Y_i = \beta_0^* + \beta_1(X_i-\media{X}) + \varepsilon_i,
  \qquad \text{con } \beta_0^* = \beta_0+\beta_1\media{X}.
\]

\begin{demostracion}[Comprobación]
  Sumando y restando $\beta_1\media{X}$ en el modelo original,
  \begin{align*}
    Y_i &= \beta_0 + \beta_1X_i + \varepsilon_i \\
        &= \beta_0 + \beta_1\media{X} - \beta_1\media{X} + \beta_1X_i + \varepsilon_i \\
        &= \underbrace{(\beta_0 + \beta_1\media{X})}_{\beta_0^*}
          + \beta_1(X_i-\media{X}) + \varepsilon_i .
  \end{align*}
  Ambos modelos representan la misma recta, expresada respecto de un origen distinto en el
  eje $X$.
\end{demostracion}

La interpretación de los parámetros es la siguiente:
\begin{itemize}
  \item $\beta_0^*$ (\emph{intercept}) es el valor esperado de la respuesta cuando $X=\media{X}$.
  \item $\beta_1$ (pendiente) es el cambio en la media de la respuesta cuando $X$ se desvía de
    su media 1 unidad. Es \textbf{el mismo parámetro} que en el modelo con predictor $X_i$.
\end{itemize}

\begin{nota}
  En el ejemplo de los conductores el \emph{intercept} $\beta_0$ no era interpretable, mientras que
  $\beta_0^*$ sí lo es: representa la distancia media para un conductor de edad media.
\end{nota}

\subsection{El proceso de modelización}

\begin{figure}[H]
  \centering
  \begin{tikzpicture}[node distance=0.55cm,
      paso/.style={draw, fill=gray!15, minimum width=4cm, minimum height=0.8cm, align=center,
                   font=\small},
      ext/.style={draw, rounded corners=8pt, fill=gray!30, minimum width=2cm,
                  minimum height=0.7cm, font=\small},
      dec/.style={draw, diamond, aspect=2.2, fill=gray!15, align=center, font=\small,
                  inner sep=1pt},
      fl/.style={-{Stealth}}]
    \node[ext] (ini) {Comienzo};
    \node[paso, below=of ini] (eda) {Análisis exploratorio de datos};
    \node[paso, below=of eda] (aju) {Ajuste de posibles\\ modelos de regresión};
    \node[dec, below=of aju] (ok) {¿Se ajustan bien\\ a los datos?};
    \node[paso, left=1.2cm of ok, minimum width=3.3cm] (rev) {Revisar los modelos\\ y plantear nuevos};
    \node[paso, below=of ok] (mej) {Elegir el ``mejor''};
    \node[paso, below=of mej] (inf) {Hacer inferencias\\ con esos modelos};
    \node[ext, below=of inf] (fin) {Fin};
    \draw[fl] (ini) -- (eda); \draw[fl] (eda) -- (aju); \draw[fl] (aju) -- (ok);
    \draw[fl] (ok) -- node[above, font=\small] {No} (rev);
    \draw[fl] (rev) |- ($(aju.south)!0.5!(ok.north)$);
    \draw[fl] (ok) -- node[right, font=\small] {Sí} (mej);
    \draw[fl] (mej) -- (inf); \draw[fl] (inf) -- (fin);
  \end{tikzpicture}
  \caption{Etapas de un análisis de regresión (compárese con las etapas del ajuste de un modelo
    estadístico, \cref{sec:t1-modelos}).}
\end{figure}

% ==============================================================================
%  2.3  Estimación por mínimos cuadrados. Propiedades de los estimadores
% ==============================================================================
\section{Estimación por mínimos cuadrados}
\label{sec:t2-mc}

``Ajustar'' la recta de regresión es encontrar los \textbf{mejores estimadores} de los
coeficientes $\beta_0$ y $\beta_1$, para lo que es necesario fijar un criterio que permita
comparar rectas.

\begin{ejemplo}[Dos rectas para tres puntos][ej:t2-tres-puntos]
  Se dispone de una muestra de $n=3$ puntos, $(20,5)$, $(55,12)$ y $(30,10)$, y se consideran
  dos rectas:
  \[
    M_1\colon\ \est{Y} = 9 + 0\cdot X,
    \qquad
    M_2\colon\ \est{Y} = 2.81 + 0.177X .
  \]
  \begin{center}
    \begin{minipage}[t]{0.46\linewidth}
      \centering
      \begin{tikzpicture}
        \begin{axis}[width=5.5cm, height=4.5cm, xmin=0, xmax=70, ymin=0, ymax=14,
            xtick={20,40,60}, ytick={3,6,9,12}, tick label style={font=\scriptsize}]
          \addplot[recta, domain=5:68] {9};
          \addplot[puntos, mark size=1.6pt] coordinates {(20,5) (55,12) (30,10)};
          \draw[red] (axis cs:20,5) -- (axis cs:20,9);
          \draw[red] (axis cs:55,12) -- (axis cs:55,9);
          \draw[red] (axis cs:30,10) -- (axis cs:30,9);
        \end{axis}
      \end{tikzpicture}\\[2pt]
      {\small Recta $M_1$: $Q_1=26$.}
    \end{minipage}\hfill
    \begin{minipage}[t]{0.46\linewidth}
      \centering
      \begin{tikzpicture}
        \begin{axis}[width=5.5cm, height=4.5cm, xmin=0, xmax=70, ymin=0, ymax=14,
            xtick={20,40,60}, ytick={3,6,9,12}, tick label style={font=\scriptsize}]
          \addplot[recta, domain=5:68] {2.8077+0.17692*x};
          \addplot[puntos, mark size=1.6pt] coordinates {(20,5) (55,12) (30,10)};
          \draw[red] (axis cs:20,5) -- (axis cs:20,6.346);
          \draw[red] (axis cs:55,12) -- (axis cs:55,12.538);
          \draw[red] (axis cs:30,10) -- (axis cs:30,8.115);
        \end{axis}
      \end{tikzpicture}\\[2pt]
      {\small Recta $M_2$: $Q_2=5.7$.}
    \end{minipage}
  \end{center}
  Los segmentos rojos son las desviaciones $y_i-\est{y}_i$. La suma de sus cuadrados es:
  \[
    Q_1 = (-4)^2+3^2+1^2 = 26,
    \qquad
    Q_2 = (-1.35)^2+(-0.54)^2+1.88^2 \approx 5.7 .
  \]
  Existen infinitas rectas que podrían ajustarse, una por cada par $(\beta_0,\beta_1)$. Según
  el criterio de mínimos cuadrados (\cref{def:t2-mc}) es preferible la de menor suma de
  cuadrados de las desviaciones, $M_2$ ($Q_2=5.7<26=Q_1$). De hecho, $M_2$ es la recta de
  mínimos cuadrados de estos datos (\cref{ej:t2-tres-puntos-mc}).
\end{ejemplo}

\subsection{El criterio de mínimos cuadrados}

\begin{definicion}[Estimadores de mínimos cuadrados][def:t2-mc]
  Dadas las observaciones $(x_i,y_i)$, las desviaciones de $y_i$ respecto de su valor esperado
  son $y_i-(\beta_0+\beta_1x_i)$. Los \textbf{estimadores mínimo cuadráticos} (MC) de $\beta_0$
  y $\beta_1$ son los valores $\est{\beta}_0$, $\est{\beta}_1$ que minimizan la suma de los
  cuadrados de las $n$ desviaciones:
  \[
    Q(\beta_0,\beta_1) = \sumi (y_i-\beta_0-\beta_1x_i)^2,
    \qquad
    (\est{\beta}_0,\est{\beta}_1) = \arg\min_{\beta_0,\beta_1} Q(\beta_0,\beta_1).
  \]
\end{definicion}

\paragraph{Ecuaciones normales.} Derivando $Q$ e igualando a cero:
\begin{align*}
  \frac{\partial Q}{\partial\beta_0} &= -2\sumi (y_i-\beta_0-\beta_1x_i) = 0
  &&\Longrightarrow\quad \sumi y_i = n\est{\beta}_0 + \est{\beta}_1\sumi x_i, \\
  \frac{\partial Q}{\partial\beta_1} &= -2\sumi x_i(y_i-\beta_0-\beta_1x_i) = 0
  &&\Longrightarrow\quad \sumi x_iy_i = \est{\beta}_0\sumi x_i + \est{\beta}_1\sumi x_i^2 .
\end{align*}
Dividiendo la primera entre $n$ se obtiene $\est{\beta}_0$; sustituyéndolo en la segunda,
$\sumi x_iy_i - n\media{x}\,\media{y} = \est{\beta}_1\left(\sumi x_i^2 - n\media{x}^2\right)$,
y como $\sumi x_iy_i - n\media{x}\,\media{y}=SS_{XY}$ y $\sumi x_i^2 - n\media{x}^2=SS_X$:

\begin{teorema}[Estimadores de mínimos cuadrados][teo:t2-estimadores-mc]
  \[
    \boxed{\ \est{\beta}_0 = \media{y} - \est{\beta}_1\media{x}\ }
    \qquad\qquad
    \boxed{\ \est{\beta}_1 = \frac{\sumi (x_i-\media{x})(y_i-\media{y})}{\sumi (x_i-\media{x})^2}
    = \frac{SS_{XY}}{SS_X}\ }
  \]
  donde $\media{x}$ e $\media{y}$ son las medias de las $x_i$ y de las $y_i$ y $SS_X=SS_{xx}$,
  $SS_{XY}=SS_{xy}$ con la notación de la \cref{def:SS-MS}.
\end{teorema}

\begin{ejemplo}[Tres puntos: recta de mínimos cuadrados][ej:t2-tres-puntos-mc]
  Para los datos del \cref{ej:t2-tres-puntos}: $\media{x}=35$, $\media{y}=9$,
  \[
    SS_X = (-15)^2+20^2+(-5)^2 = 650, \qquad
    SS_{XY} = (-15)(-4)+(20)(3)+(-5)(1) = 115,
  \]
  así que $\est{\beta}_1 = 115/650 = 0.177$ y $\est{\beta}_0 = 9-0.177\cdot35 = 2.81$: es la
  recta $M_2$.
\end{ejemplo}

\subsection{Propiedades de los estimadores: teorema de Gauss--Markov}

La justificación de que los estimadores mínimo cuadráticos son los ``mejores'' la proporciona
el siguiente teorema.

\begin{teorema}[Gauss--Markov][teo:gauss-markov]
  En las condiciones del modelo de regresión, los estimadores mínimo cuadráticos
  $\est{\beta}_0$ y $\est{\beta}_1$:
  \begin{itemize}
    \item son \textbf{insesgados}: $\E[\est{\beta}_0]=\beta_0$ y $\E[\est{\beta}_1]=\beta_1$;
    \item tienen \textbf{mínima varianza} entre todos los estimadores \textbf{lineales
      insesgados} de $\beta_0$ y $\beta_1$.
  \end{itemize}
\end{teorema}

El significado de cada término del teorema es el siguiente:
\begin{itemize}
  \item \emph{Lineal}: el estimador es combinación lineal de las observaciones,
    $\sum_i c_iY_i$ con $c_i$ constantes. Los estimadores MC lo son; por ejemplo,
    $\est{\beta}_1=\sum_i k_iY_i$ con $k_i=(x_i-\media{x})/SS_X$.
  \item \emph{Insesgado}: su esperanza, sobre todas las muestras posibles, coincide con el
    parámetro.
  \item \emph{Mínima varianza}: entre todos los estimadores lineales insesgados, los MC son los
    de menor variabilidad muestral, es decir, los más precisos.
\end{itemize}
Por ello se denominan \textbf{mejores estimadores lineales insesgados} (ELIO; en inglés,
BLUE: \emph{Best Linear Unbiased Estimators}).

\paragraph{Estimación de la respuesta media.} A partir de los coeficientes estimados se obtiene
la \textbf{recta de regresión estimada} $\est{Y}=\est{\beta}_0+\est{\beta}_1X$, donde $\est{Y}$ es el
valor estimado de la recta para un valor $X$ del regresor.
\begin{itemize}
  \item $\E(Y)=\beta_0+\beta_1X$ es la media de la distribución de $Y$ condicionada al valor
    $X$: la \textbf{respuesta media}.
  \item $\est{Y}$ es un estimador puntual de $\E(Y)$, y es el \textbf{estimador insesgado de
    varianza mínima en la clase de los estimadores lineales insesgados} (extensión del teorema
    de Gauss--Markov).
\end{itemize}

En consecuencia, cada valor ajustado $\est{y}_i=\est{\beta}_0+\est{\beta}_1x_i$ es la
\textbf{estimación de la respuesta media} $\E(Y\cond X=x_i)$, es decir, el valor esperado de la
respuesta para los individuos con $X=x_i$ según el modelo ajustado. No es una predicción exacta
de $y_i$: ambos difieren en el residuo $e_i$. En el ejemplo de los conductores,
$\est{y}=396$ es la distancia media estimada para los conductores de 60 años.

\subsection{Propiedades de la recta de regresión ajustada}

\begin{proposicion}[Propiedades de la recta de regresión][prop:t2-propiedades-recta]
  Para la recta ajustada por mínimos cuadrados:
  \begin{enumerate}
    \item La suma de los residuos es 0: $\sumi e_i=0$.
    \item La suma de los residuos al cuadrado, $\sumi e_i^2$, es mínima.
    \item La suma de los valores observados es igual a la de los ajustados:
      $\sumi y_i=\sumi \est{y}_i$.
    \item $\sumi x_ie_i=0$.
    \item $\sumi \est{y}_ie_i=0$.
    \item La recta de regresión pasa por el punto $(\media{x},\media{y})$.
  \end{enumerate}
\end{proposicion}

\begin{demostracion}
  Se parte de que $e_i=y_i-\est{\beta}_0-\est{\beta}_1x_i$ y de que $\est{\beta}_0$ y
  $\est{\beta}_1$ cumplen las ecuaciones normales:
  \[
    \text{(EN1)}\ \ \sumi (y_i-\est{\beta}_0-\est{\beta}_1x_i)=0,
    \qquad
    \text{(EN2)}\ \ \sumi x_i(y_i-\est{\beta}_0-\est{\beta}_1x_i)=0 .
  \]
  \begin{enumerate}
    \item Es exactamente (EN1): $\sumi e_i=0$.
    \item $\sumi e_i^2 = Q(\est{\beta}_0,\est{\beta}_1)$, y $(\est{\beta}_0,\est{\beta}_1)$ es,
      por definición (\cref{def:t2-mc}), el punto que minimiza $Q$.
    \item Como $y_i=\est{y}_i+e_i$, sumando y usando 1:
      $\sumi y_i=\sumi\est{y}_i+\sumi e_i=\sumi\est{y}_i$.
    \item Es exactamente (EN2): $\sumi x_ie_i=0$.
    \item Usando 1 y 4:
      $\sumi \est{y}_ie_i=\sumi(\est{\beta}_0+\est{\beta}_1x_i)e_i
      =\est{\beta}_0\sumi e_i+\est{\beta}_1\sumi x_ie_i = 0$.
    \item Evaluando la recta en $x=\media{x}$ y usando
      $\est{\beta}_0=\media{y}-\est{\beta}_1\media{x}$:
      $\est{\beta}_0+\est{\beta}_1\media{x}=\media{y}-\est{\beta}_1\media{x}+\est{\beta}_1\media{x}
      =\media{y}$.
  \end{enumerate}
\end{demostracion}

\subsection{Estimación de \texorpdfstring{$\sigma^2$}{sigma cuadrado}}
\label{sec:t2-sigma}

\paragraph{Caso general.} Sea $Y$ una variable aleatoria con varianza $\sigma^2$. Un buen
estimador de $\sigma^2$ es la varianza muestral corregida, que es una suma de cuadrados
dividida por sus grados de libertad (\cref{def:SS-MS}):
\[
  \est{\sigma}^2 = S_Y^2 = \frac{1}{n-1}\sumi (y_i-\media{y})^2 = \frac{SS_Y}{n-1}.
\]
Se divide entre $n-1$ porque \textbf{se pierde un grado de libertad al estimar la media
poblacional} mediante $\media{Y}$.

\begin{nota}[Grados de libertad]
  Las $n$ desviaciones $y_i-\media{y}$ no son libres, ya que siempre suman 0: conocidas $n-1$
  de ellas, la última queda determinada. En general, cada parámetro que se estima a partir de
  los datos para calcular las desviaciones resta un grado de libertad.
\end{nota}

\paragraph{En el modelo de regresión.}
\begin{itemize}
  \item Cada $Y_i$ es una variable aleatoria con varianza $\sigma^2$ (la misma que el error
    $\varepsilon_i$), por lo que $\sigma^2$ se estima como una suma de cuadrados dividida por
    sus grados de libertad.
  \item Ahora, sin embargo, \textbf{la media de cada $Y_i$ depende de $X_i$}, de modo que la
    desviación de cada $y_i$ debe calcularse respecto de \emph{su} media estimada, $\est{y}_i$;
    es decir, a partir de los residuos $e_i=y_i-\est{y}_i$.
  \item La suma de cuadrados correspondiente es la \textbf{suma de cuadrados del error}:
    \[
      SSE = \sumi (y_i-\est{y}_i)^2 = \sumi e_i^2 .
    \]
  \item $SSE$ tiene $n-2$ grados de libertad: se pierden 2 al estimar $\beta_0$ y $\beta_1$,
    necesarios para obtener los $\est{y}_i$.
\end{itemize}

\begin{proposicion}[Estimador insesgado de $\sigma^2$][prop:t2-mse]
  Un estimador insesgado de $\sigma^2$ es la \textbf{media cuadrática del error}:
  \[
    MSE = \frac{SSE}{n-2} = \frac{1}{n-2}\sumi (y_i-\est{y}_i)^2 = \frac{1}{n-2}\sumi e_i^2 .
  \]
\end{proposicion}

% ==============================================================================
%  2.4  Estimación por máxima verosimilitud. Propiedades de los estimadores
% ==============================================================================
\section{Estimación por máxima verosimilitud}
\label{sec:t2-mv}

Si además se utiliza que los errores siguen una distribución \textbf{normal}, los parámetros
$\beta_0$, $\beta_1$ y $\sigma^2$ pueden estimarse por el \textbf{método de máxima
verosimilitud} (MV).

\paragraph{Verosimilitud.} En el modelo de regresión lineal simple,
$Y_i\sim\Normal(\beta_0+\beta_1X_i,\sigma^2)$ (\cref{prop:t2-dist-Y}), con densidad
\[
  f_i = \frac{1}{\sqrt{2\pi}\,\sigma}
        \exp\left\{-\frac12\left(\frac{Y_i-(\beta_0+\beta_1X_i)}{\sigma}\right)^2\right\}.
\]
Como las $Y_i$ son independientes, la \textbf{función de verosimilitud} de la muestra
$(Y_1,\ldots,Y_n)$ es el producto de las densidades:
\begin{align*}
  L(\beta_0,\beta_1,\sigma^2)
  &= \prod_{i=1}^{n} \frac{1}{(2\pi\sigma^2)^{1/2}}
     \exp\left\{-\frac{1}{2\sigma^2}(Y_i-\beta_0-\beta_1X_i)^2\right\} \\
  &= \frac{1}{(2\pi\sigma^2)^{n/2}}
     \exp\left\{-\frac{1}{2\sigma^2}\sumi (Y_i-\beta_0-\beta_1X_i)^2\right\}.
\end{align*}
La \textbf{log-verosimilitud} es
\[
  \ln L(\beta_0,\beta_1,\sigma^2)
  = -\frac{n}{2}\ln(2\pi) - \frac{n}{2}\ln(\sigma^2)
    - \frac{1}{2\sigma^2}\sumi (Y_i-\beta_0-\beta_1X_i)^2 .
\]

\paragraph{Maximización.} Se buscan los valores de $\beta_0$, $\beta_1$ y $\sigma^2$ que
maximizan $L$ o, equivalentemente, $\ln L$. Derivando respecto de cada parámetro e igualando a
0 se obtiene:
\begin{align*}
  \frac{\partial \ln L}{\partial\beta_0} &= \frac{1}{\sigma^2}\sumi (Y_i-\beta_0-\beta_1X_i) = 0
  &&\Longrightarrow\ \sumi \big(Y_i-\est{\beta}_0-\est{\beta}_1X_i\big) = 0, \\
  \frac{\partial \ln L}{\partial\beta_1} &= \frac{1}{\sigma^2}\sumi X_i(Y_i-\beta_0-\beta_1X_i) = 0
  &&\Longrightarrow\ \sumi X_i\big(Y_i-\est{\beta}_0-\est{\beta}_1X_i\big) = 0, \\
  \frac{\partial \ln L}{\partial\sigma^2} &= -\frac{n}{2\sigma^2}
     + \frac{1}{2\sigma^4}\sumi (Y_i-\beta_0-\beta_1X_i)^2 = 0
  &&\Longrightarrow\ \frac{n}{2\est{\sigma}^2}
     = \frac{1}{2\est{\sigma}^4}\sumi \big(Y_i-\est{\beta}_0-\est{\beta}_1X_i\big)^2 .
\end{align*}

\begin{nota}
  Las dos primeras ecuaciones coinciden con \textbf{las ecuaciones normales} de mínimos
  cuadrados. Para cualquier $\sigma^2$ fijo, maximizar $\ln L$ en $(\beta_0,\beta_1)$ equivale a
  minimizar $\sum (Y_i-\beta_0-\beta_1X_i)^2=Q$, ya que es el único término que depende de
  ellos y aparece con signo negativo.
\end{nota}

\begin{teorema}[Estimadores de máxima verosimilitud][teo:t2-emv]
  \begin{itemize}
    \item De las dos primeras ecuaciones:
      \[
        \est{\beta}_0 = \media{y}-\est{\beta}_1\media{x},
        \qquad
        \est{\beta}_1 = \frac{\sumi (x_i-\media{x})(y_i-\media{y})}{\sumi (x_i-\media{x})^2}
        = \frac{SS_{XY}}{SS_X}.
      \]
      Son \textbf{los mismos estimadores que los de mínimos cuadrados}
      (\cref{teo:t2-estimadores-mc}) y, por tanto, tienen sus mismas propiedades.
    \item De la tercera:
      \[
        \est{\sigma}^2 = \frac1n\sumi \big(y_i-\est{\beta}_0-\est{\beta}_1x_i\big)^2
        = \frac1n\sumi e_i^2 = \frac{SSE}{n}.
      \]
  \end{itemize}
\end{teorema}

\paragraph{El estimador de $\sigma^2$.}
\begin{itemize}
  \item El estimador máximo verosímil de $\sigma^2$ es \textbf{sesgado}: como
    $\est{\sigma}^2=\frac{n-2}{n}MSE$ y $MSE$ es insesgado (\cref{prop:t2-mse}),
    $\E[\est{\sigma}^2]=\frac{n-2}{n}\sigma^2<\sigma^2$ (subestima $\sigma^2$).
  \item Habitualmente se usa $MSE$ como estimador insesgado:
    $MSE=\dfrac{n}{n-2}\,\est{\sigma}^2$.
  \item Si $n$ es grande, $\frac{n}{n-2}\approx1$ y la diferencia entre ambos es pequeña.
\end{itemize}

\begin{resumen}[Resumen: mínimos cuadrados frente a máxima verosimilitud]
  \centering
  \small
  \renewcommand{\arraystretch}{1.5}
  \begin{tabular}{lcc}
    \toprule
    & \textbf{Mínimos cuadrados} & \textbf{Máxima verosimilitud} \\
    \midrule
    Requiere normalidad de los errores & No & Sí \\
    $\est{\beta}_1$ & $SS_{XY}/SS_X$ & $SS_{XY}/SS_X$ \\
    $\est{\beta}_0$ & $\media{y}-\est{\beta}_1\media{x}$ & $\media{y}-\est{\beta}_1\media{x}$ \\
    $\sigma^2$ & $MSE=\dfrac{SSE}{n-2}$ (insesgado) & $\est{\sigma}^2=\dfrac{SSE}{n}$ (sesgado) \\
    \bottomrule
  \end{tabular}
\end{resumen}


% ==============================================================================
%  TEMA 3 — Inferencias en regresión
%  Fuente: diapositivas T3_Inferencias_regresion.pdf
% ==============================================================================
\chapter{Inferencias en regresión}
\label{cap:tema3}

% ==============================================================================
%  3.1  Inferencias sobre los coeficientes de regresión
% ==============================================================================
\section{Inferencias sobre los coeficientes de regresión}
\label{sec:t3-coeficientes}

En todo el tema se supone que es aplicable el modelo de regresión lineal simple con errores
normales (\cref{def:rls}):
\[
  Y_i = \beta_0 + \beta_1X_i + \varepsilon_i, \qquad
  \varepsilon_i \iid \Normal(0,\sigma^2), \quad i=1,\ldots,n .
\]
Hacer inferencias sobre $\beta_0$ y $\beta_1$ consiste en cuantificar la incertidumbre de sus
estimaciones mediante intervalos de confianza y en resolver contrastes de hipótesis sobre sus
valores. Los contrastes de mayor interés son los de la forma $H_0\colon\beta_i=0$, y en
particular
\[
  H_0\colon \beta_1 = 0 \qquad\text{frente a}\qquad H_1\colon \beta_1\neq 0 .
\]
Si $\beta_1=0$, entonces $\E(Y)=\beta_0+0\cdot X=\beta_0$: la recta de regresión es
horizontal y la media de $Y$ no varía con $X$. Más aún, bajo el modelo las $Y_i$ son
independientes con $Y_i\sim\Normal(\beta_0+\beta_1X_i,\sigma^2)$ (\cref{prop:t2-dist-Y}); si
$\beta_1=0$, todas siguen la misma distribución $\Normal(\beta_0,\sigma^2)$, que no depende de
$X$, de modo que \textbf{no hay relación de ningún tipo entre $X$ e $Y$}.

Esta conclusión presupone que el modelo es correcto. Si falla la hipótesis de linealidad, $Y$
puede depender de $X$ de forma no lineal (por ejemplo, cuadrática) aunque la recta que mejor se
ajusta a los datos tenga pendiente nula; por eso es necesario validar el modelo
(\cref{sec:t4-objetivos}).

Para resolver estos contrastes y construir los intervalos es necesario conocer la distribución
en el muestreo de los estimadores $\est{\beta}_0$ y $\est{\beta}_1$, es decir, la distribución de
los valores que toman en muestras repetidas con los valores de $X$ fijos.

\subsection{Distribución en el muestreo de \texorpdfstring{$\est{\beta}_1$}{beta1}}

\begin{proposicion}[$\est{\beta}_1$ como combinación lineal de las $Y_i$][prop:t3-ki]
  El estimador $\est{\beta}_1 = SS_{XY}/\SSX$ puede escribirse como
  \[
    \est{\beta}_1 = \sumi k_iY_i,
    \qquad
    k_i = \frac{X_i-\media{X}}{\sumi (X_i-\media{X})^2} = \frac{X_i-\media{X}}{\SSX},
  \]
  donde los $k_i$ son constantes que dependen únicamente de las $X_i$. Además:
  \[
    \text{(i) } \sumi k_i = 0, \qquad
    \text{(ii) } \sumi k_iX_i = 1, \qquad
    \text{(iii) } \sumi k_i^2 = \frac{1}{\SSX}.
  \]
\end{proposicion}

\begin{demostracion}
  \emph{Expresión de $\est{\beta}_1$.} La suma de las desviaciones respecto de la media es nula:
  \[
    \sumi (X_i-\media{X}) = \sumi X_i - n\media{X} = n\media{X}-n\media{X} = 0 .
  \]
  Por tanto,
  \begin{align*}
    SS_{XY} &= \sumi (X_i-\media{X})(Y_i-\media{Y}) \\
            &= \sumi (X_i-\media{X})Y_i - \media{Y}\sumi (X_i-\media{X}) \\
            &= \sumi (X_i-\media{X})Y_i ,
  \end{align*}
  y dividiendo por $\SSX$,
  \begin{align*}
    \est{\beta}_1 &= \frac{SS_{XY}}{\SSX}
      = \sumi \frac{X_i-\media{X}}{\SSX}\,Y_i \\
    &= \sumi k_iY_i .
  \end{align*}

  \emph{Propiedad (i).}
  \[
    \sumi k_i = \frac{1}{\SSX}\sumi (X_i-\media{X}) = 0 .
  \]

  \emph{Propiedad (ii).} Restando el término nulo $\media{X}\displaystyle\sumi (X_i-\media{X})$,
  \begin{align*}
    \sumi k_iX_i &= \frac{1}{\SSX}\sumi (X_i-\media{X})X_i \\
    &= \frac{1}{\SSX}\left(\sumi (X_i-\media{X})X_i - \media{X}\sumi (X_i-\media{X})\right) \\
    &= \frac{1}{\SSX}\sumi (X_i-\media{X})(X_i-\media{X}) \\
    &= \frac{\SSX}{\SSX} = 1 .
  \end{align*}

  \emph{Propiedad (iii).}
  \begin{align*}
    \sumi k_i^2 &= \sumi \frac{(X_i-\media{X})^2}{\SSX^2} \\
    &= \frac{1}{\SSX^2}\sumi (X_i-\media{X})^2 \\
    &= \frac{\SSX}{\SSX^2} = \frac{1}{\SSX} .
  \end{align*}
\end{demostracion}

\begin{teorema}[Distribución de $\est{\beta}_1$][teo:t3-dist-b1]
  \[
    \est{\beta}_1 \sim \Normal\!\left(\beta_1,\ \frac{\sigma^2}{\SSX}\right),
    \qquad
    \frac{\est{\beta}_1-\beta_1}{\sqrt{MSE/\SSX}} \sim \tdist{n-2} .
  \]
\end{teorema}

\begin{demostracion}
  \emph{Normalidad.} Por la \cref{prop:t3-ki}, $\est{\beta}_1$ es combinación lineal de las
  $Y_i$, que son normales independientes; luego es normal.

  \emph{Insesgadez.} Usando las propiedades (i) y (ii) de los $k_i$,
  \begin{align*}
    \E(\est{\beta}_1) &= \E\left(\sumi k_iY_i\right) = \sumi k_i\E(Y_i) \\
    &= \sumi k_i(\beta_0+\beta_1X_i) \\
    &= \beta_0\sumi k_i + \beta_1\sumi k_iX_i \\
    &= \beta_0\cdot0+\beta_1\cdot1 = \beta_1 .
  \end{align*}

  \emph{Varianza.} Las $Y_i$ son independientes con varianza $\sigma^2$ y los $k_i$ son
  constantes, así que las covarianzas se anulan; usando la propiedad (iii),
  \begin{align*}
    \Var(\est{\beta}_1) &= \Var\left(\sumi k_iY_i\right) = \sumi k_i^2\Var(Y_i) \\
    &= \sigma^2\sumi k_i^2 \\
    &= \frac{\sigma^2}{\SSX} .
  \end{align*}
  Al estimar $\sigma^2$ por $MSE$, insesgado (\cref{prop:t2-mse}), se obtiene el estimador
  insesgado $\widehat{\Var}(\est{\beta}_1)=MSE/\SSX$.

  \emph{Distribución $t$.} Se utiliza que
  \[
    \frac{(n-2)MSE}{\sigma^2} = \frac{SSE}{\sigma^2} \sim \chisq{n-2}
  \]
  y que $SSE$ es independiente de $\est{\beta}_0$ y $\est{\beta}_1$ (resultado que se admite sin
  demostración). Entonces
  \begin{align*}
    \frac{\est{\beta}_1-\beta_1}{\sqrt{MSE/\SSX}}
    &= \frac{\dfrac{\est{\beta}_1-\beta_1}{\sqrt{\sigma^2/\SSX}}}
            {\sqrt{\dfrac{MSE}{\sigma^2}}}
     = \frac{\dfrac{\est{\beta}_1-\beta_1}{\sqrt{\sigma^2/\SSX}}}
            {\sqrt{\dfrac{SSE/\sigma^2}{n-2}}} ,
  \end{align*}
  que es el cociente entre una $\Normal(0,1)$ y la raíz de una $\chisq{n-2}$ independiente
  dividida por sus grados de libertad; luego sigue una $\tdist{n-2}$.
\end{demostracion}

\subsection{Intervalo de confianza para \texorpdfstring{$\beta_1$}{beta1}}

Del \cref{teo:t3-dist-b1},
\[
  P\left(\tq{n-2}{\alpha/2} \le \frac{\est{\beta}_1-\beta_1}{\sqrt{MSE/\SSX}}
  \le \tq{n-2}{1-\alpha/2}\right) = 1-\alpha .
\]
Por la simetría de la $t$ de Student, $\tq{n-2}{\alpha/2}=-\tq{n-2}{1-\alpha/2}$; despejando
$\beta_1$,
\[
  1-\alpha = P\left(\est{\beta}_1-\tq{n-2}{1-\alpha/2}\sqrt{\frac{MSE}{\SSX}}
  \le \beta_1 \le
  \est{\beta}_1+\tq{n-2}{1-\alpha/2}\sqrt{\frac{MSE}{\SSX}}\right).
\]

\begin{resumen}[Intervalo de confianza de nivel $1-\alpha$ para $\beta_1$]
  \[
    \beta_1 \in \left(\est{\beta}_1 \pm \tq{n-2}{1-\alpha/2}\sqrt{\frac{MSE}{\SSX}}\right).
  \]
  La cantidad $\sqrt{MSE/\SSX}$ es el \textbf{error estándar} de $\est{\beta}_1$.
\end{resumen}

\begin{ejemplo}[Volumen de madera][ej:arboles]
  Se quiere estimar la cosecha potencial de madera de un bosque. En una muestra de árboles se
  toman medidas no destructivas: HT (altura, en pies), DBH (diámetro del tronco a la altura del
  pecho, en pulgadas) y D16 (diámetro del tronco a 16 pies de altura, en pulgadas), con el fin
  de predecir el volumen de madera VOL (pies cúbicos) de cada árbol. Con $n=20$ árboles se
  ajusta un modelo de regresión lineal simple con respuesta VOL y regresor DBH, por ser la
  medida más fácil de obtener:
  \[
    \est{Y} = -45.57 + 6.916\,X, \qquad MSE = 36.587, \qquad \SSX = 64.512 .
  \]
  \begin{center}
    \begin{tikzpicture}
      \begin{axis}[width=7cm, height=5cm, xlabel={DBH (pulgadas)}, ylabel={VOL (pies$^3$)},
          xmin=9.5, xmax=20, ymin=15, ymax=105, label style={font=\small},
          tick label style={font=\footnotesize}]
        \addplot[puntos, mark size=1.3pt] table[x=dbh, y=vol] {datos/tema3/arboles.dat};
        \addplot[recta, domain=10:19.3] {-45.57+6.916*x};
      \end{axis}
    \end{tikzpicture}\\
    {\footnotesize (Datos reconstruidos a partir del gráfico de las diapositivas.)}
  \end{center}
  \textbf{IC del 95\,\% para $\beta_1$.} Con $\tq{18}{0.975}=2.101$ y error estándar
  $\sqrt{36.587/64.512}=0.7531$,
  \[
    \beta_1 \in (6.916 \pm 2.101\cdot0.7531) = (6.916\pm1.582) = (5.334,\ 8.498).
  \]
  Con una confianza del 95\,\%, el incremento medio poblacional del volumen de madera por cada
  pulgada de aumento del diámetro a la altura del pecho está entre 5.334 y 8.498 pies cúbicos.
  Como $0\notin(5.334,\,8.498)$, se rechaza $H_0\colon\beta_1=0$ a nivel $\alpha=0.05$: existe
  relación lineal entre DBH y VOL.
\end{ejemplo}

\subsection{Contrastes de hipótesis para \texorpdfstring{$\beta_1$}{beta1}}

Todos se basan en el estadístico del \cref{teo:t3-dist-b1}.

\paragraph{Contraste bilateral.} El más relevante es $H_0\colon\beta_1=0$ frente a
$H_1\colon\beta_1\neq0$, ya que si no se rechaza $H_0$ se concluye que no hay relación lineal
entre $X$ e $Y$. Se calcula el valor observado del estadístico
\[
  t^* = \frac{\est{\beta}_1-0}{\sqrt{MSE/\SSX}}
\]
y se compara con el cuantil de la $\tdist{n-2}$:
\begin{itemize}
  \item si $|t^*|\le\tq{n-2}{1-\alpha/2}$, no se rechaza $H_0$ a nivel $\alpha$;
  \item si $|t^*|>\tq{n-2}{1-\alpha/2}$, se rechaza $H_0$ a nivel $\alpha$.
\end{itemize}

\paragraph{Contraste unilateral.} Para $H_0\colon\beta_1\le0$ frente a $H_1\colon\beta_1>0$ se usa
el mismo estadístico, comparándolo ahora con el cuantil $1-\alpha$:
\begin{itemize}
  \item si $t^*\le\tq{n-2}{1-\alpha}$, no se rechaza $H_0$ a nivel $\alpha$;
  \item si $t^*>\tq{n-2}{1-\alpha}$, se rechaza $H_0$ a nivel $\alpha$.
\end{itemize}
(Análogamente para $H_1\colon\beta_1<0$, rechazando si $t^*<-\tq{n-2}{1-\alpha}$.)

\paragraph{Valor específico $b_1$.} Para $H_0\colon\beta_1=b_1$ frente a $H_1\colon\beta_1\neq b_1$
(o las versiones unilaterales) se usa
\[
  t^* = \frac{\est{\beta}_1-b_1}{\sqrt{MSE/\SSX}},
\]
con las mismas reglas de decisión.

\begin{ejemplo}[Volumen de madera: contrastes sobre $\beta_1$]
  \begin{itemize}
    \item \emph{Bilateral}, $H_0\colon\beta_1=0$:
      $t^* = 6.916/0.7531 = 9.18$. Como $|t^*|>\tq{18}{0.975}=2.101$, se rechaza $H_0$: hay
      relación lineal entre el diámetro del tronco y el volumen de madera.
    \item \emph{Unilateral}, $H_0\colon\beta_1\le0$ frente a $H_1\colon\beta_1>0$: el estadístico
      es el mismo, $t^*=9.18>\tq{18}{0.95}=1.73$, y se rechaza $H_0$: $\beta_1$ es positivo, es
      decir, el volumen aumenta significativamente con el diámetro.
    \item \emph{Valor específico}, $H_0\colon\beta_1=6$ frente a $H_1\colon\beta_1\neq6$:
      $t^* = (6.916-6)/0.7531 = 1.22 \le 2.101$, y no se rechaza $H_0$: no hay evidencia de que
      el incremento medio de volumen por pulgada de diámetro sea distinto de 6 pies cúbicos.
  \end{itemize}
\end{ejemplo}

\subsection{Inferencias sobre \texorpdfstring{$\beta_0$}{beta0}}

\begin{teorema}[Distribución de $\est{\beta}_0$][teo:t3-dist-b0]
  El estimador $\est{\beta}_0=\media{Y}-\est{\beta}_1\media{X}$ verifica
  \[
    \est{\beta}_0 \sim \Normal\!\left(\beta_0,\ \sigma^2\left(\frac1n+\frac{\media{X}^2}{\SSX}\right)\right),
    \qquad
    \frac{\est{\beta}_0-\beta_0}{\raizMSE{\frac1n+\frac{\media{X}^2}{\SSX}}} \sim \tdist{n-2}.
  \]
\end{teorema}

\begin{demostracion}
  \emph{Normalidad.} Escribiendo $\media{Y}$ y $\est{\beta}_1$ como combinaciones lineales de
  las $Y_i$ (\cref{prop:t3-ki}),
  \begin{align*}
    \est{\beta}_0 &= \media{Y}-\est{\beta}_1\media{X}
      = \sumi \frac1n\,Y_i - \media{X}\sumi k_iY_i \\
    &= \sumi\left(\frac1n - \media{X}k_i\right)Y_i \\
    &= \sumi\left(\frac1n - \frac{\media{X}(X_i-\media{X})}{\SSX}\right)Y_i ,
  \end{align*}
  que es combinación lineal de normales independientes; luego es normal.

  \emph{Insesgadez.} Como $\E(\est{\beta}_1)=\beta_1$,
  \begin{align*}
    \E(\est{\beta}_0) &= \E(\media{Y}) - \media{X}\,\E(\est{\beta}_1) \\
    &= \frac1n\sumi\E(Y_i) - \beta_1\media{X} \\
    &= \frac1n\sumi(\beta_0+\beta_1X_i) - \beta_1\media{X} \\
    &= \beta_0 + \beta_1\media{X} - \beta_1\media{X} \\
    &= \beta_0 .
  \end{align*}

  \emph{Varianza.} Por la independencia de las $Y_i$ y usando las propiedades (i) y (iii) de
  la \cref{prop:t3-ki},
  \begin{align*}
    \Var(\est{\beta}_0) &= \sumi\left(\frac1n-\media{X}k_i\right)^2\Var(Y_i) \\
    &= \sigma^2\sumi\left(\frac{1}{n^2} - \frac{2\media{X}}{n}\,k_i + \media{X}^2k_i^2\right) \\
    &= \sigma^2\left(\frac{n}{n^2} - \frac{2\media{X}}{n}\sumi k_i
       + \media{X}^2\sumi k_i^2\right) \\
    &= \sigma^2\left(\frac1n - 0 + \frac{\media{X}^2}{\SSX}\right)
     = \sigma^2\left(\frac1n + \frac{\media{X}^2}{\SSX}\right).
  \end{align*}
  Sustituyendo $\sigma^2$ por $MSE$ se obtiene el estimador insesgado
  \[
    \widehat{\Var}(\est{\beta}_0)=MSE\left(\frac1n+\frac{\media{X}^2}{\SSX}\right)
  \]
  y, por el mismo argumento que en el \cref{teo:t3-dist-b1}, la distribución $\tdist{n-2}$ del
  estadístico tipificado.
\end{demostracion}

\begin{resumen}[Inferencias sobre $\beta_0$]
  \begin{itemize}
    \item IC de nivel $1-\alpha$:
      $\displaystyle \beta_0 \in \left(\est{\beta}_0 \pm \tq{n-2}{1-\alpha/2}
      \raizMSE{\frac1n+\frac{\media{X}^2}{\SSX}}\right)$.
    \item Contraste de $H_0\colon\beta_0=b_0$:
      $\displaystyle t^* = \frac{\est{\beta}_0-b_0}{\raizMSE{\frac1n+\frac{\media{X}^2}{\SSX}}}$.
      En el contraste bilateral se rechaza $H_0$ si $|t^*|>\tq{n-2}{1-\alpha/2}$; en el
      unilateral ($H_1\colon\beta_0>b_0$), si $t^*>\tq{n-2}{1-\alpha}$.
  \end{itemize}
\end{resumen}

\subsection{Consideraciones sobre las inferencias}

\begin{itemize}
  \item \textbf{Normalidad.} Aunque las $Y_i$ no sean normales, las distribuciones de
    $\est{\beta}_0$ y $\est{\beta}_1$ son aproximadamente normales, en las condiciones del
    teorema central del límite.
  \item Se asume que $X$ es constante, y las inferencias son válidas en el \textbf{rango de
    valores de $X$ observado} en la muestra.
  \item Las varianzas de $\est{\beta}_0$ y $\est{\beta}_1$ dependen de $\SSX=\sum(X_i-\media{X})^2$:
    \textbf{mayor dispersión en $X$ implica menor variabilidad de los estimadores}. Esto es útil
    al planificar la recogida de datos.
\end{itemize}

\subsection{Inferencia simultánea: el problema de las comparaciones múltiples}
\label{sec:t3-bonferroni}

Hasta ahora se ha hecho inferencia sobre cada coeficiente por separado, fijando un nivel de
significación (por ejemplo, $\alpha=0.05$) en cada contraste o intervalo. Ese nivel está
garantizado para cada coeficiente por separado, pero no para el conjunto de las inferencias:
aunque fueran independientes, la confianza conjunta de dos intervalos del 95\,\% sería
$0.95^2\approx0.90$; y como $\est{\beta}_0$ y $\est{\beta}_1$ no son independientes (salvo si
$\media{X}=0$, \cref{prop:t5-var-simple}), ni siquiera puede calcularse como un producto.
Aparece así el problema de las \textbf{comparaciones múltiples}:
\begin{itemize}
  \item la confianza conjunta de varios IC simultáneos es menor que la de cada uno (menor del
    95\,\%);
  \item el nivel de significación del contraste simultáneo es mayor que el de cada contraste
    (mayor que 0.05): aumentan los falsos positivos.
\end{itemize}

\begin{ejemplo}[Probabilidad de falso positivo]
  Al lanzar un par de dados, llamemos ``error de tipo I'' a obtener dos unos, con probabilidad
  $\alpha=\frac16\cdot\frac16\approx0.028$. En una sola tirada la probabilidad de ese resultado
  es $\alpha$,
  pero al repetir las tiradas es cada vez más probable que alguna dé dos unos.

  Del mismo modo, si se realizan 20 contrastes independientes, cada uno a nivel $\alpha=0.05$,
  \begin{align*}
    P(\text{al menos un contraste significativo por azar})
    &= 1 - P(\text{ninguno significativo por azar}) \\
    &= 1-(1-0.05)^{20} \approx 0.64 .
  \end{align*}
\end{ejemplo}

Una solución a este problema es el \textbf{ajuste de Bonferroni}: modificar el nivel de
significación o de confianza de cada inferencia individual para conseguir el nivel deseado en la
inferencia simultánea.

\begin{proposicion}[Desigualdad de Bonferroni][prop:bonferroni]
  Para sucesos $A_1,\ldots,A_m$,
  \(
    P\left(\bigcup_{i=1}^m A_i\right) \le \sum_{i=1}^m P(A_i).
  \)
\end{proposicion}

\paragraph{Caso de $\beta_0$ y $\beta_1$.} Sea $A_0$ ($A_1$) el suceso ``el IC individual de
nivel $1-\alpha$ para $\beta_0$ ($\beta_1$) no contiene a $\beta_0$ ($\beta_1$)'', de modo que
$P(A_0)=P(A_1)=\alpha$. La probabilidad de que \emph{ambos} intervalos contengan a su parámetro
es
\[
  P(\overline{A_0}\cap\overline{A_1}) = 1-P(A_0\cup A_1)
  = 1-P(A_0)-P(A_1)+P(A_0\cap A_1) \ \ge\ 1-2\alpha ,
\]
porque $P(A_0\cap A_1)\ge0$. Así, dos IC individuales del 95\,\% garantizan una confianza
simultánea de al menos el 90\,\%, y dos IC individuales del 97.5\,\% garantizan una confianza
simultánea de al menos el 95\,\%.

\paragraph{Caso general.} Se construyen $m$ intervalos individuales de nivel $1-\alpha^*$,
\[
  C_i\colon\quad \beta_i \in \left(\est{\beta}_i \pm \tq{n-2}{1-\alpha^*/2}
  \sqrt{\widehat{\Var}(\est{\beta}_i)}\right),
  \qquad P(\beta_i\in C_i)=1-\alpha^* .
\]
Por la desigualdad de Bonferroni, la confianza simultánea cumple
\[
  P\left(\bigcap_{i=1}^m C_i\right)
  = 1-P\left(\bigcup_{i=1}^m \overline{C_i}\right)
  \ \ge\ 1-\sum_{i=1}^m P(\overline{C_i}) = 1-m\alpha^* .
\]

\begin{resumen}[Ajuste de Bonferroni]
  Para garantizar un nivel de confianza simultáneo $1-\alpha$ en $m$ intervalos (o un nivel de
  significación global $\alpha$ en $m$ contrastes) basta tomar en cada uno
  \[
    \alpha^* = \frac{\alpha}{m}.
  \]
\end{resumen}

% ==============================================================================
%  3.2  Inferencias sobre la respuesta media
%  3.3  Inferencias sobre una nueva observación
% ==============================================================================
\section{Inferencias sobre la respuesta media}
\label{sec:t3-respuesta-media}

Un objetivo habitual de los modelos de regresión es estimar la respuesta media para un valor
dado $x_h$ de $X$. En el \cref{ej:arboles}, interesa estimar el volumen medio de madera de los
árboles con un determinado diámetro, para decidir si merece la pena talarlos.

Denotamos $\E(Y_h)=\E(Y\cond X=x_h)=\beta_0+\beta_1x_h$. Su \textbf{estimador puntual} es el
valor de la recta ajustada en $x_h$:
\[
  \est{Y}_h = \est{\beta}_0 + \est{\beta}_1x_h .
\]

\begin{teorema}[Distribución de $\est{Y}_h$][teo:t3-dist-Yh]
  \begin{gather*}
    \est{Y}_h \sim \Normal\!\left(\E(Y_h),\ \sigma^2\left(\frac1n+\frac{(x_h-\media{X})^2}{\SSX}\right)\right),
    \\
    \frac{\est{Y}_h-\E(Y_h)}{\raizMSE{\frac1n+\frac{(x_h-\media{X})^2}{\SSX}}} \sim \tdist{n-2}.
  \end{gather*}
\end{teorema}

\begin{demostracion}
  \emph{Normalidad.} $\est{Y}_h$ es combinación lineal de $\est{\beta}_0$ y $\est{\beta}_1$, y por
  tanto de las $Y_i$, que son normales independientes; luego es normal.

  \emph{Insesgadez.}
  \begin{align*}
    \E(\est{Y}_h) &= \E(\est{\beta}_0)+\E(\est{\beta}_1)\,x_h \\
    &= \beta_0+\beta_1x_h \\
    &= \E(Y_h) .
  \end{align*}

  \emph{Varianza.} Usando $\est{\beta}_0=\media{Y}-\est{\beta}_1\media{X}$,
  \begin{align*}
    \est{Y}_h &= \est{\beta}_0+\est{\beta}_1x_h \\
    &= \media{Y}-\est{\beta}_1\media{X}+\est{\beta}_1x_h \\
    &= \media{Y}+\est{\beta}_1(x_h-\media{X}) ,
  \end{align*}
  y, como $x_h$ y $\media{X}$ son constantes,
  \[
    \Var(\est{Y}_h) = \Var(\media{Y}) + (x_h-\media{X})^2\,\Var(\est{\beta}_1)
    + 2(x_h-\media{X})\Cov(\media{Y},\est{\beta}_1) .
  \]
  La covarianza es nula: escribiendo $\media{Y}$ y $\est{\beta}_1$ como combinaciones lineales de
  las $Y_i$, que son independientes con varianza $\sigma^2$, y usando
  la propiedad (i) de los $k_i$ (\cref{prop:t3-ki}),
  \begin{align*}
    \Cov(\media{Y},\est{\beta}_1) &= \Cov\left(\sumi \frac1n\,Y_i,\ \sumi k_iY_i\right) \\
    &= \sumi \frac1n\,k_i\,\Var(Y_i) \\
    &= \frac{\sigma^2}{n}\sumi k_i = 0 .
  \end{align*}
  Por tanto, con $\Var(\media{Y})=\sigma^2/n$ y $\Var(\est{\beta}_1)=\sigma^2/\SSX$
  (\cref{teo:t3-dist-b1}),
  \begin{align*}
    \Var(\est{Y}_h) &= \frac{\sigma^2}{n} + (x_h-\media{X})^2\,\frac{\sigma^2}{\SSX} \\
    &= \sigma^2\left(\frac1n+\frac{(x_h-\media{X})^2}{\SSX}\right).
  \end{align*}
  Al sustituir $\sigma^2$ por $MSE$ se obtiene la distribución $\tdist{n-2}$, por el mismo
  argumento que en el \cref{teo:t3-dist-b1} ($\est{Y}_h$ es función de $\est{\beta}_0$ y
  $\est{\beta}_1$, independientes de $SSE$).
\end{demostracion}

\begin{resumen}[IC de nivel $1-\alpha$ para la respuesta media $\E(Y_h)$]
  \[
    \E(Y_h) \in \left(\est{Y}_h \pm \tq{n-2}{1-\alpha/2}
    \raizMSE{\frac1n+\frac{(x_h-\media{X})^2}{\SSX}}\right).
  \]
\end{resumen}

La varianza de $\est{Y}_h$ es tanto menor cuanto más próximo está $x_h$ a la media muestral
$\media{X}$: los intervalos más precisos se obtienen en el punto $(\media{X},\media{Y})$ y se
ensanchan al alejarse de él.

% ------------------------------------------------------------------------------
\section{Inferencias sobre una nueva observación}
\label{sec:t3-prediccion}

Se quiere ahora predecir el valor de una \textbf{nueva observación} de $Y$, $\Ynew$, cuando
$X=x_h$. Por ejemplo, se mide el diámetro de un árbol concreto, 10.2 pulgadas, y se quiere
predecir el volumen de madera de \emph{ese} árbol.

\begin{itemize}
  \item La nueva observación corresponde a un nuevo individuo, al que se asume aplicable el
    modelo de regresión ajustado.
  \item Hay que distinguir entre la predicción de la respuesta media, $\E(Y_h)$, y la de una
    nueva observación, $\Ynew$:
    \begin{itemize}
      \item $\E(Y_h)$ es la media de la distribución de $Y$ condicionada a $X=x_h$;
      \item $\Ynew$ es un valor de una variable aleatoria con la distribución de $Y$
        condicionada a $X=x_h$, y difícilmente coincidirá con $\E(Y_h)$.
    \end{itemize}
\end{itemize}
El \textbf{estimador puntual} es el mismo que para la respuesta media,
\[
  \est{Y}_{h(\mathrm{new})} = \est{Y}_h = \est{\beta}_0 + \est{\beta}_1x_h ,
\]
pero el error que se comete es distinto, y por tanto también su distribución.

\paragraph{Parámetros conocidos.} Si $\beta_0$, $\beta_1$ y $\sigma^2$ fueran conocidos, la única
fuente de incertidumbre sería la variabilidad de $Y$ alrededor de su media:
$\Ynew\sim\Normal(\beta_0+\beta_1x_h,\sigma^2)$, y un intervalo de nivel $1-\alpha$ sería
$\E(Y_h)\pm z_{1-\alpha/2}\,\sigma$.

\begin{ejemplo}[Volumen de madera con parámetros conocidos]
  Supongamos $\beta_0=-45.57$, $\beta_1=6.916$ y $\sigma^2=36.587$ ($\sigma=6.049$).
  Para un árbol
  de diámetro $x_h=10.2$ pulgadas, el volumen medio es
  $\E(Y_h)=-45.57+6.916\cdot10.2=24.97$ pies cúbicos. Si el modelo es adecuado (errores normales,
  independientes y con varianza constante), el volumen de ese árbol sigue una
  $\Normal(24.97,\ 36.587)$, y un intervalo del 95\,\% para él es
  \[
    24.97 \pm 1.96\cdot6.049 = (13.11,\ 36.83).
  \]
\end{ejemplo}

\paragraph{Parámetros desconocidos.} En la práctica los parámetros se estiman, y al estimar
$\E(Y_h)$ se comete un error: el centro del intervalo para $\Ynew$ es a su vez incierto (no se
conoce con certeza la localización de la distribución de $Y$). Hay dos fuentes de
variabilidad:
\begin{enumerate}
  \item la asociada a la estimación de la respuesta media $\E(Y_h)$;
  \item la asociada a la distribución de $Y$ alrededor de su media.
\end{enumerate}

\begin{figure}[H]
  \centering
  \begin{tikzpicture}[x=1cm, y=2.2cm]
    \draw[-{Stealth}] (-4.8,0) -- (4.8,0);
    \foreach \c in {-2,2} {
      \draw[thick, domain=\c-2.6:\c+2.6, samples=60] plot (\x, {exp(-(\x-\c)^2/1.3)});
      \draw[dashed] (\c,0) -- (\c,1);
      \draw (\c-1.5,0) -- (\c-1.5,1.35);
      \draw (\c+1.5,0) -- (\c+1.5,1.35);
      \draw[{Stealth}-{Stealth}] (\c-1.5,1.25) -- (\c+1.5,1.25);
      \node[font=\scriptsize, fill=white, align=center] at (\c,1.25)
        {límites de predicción\\ si $\E(Y_h)$ está aquí};
    }
    \draw (0,0.05) -- (0,-0.05) node[below, font=\small] {$\est{Y}_h$};
    \draw[{Stealth}-{Stealth}] (-2,-0.45) -- (2,-0.45)
      node[midway, below, font=\scriptsize] {límites de confianza para $\E(Y_h)$};
  \end{tikzpicture}
  \caption{La predicción de una nueva observación combina la incertidumbre sobre
    $\E(Y_h)$ con la variabilidad de $Y$ alrededor de su media.}
\end{figure}

La cantidad relevante es el \textbf{error de predicción} $\Ynew-\est{Y}_h$.
Como $\Ynew$ es
independiente de las observaciones de la muestra, y por tanto de $\est{Y}_h$,
\begin{align*}
  \Var(\Ynew-\est{Y}_h) &= \Var(\Ynew) + \Var(\est{Y}_h)
  = \sigma^2 + \sigma^2\left(\frac1n+\frac{(x_h-\media{X})^2}{\SSX}\right) \\
  &= \sigma^2\left(1+\frac1n+\frac{(x_h-\media{X})^2}{\SSX}\right),
\end{align*}
y el error de predicción tiene media $\E(Y_h)-\E(Y_h)=0$. Es normal, por ser diferencia de
normales independientes, y es independiente de $SSE$, ya que $\Ynew$ es independiente de la
muestra y $\est{Y}_h$ lo es de $SSE$; por ello, al sustituir $\sigma^2$ por $MSE$ se obtiene de
nuevo una $\tdist{n-2}$.

\begin{teorema}[Intervalo de predicción][teo:t3-prediccion]
  \begin{gather*}
    \Ynew-\est{Y}_h \sim \Normal\!\left(0,\ \sigma^2\left(1+\frac1n+\frac{(x_h-\media{X})^2}{\SSX}\right)\right),
    \\
    \frac{\Ynew-\est{Y}_h}{\raizMSE{1+\frac1n+\frac{(x_h-\media{X})^2}{\SSX}}} \sim \tdist{n-2},
  \end{gather*}
  y el \textbf{intervalo de predicción} de nivel $1-\alpha$ para $\Ynew$ es
  \[
    \Ynew \in \left(\est{Y}_h \pm \tq{n-2}{1-\alpha/2}
    \raizMSE{1+\frac1n+\frac{(x_h-\media{X})^2}{\SSX}}\right).
  \]
\end{teorema}

El intervalo para una nueva observación tiene el mismo centro que el de la respuesta media, pero
mayor varianza. Por tanto, el IC para la respuesta media está siempre contenido en el intervalo de
predicción.

\subsection{Media de \texorpdfstring{$m$}{m} nuevas observaciones}

En ocasiones interesa predecir la media $\Ymnew$ de $m$ nuevas observaciones de $Y$ para
$X=x_h$. Se asume que el modelo es aplicable a los $m$ nuevos individuos, que son por tanto
independientes. Las fuentes de variabilidad son la estimación de $\E(Y_h)$ y la variabilidad de
la media de $m$ observaciones de $Y\cond X=x_h$, con $\Var(\Ymnew)=\sigma^2/m$. Así,
\begin{align*}
  \Var(\Ymnew-\est{Y}_h) &= \frac{\sigma^2}{m} + \sigma^2\left(\frac1n+\frac{(x_h-\media{X})^2}{\SSX}\right) \\
  &= \sigma^2\left(\frac1m+\frac1n+\frac{(x_h-\media{X})^2}{\SSX}\right).
\end{align*}

\begin{resumen}[Intervalo de predicción para la media de $m$ nuevas observaciones]
  \begin{gather*}
    \frac{\Ymnew-\est{Y}_h}{\raizMSE{\frac1m+\frac1n+\frac{(x_h-\media{X})^2}{\SSX}}}\sim\tdist{n-2},
    \\
    \Ymnew \in \left(\est{Y}_h \pm \tq{n-2}{1-\alpha/2}
    \raizMSE{\frac1m+\frac1n+\frac{(x_h-\media{X})^2}{\SSX}}\right).
  \end{gather*}
  Con $m=1$ se recupera el intervalo de predicción del \cref{teo:t3-prediccion}.
\end{resumen}

\subsection{Predicción inversa}

En ocasiones interesa \textbf{predecir el valor de $X$} a partir de un valor observado de $Y$,
utilizando el modelo de $Y$ sobre $X$.

\begin{ejemplo}[Calibración de un aparato de medida][ej:t3-calibracion]
  Se realiza un experimento con 12 individuos para evaluar un aparato que mide la concentración
  de azúcar (galactosa) en sangre. A partir de concentraciones conocidas ($X$) se registran las
  mediciones del aparato ($Y$) y se ajusta el modelo
  \[
    \est{Y} = -0.1 + 1.017\,X .
  \]
  Para un nuevo individuo solo se dispone de la medición del aparato, $\Ynew$, y se quiere
  estimar su verdadera concentración $\Xnew$.
\end{ejemplo}

Despejando en la recta ajustada se obtiene
$\est{X}=(Y+0.1)/1.017=0.098+0.983\,Y$. Esta recta \emph{no} es la recta de regresión de $X$
sobre $Y$: la de $X$ sobre $Y$ se obtendría minimizando las desviaciones en la dirección de $X$,
con pendiente $SS_{XY}/SS_Y$, mientras que al despejar se obtiene pendiente
$1/\est{\beta}_1=\SSX/SS_{XY}$; ambas coinciden solo si $r_{XY}^2=1$. Además, en este problema $X$
es la variable controlada (no aleatoria) e $Y$ la aleatoria, por lo que el modelo adecuado es el
de $Y$ sobre $X$. La estimación de $X$ se hace despejando en él, lo que se denomina
\textbf{predicción inversa}.

\begin{definicion}[Predicción inversa][def:t3-prediccion-inversa]
  Conocido el modelo estimado $\est{Y}=\est{\beta}_0+\est{\beta}_1X$ y observado $\Ynew$ en un
  nuevo individuo, el estimador puntual de $\Xnew$ es
  \[
    \est{X}_{h(\mathrm{new})} = \frac{\Ynew-\est{\beta}_0}{\est{\beta}_1}, \qquad \est{\beta}_1\neq0 .
  \]
\end{definicion}

Para hacer inferencia sobre $\Xnew$ se usan las siguientes propiedades (aproximadas, ya que
$\est{X}_{h(\mathrm{new})}$ es un cociente de variables aleatorias):
\begin{itemize}
  \item es aproximadamente normal;
  \item es aproximadamente insesgado:
    \[
      \E\big(\est{X}_{h(\mathrm{new})}\big)\approx
      \frac{\E(\Ynew)-\E(\est{\beta}_0)}{\E(\est{\beta}_1)}
      = \frac{\beta_0+\beta_1\Xnew-\beta_0}{\beta_1} = \Xnew ;
    \]
  \item su varianza es aproximadamente
    $\Var\big(\est{X}_{h(\mathrm{new})}\big)\approx
    \dfrac{\sigma^2}{\beta_1^2}\left(1+\dfrac1n+\dfrac{(\Xnew-\media{X})^2}{\SSX}\right)$.
\end{itemize}

\begin{resumen}[IC aproximado de nivel $1-\alpha$ para $\Xnew$]
  \[
    \Xnew \in \left(\est{X}_{h(\mathrm{new})} \pm \tq{n-2}{1-\alpha/2}
    \sqrt{\frac{MSE}{\est{\beta}_1^2}\left(1+\frac1n+
    \frac{\big(\est{X}_{h(\mathrm{new})}-\media{X}\big)^2}{\SSX}\right)}\,\right).
  \]
\end{resumen}

% ==============================================================================
%  3.4  Bandas de confianza para la recta de regresión
% ==============================================================================
\section{Bandas de confianza para la recta de regresión}
\label{sec:t3-bandas}

Se puede construir un IC para $\E(Y_h)$ en cada valor $x_h$. Repitiéndolo a lo largo de todo el
rango de $X$ observado se obtiene una \textbf{banda de confianza} para la recta de regresión. La
forma más sencilla es construir un IC de nivel $1-\alpha$ para cada valor de $X$: es la
\textbf{banda de confianza puntual} de nivel $1-\alpha$ para $\E(Y)$.

Como ocurre con cada IC para la respuesta media, la banda es más estrecha para valores de $X$
cercanos a $\media{X}$. Sin embargo, la banda puntual no garantiza el nivel $1-\alpha$ para
\emph{toda} la recta a la vez: por el problema de las comparaciones múltiples
(\cref{sec:t3-bonferroni}), su \textbf{nivel de confianza simultáneo es menor que $1-\alpha$}.

Interesa, por tanto, una \textbf{banda de confianza simultánea}: una región que, con confianza
$1-\alpha$, contenga a la recta de regresión completa. Esta región tiene además nivel al menos
$1-\alpha$ para cada valor de $X$ en particular.

El ajuste de Bonferroni solo es aplicable a un número finito $m$ de valores $x_h$ (tomando
$\alpha^*=\alpha/m$ en cada uno); para la recta completa, que involucra infinitos valores de $X$,
no sirve, y se utiliza la banda de Working--Hotelling.

\begin{definicion}[Banda de confianza de Working--Hotelling][def:working-hotelling]
  La banda de confianza de Working--Hotelling de nivel $1-\alpha$ para la recta de regresión
  está formada, para cada valor $x_h$ de $X$ (observado o no, entre el mínimo y el máximo
  observados), por los intervalos
  \[
    \est{Y}_h \pm W\raizMSE{\frac1n+\frac{(x_h-\media{X})^2}{\SSX}},
    \qquad
    W=\sqrt{2\,\Fq{2}{n-2}{1-\alpha}},
  \]
  donde $\Fq{2}{n-2}{1-\alpha}$ es el cuantil $1-\alpha$ de la distribución $F$ de
  Fisher--Snedecor con 2 y $n-2$ grados de libertad.
\end{definicion}

\begin{itemize}
  \item La expresión es la del IC para la respuesta media, sustituyendo el cuantil de la $t$ de
    Student por $W$, que es mayor y cubre el nivel simultáneo. Por ello la banda de
    Working--Hotelling es más ancha para todos los valores de $X$, y su nivel para cada $X$
    particular es $\ge1-\alpha$.
  \item Para obtener IC de nivel conjunto (simultáneos) para varios valores de $X$ debe usarse
    la banda de Working--Hotelling y no los IC individuales basados en la $t$.
  \item La banda es válida en el rango de $X$ en que el modelo lineal es adecuado, es decir,
    entre el mínimo y el máximo observados: el \emph{scope} del modelo.
  \item Los límites de la banda son una \textbf{hipérbola}: la región es más estrecha cerca de
    $\media{X}$ y más ancha al alejarse de él.
\end{itemize}

\begin{figure}[H]
  \centering
  \begin{tikzpicture}
    \begin{axis}[width=9cm, height=6.5cm, xlabel={DBH}, ylabel={VOL}, xmin=9.8, xmax=19.4,
        ymin=0, ymax=110, label style={font=\small}, tick label style={font=\footnotesize},
        legend style={font=\scriptsize, at={(0.98,0.02)}, anchor=south east, draw=none,
          fill=white, fill opacity=0.8, text opacity=1}]
      \addplot[puntos, mark size=1.2pt] table[x=dbh, y=vol] {datos/tema3/arboles.dat};
      \addlegendentry{datos}
      \addplot[black, thick, no marks] table[x=x, y=fit] {datos/tema3/bandas.dat};
      \addlegendentry{recta ajustada}
      \addplot[NavyBlue, dashed, no marks] table[x=x, y=cilo] {datos/tema3/bandas.dat};
      \addlegendentry{IC puntual ($t$)}
      \addplot[NavyBlue, dashed, no marks, forget plot] table[x=x, y=cihi] {datos/tema3/bandas.dat};
      \addplot[Maroon, densely dotted, thick, no marks] table[x=x, y=whlo] {datos/tema3/bandas.dat};
      \addlegendentry{Working--Hotelling}
      \addplot[Maroon, densely dotted, thick, no marks, forget plot] table[x=x, y=whhi] {datos/tema3/bandas.dat};
    \end{axis}
  \end{tikzpicture}
  \caption{Banda de confianza puntual del 95\,\% basada en la $t$ y banda simultánea de
    Working--Hotelling, más ancha, en el ejemplo de los árboles (datos reconstruidos).}
\end{figure}

\paragraph{Bandas de predicción.} De forma análoga se construyen bandas para los valores
predichos a partir de los intervalos de predicción (\cref{teo:t3-prediccion}). Son mucho más
anchas que las de la respuesta media. En ambos casos, las estimaciones son más precisas para
valores de $X$ cercanos a $\media{X}$.

\begin{figure}[H]
  \centering
  \begin{tikzpicture}
    \begin{axis}[width=9cm, height=6.5cm, xlabel={DBH}, ylabel={VOL}, xmin=9.8, xmax=19.4,
        ymin=0, ymax=110, label style={font=\small}, tick label style={font=\footnotesize},
        legend style={font=\scriptsize, at={(0.98,0.02)}, anchor=south east, draw=none,
          fill=white, fill opacity=0.8, text opacity=1}]
      \addplot[puntos, mark size=1.2pt, forget plot] table[x=dbh, y=vol] {datos/tema3/arboles.dat};
      \addplot[black, thick, no marks, forget plot] table[x=x, y=fit] {datos/tema3/bandas.dat};
      \addplot[NavyBlue, dashed, no marks] table[x=x, y=cilo] {datos/tema3/bandas.dat};
      \addlegendentry{IC para $\E(Y_h)$}
      \addplot[NavyBlue, dashed, no marks, forget plot] table[x=x, y=cihi] {datos/tema3/bandas.dat};
      \addplot[gray, no marks] table[x=x, y=pilo] {datos/tema3/bandas.dat};
      \addlegendentry{intervalo de predicción}
      \addplot[gray, no marks, forget plot] table[x=x, y=pihi] {datos/tema3/bandas.dat};
    \end{axis}
  \end{tikzpicture}
  \caption{Banda del 95\,\% para la respuesta media y banda de predicción del 95\,\% en el
    ejemplo de los árboles (datos reconstruidos).}
\end{figure}

% ==============================================================================
%  3.5  ANOVA en regresión simple. Descomposición de la variabilidad
%  3.6  El contraste de la regresión. F-test general
% ==============================================================================
\section{ANOVA en regresión simple}
\label{sec:t3-anova}

El análisis de regresión puede abordarse desde la perspectiva del \textbf{análisis de la
varianza (ANOVA)}, que se basa en la descomposición de la variabilidad de la respuesta $Y$. La
variabilidad se mide mediante las desviaciones de cada observación respecto de su media,
$Y_i-\media{Y}$, que se descomponen en
\[
  Y_i-\media{Y} = (Y_i-\est{Y}_i) + (\est{Y}_i-\media{Y}) .
\]

\begin{figure}[H]
  \centering
  \begin{tikzpicture}
    \begin{axis}[width=9.5cm, height=6.5cm, xmin=0, xmax=10, ymin=0, ymax=31,
        xlabel={$x$}, ylabel={$y$}, tick label style={font=\footnotesize}, clip=false]
      \addplot[Maroon, no marks, domain=0:10] {30-3*x};
      \addplot[black, dashed, no marks, domain=0:10] {10};
      \node[font=\scriptsize, anchor=west] at (axis cs:10.1,10) {$\media{Y}$};
      \node[font=\scriptsize, Maroon, anchor=south west] at (axis cs:5.6,14.5) {$\est{Y}=\est{\beta}_0+\est{\beta}_1x$};
      \addplot[only marks, mark=*, mark size=2pt, NavyBlue] coordinates {(3,28)};
      \draw[NavyBlue, thick] (axis cs:3,28) -- (axis cs:3,21);
      \draw[ForestGreen, thick] (axis cs:3,21) -- (axis cs:3,10);
      \draw[decorate, decoration={brace, amplitude=4pt}] (axis cs:3.15,28) -- (axis cs:3.15,21)
        node[midway, right=4pt, font=\scriptsize] {$Y_i-\est{Y}_i$};
      \draw[decorate, decoration={brace, amplitude=4pt}] (axis cs:3.15,21) -- (axis cs:3.15,10)
        node[midway, right=4pt, font=\scriptsize] {$\est{Y}_i-\media{Y}$};
      \draw[decorate, decoration={brace, amplitude=5pt, mirror}] (axis cs:2.85,28) -- (axis cs:2.85,10)
        node[midway, left=5pt, font=\scriptsize] {$Y_i-\media{Y}$};
    \end{axis}
  \end{tikzpicture}
  \caption{Descomposición de la desviación de una observación respecto de la media.}
\end{figure}

\begin{teorema}[Descomposición de la variabilidad][teo:t3-sst]
  \[
    \underbrace{\sumi (Y_i-\media{Y})^2}_{SST}
    = \underbrace{\sumi (\est{Y}_i-\media{Y})^2}_{SSR}
    + \underbrace{\sumi (Y_i-\est{Y}_i)^2}_{SSE},
  \]
  es decir, variabilidad total de $Y$ = variabilidad explicada por la regresión + variabilidad no
  explicada por la regresión.
\end{teorema}

\begin{demostracion}
  Elevando al cuadrado $Y_i-\media{Y}=(Y_i-\est{Y}_i)+(\est{Y}_i-\media{Y})$ y sumando,
  \[
    SST = SSE + SSR + 2\sumi e_i(\est{Y}_i-\media{Y}),
  \]
  con $e_i=Y_i-\est{Y}_i$. Por las propiedades 1 y 5 de la recta de regresión
  (\cref{prop:t2-propiedades-recta}), $\sum e_i=0$ y $\sum\est{Y}_ie_i=0$, luego
  \(
    \sumi e_i(\est{Y}_i-\media{Y}) = \sumi \est{Y}_ie_i - \media{Y}\sumi e_i = 0 .
  \)
\end{demostracion}

\subsection{Los tres componentes}

\paragraph{Variabilidad total, $SST=\sum(Y_i-\media{Y})^2$.} Sin tener en cuenta $X$, el mejor
predictor de $Y$ es $\media{Y}$, y $SST$ es la suma de las desviaciones cuadráticas respecto de
él. Tiene $n-1$ grados de libertad, ya que hay que estimar $\media{Y}$ (equivalentemente, las
desviaciones $Y_i-\media{Y}$ suman 0). La media cuadrática $MST=SST/(n-1)$ es el estimador
habitual de la varianza cuando no hay ningún predictor.

\paragraph{Variabilidad explicada por el modelo, $SSR=\sum(\est{Y}_i-\media{Y})^2$.}
\begin{itemize}
  \item $\est{Y}_i-\media{Y}$ son las desviaciones de la recta ajustada (que tiene en cuenta $X$)
    respecto del mejor predictor sin tener en cuenta $X$.
  \item Si no hay relación entre $X$ e $Y$ ($\est{\beta}_1=0$), entonces
    $\est{Y}_i=\est{\beta}_0=\media{Y}$ para todo $i$ y $SSR=0$. En otro caso, $SSR>0$.
  \item Tiene \textbf{1 grado de libertad}: aunque hay $n$ desviaciones, todos los $\est{Y}_i$ se
    calculan con la misma recta, con 2 grados de libertad (\emph{intercept} y pendiente), y se pierde
    uno porque las desviaciones $\est{Y}_i-\media{Y}$ suman 0. La media cuadrática es $MSR=SSR/1$.
\end{itemize}

\paragraph{Variabilidad no explicada o residual, $SSE=\sum(Y_i-\est{Y}_i)^2$.}
\begin{itemize}
  \item Al usar $X$ como predictor, la variabilidad queda representada por las desviaciones de
    cada observación a la recta ajustada, $Y_i-\est{Y}_i$.
  \item $SSE=0$ si todas las observaciones coinciden con sus valores ajustados, y es mayor
    cuanto mayor es la dispersión alrededor de la recta.
  \item Tiene $n-2$ grados de libertad, ya que hay que estimar los dos coeficientes para calcular
    los $\est{Y}_i$. La media cuadrática $MSE=SSE/(n-2)$ es la mejor estimación de la varianza de
    $Y$ una vez que se condiciona a la variable explicativa.
\end{itemize}
Los grados de libertad también se descomponen: $n-1 = 1 + (n-2)$.

\subsection{El contraste de la regresión}

Las medias cuadráticas son variables aleatorias, puesto que lo son las $Y_i$. Sus esperanzas son
\[
  \E(MSE) = \sigma^2, \qquad \E(MSR) = \sigma^2 + \beta_1^2\,\SSX .
\]
Por tanto:
\begin{itemize}
  \item $MSE$ es un estimador insesgado de $\sigma^2$;
  \item si $\beta_1=0$ (no hay relación lineal), $MSR$ también es un estimador insesgado de
    $\sigma^2$;
  \item si $\beta_1\neq0$, $\E(MSR)>\E(MSE)$, ya que $\beta_1^2\SSX>0$.
\end{itemize}
Esto sugiere usar el cociente $MSR/MSE$ para contrastar $\beta_1=0$: cuanto mayor sea, más
evidencia habrá en contra de $\beta_1=0$.

\paragraph{Distribución bajo $H_0\colon\beta_1=0$.}
\begin{itemize}
  \item Si las $Y_i$ son independientes y proceden de la misma normal con varianza $\sigma^2$,
    $SST$ es una suma de cuadrados con $n-1$ grados de libertad y $SST/\sigma^2\sim\chisq{n-1}$.
  \item $SST$ es la suma de $SSR$ y $SSE$, cuyos grados de libertad, 1 y $n-2$, suman $n-1$.
  \item Si $\beta_1=0$, todas las $Y_i$ tienen la misma media $\beta_0$ y la misma varianza
    $\sigma^2$, y $SSR/\sigma^2$ y $SSE/\sigma^2$ son variables independientes con distribuciones
    $\chisq{1}$ y $\chisq{n-2}$, respectivamente.
\end{itemize}
En consecuencia,
\[
  F = \frac{MSR}{MSE} = \frac{\dfrac{SSR}{\sigma^2}\Big/1}{\dfrac{SSE}{\sigma^2}\Big/(n-2)}
  = \frac{\chisq{1}/1}{\chisq{n-2}/(n-2)} \bajoHo \Fdist{1}{n-2} .
\]

\begin{resumen}[Contraste de la regresión (F-test)]
  Para $H_0\colon\beta_1=0$ frente a $H_1\colon\beta_1\neq0$, con $F^*$ el valor observado de
  $F=MSR/MSE$:
  \begin{itemize}
    \item si $F^*\le\Fq{1}{n-2}{1-\alpha}$, no se rechaza $H_0$ a nivel $\alpha$;
    \item si $F^*>\Fq{1}{n-2}{1-\alpha}$, se rechaza $H_0$ a nivel $\alpha$.
  \end{itemize}
  En la práctica, se rechaza $H_0$ si el $p$-valor es menor que $\alpha$.
\end{resumen}

\paragraph{Equivalencia con el contraste $t$.} Para un nivel $\alpha$ dado, el test $F$ es
equivalente al \textbf{test bilateral} basado en la $t$ de Student para $H_0\colon\beta_1=0$. En
efecto, como $\est{\beta}_0=\media{Y}-\est{\beta}_1\media{X}$,
\[
  \est{Y}_i-\media{Y} = \est{\beta}_0+\est{\beta}_1x_i-\media{Y} = \est{\beta}_1(x_i-\media{X})
  \quad\Longrightarrow\quad
  SSR = \est{\beta}_1^2\sumi (x_i-\media{X})^2 = \est{\beta}_1^2\,\SSX ,
\]
y, puesto que $\widehat{\Var}(\est{\beta}_1)=MSE/\SSX$,
\[
  F = \frac{MSR}{MSE} = \frac{\est{\beta}_1^2\,\SSX}{MSE}
  = \frac{\est{\beta}_1^2}{\widehat{\Var}(\est{\beta}_1)}
  = \left(\frac{\est{\beta}_1}{\sqrt{\widehat{\Var}(\est{\beta}_1)}}\right)^2 = (t^*)^2 ,
\]
el cuadrado del estadístico $t$. La equivalencia es solo con el contraste \emph{bilateral}.

\subsection{La tabla ANOVA}

Toda la información anterior se resume en la \textbf{tabla ANOVA}:
\begin{center}
  \renewcommand{\arraystretch}{1.3}
  \begin{tabular}{lccccc}
    \toprule
    Fuente & $SS$ & g.l. & $MS$ & $F$ & $p$-valor \\
    \midrule
    Modelo (regresión) & $SSR$ & $1$ & $MSR$ & $F^*=MSR/MSE$ & $p$ \\
    Error & $SSE$ & $n-2$ & $MSE$ & & \\
    Total & $SST$ & $n-1$ & & & \\
    \bottomrule
  \end{tabular}
\end{center}

\begin{ejemplo}[Volumen de madera: tabla ANOVA][ej:t3-anova-arboles]
  \begin{center}
    \renewcommand{\arraystretch}{1.2}
    \begin{tabular}{lrrrrr}
      \toprule
      Fuente & $SS$ & g.l. & $MS$ & $F$ & $p$-valor \\
      \midrule
      Modelo & 3085.789 & 1 & 3085.789 & 84.34 & $<0.0001$ \\
      Error & 658.570 & 18 & 36.587 & & \\
      Total & 3744.358 & 19 & & & \\
      \bottomrule
    \end{tabular}
  \end{center}
  La salida del programa (SAS) incluye además los siguientes estadísticos:
  \begin{center}
    \renewcommand{\arraystretch}{1.2}
    \begin{tabular}{lr@{\qquad}lr}
      \toprule
      Root MSE ($\sqrt{MSE}$) & 6.04874 & R-Square ($R^2$) & 0.8241 \\
      Dependent Mean ($\media{Y}$) & 61.85150 & Adj R-Sq ($R^2$ ajustado) & 0.8143 \\
      Coeff Var & 9.77945 & & \\
      \bottomrule
    \end{tabular}
  \end{center}
  El coeficiente de variación es $\text{CV}=100\sqrt{MSE}/\media{Y}=100\cdot6.049/61.85=9.78$,
  y el $R^2$ ajustado es $1-\dfrac{SSE/(n-2)}{SST/(n-1)}=1-\dfrac{36.587}{3744.358/19}=0.8143$.
  El coeficiente de determinación $R^2$ se estudia en la \cref{sec:t3-correlacion}.

  Se rechaza $H_0\colon\beta_1=0$, y $F^*=84.34=9.18^2=(t^*)^2$, en concordancia con el
  contraste $t$ del \cref{ej:arboles}.
\end{ejemplo}

% ------------------------------------------------------------------------------
\section{El F-test general}
\label{sec:t3-ftest-general}

El contraste basado en el ANOVA se llama \textbf{contraste de la regresión}. En regresión simple
es equivalente al contraste $t$ de $H_0\colon\beta_1=0$, pero es un caso particular de un
contraste general para evaluar modelos lineales, que compara un \textbf{modelo completo}, con las
variables $X$ que se consideran apropiadas, con un \textbf{modelo reducido} en el que se eliminan
algunas de ellas. Es el mismo esquema de modelo restringido frente a no restringido del Tema 1
(\cref{sec:t1-medias}).

\paragraph{Modelo completo.} En regresión simple interviene una sola $X$:
$Y_i=\beta_0+\beta_1X_i+\varepsilon_i$, $\varepsilon_i\iid\Normal(0,\sigma^2)$. Se ajusta a los datos
(mínimos cuadrados, máxima verosimilitud\ldots) y su suma de cuadrados del error,
\[
  SSE_c = \sumi (Y_i-\est{Y}_i)^2 = \sumi \big(Y_i-\est{\beta}_0-\est{\beta}_1X_i\big)^2 ,
\]
mide la variabilidad de $Y$ alrededor de la recta ajustada. Tiene $gl_c=n-2$ grados de libertad.

\paragraph{Modelo reducido o restringido.} Si $\beta_1=0$, el modelo es
$Y_i=\beta_0+\varepsilon_i$, con $\est{\beta}_0=\media{Y}$ y suma de cuadrados del error
\[
  SSE_r = \sumi \big(Y_i-\est{\beta}_0\big)^2 = \sumi (Y_i-\media{Y})^2 = SST ,
\]
con $gl_r=n-1$ grados de libertad.

\paragraph{Estadístico.} Es preferible el modelo que comete menor error, y siempre
$SSE_c\le SSE_r$. Se contrasta
\[
  H_0\colon \text{modelo reducido} \qquad\text{frente a}\qquad H_1\colon\text{modelo completo}:
\]
una diferencia $SSE_r-SSE_c$ suficientemente pequeña apoya el modelo reducido; si no es pequeña,
incluir el parámetro adicional reduce sustancialmente la variabilidad de las $Y_i$ alrededor de
la recta ajustada, y es preferible el modelo completo.

\begin{resumen}[F-test general]
  \[
    F = \frac{(SSE_r-SSE_c)/(gl_r-gl_c)}{SSE_c/gl_c} .
  \]
  En regresión simple,
  \[
    F = \frac{(SST-SSE)/\big((n-1)-(n-2)\big)}{SSE/(n-2)} = \frac{SSR/1}{SSE/(n-2)}
    = \frac{MSR}{MSE},
  \]
  que es el contraste de la regresión.
\end{resumen}

% ==============================================================================
%  3.7  Medidas de la relación lineal entre dos variables.
%       Coeficiente de determinación y correlación lineal
% ==============================================================================
\section{Medidas de la relación lineal: determinación y correlación}
\label{sec:t3-correlacion}

\subsection{El coeficiente de determinación}

En la descomposición de la variabilidad (\cref{teo:t3-sst}), $SST$ mide la incertidumbre en la
predicción de $Y$ sin usar la información de $X$, y $SSE$ la incertidumbre al usarla. La
reducción de la variabilidad, $SST-SSE=SSR$, en proporción a la variabilidad total, mide cuánto
se reduce la incertidumbre al usar $X$ para predecir $Y$.

\begin{definicion}[Coeficiente de determinación][def:R2]
  \[
    R^2 = \frac{SSR}{SST} = 1-\frac{SSE}{SST}.
  \]
\end{definicion}

Como $0\le SSE\le SST$, se tiene $0\le R^2\le1$:
\begin{itemize}
  \item $R^2=0$ cuando $SSR=0$: la recta ajustada es horizontal ($\est{\beta}_1=0$,
    $\est{Y}=\media{Y}$) y no hay relación lineal entre $X$ e $Y$;
  \item $R^2=1$ cuando $SSE=0$: todas las observaciones coinciden con sus valores ajustados, y
    $X$ predice perfectamente a $Y$.
\end{itemize}
En la práctica se obtienen valores intermedios. Cuanto mayor es $R^2$, mayor es la reducción de
la variabilidad total de $Y$ al usar $X$ para predecirla.

\paragraph{Propiedades de $R^2$.}
\begin{itemize}
  \item $0\le R^2\le1$ y es adimensional.
  \item Es el cuadrado del coeficiente de correlación de Pearson de la muestra
    $\{(X_i,Y_i)\}_{i=1}^n$, y también del de los pares $\{(\est{Y}_i,Y_i)\}_{i=1}^n$.
  \item Cuanto mayor es $R^2$, mayor es la capacidad del regresor $X$ para predecir la respuesta
    $Y$, menor es $SSE$ y más cerca están los puntos de la recta ajustada.
  \item Toma el mismo valor si se usa $X$ para predecir $Y$ o $Y$ para predecir $X$.
\end{itemize}
En el \cref{ej:t3-anova-arboles}, $R^2=3085.789/3744.358=0.8241$: el diámetro del tronco explica el
82.4\,\% de la variabilidad del volumen de madera.

\subsection{El coeficiente de correlación lineal}

Cuando las dos variables son aleatorias, la medida habitual de la relación lineal entre ellas es
el coeficiente de correlación de Pearson muestral (\cref{def:pearson}), que en función de las
sumas de cuadrados y de la pendiente estimada se escribe
\[
  r_{XY} = \frac{\sumi (x_i-\media{X})(y_i-\media{Y})}
                {\sqrt{\sumi (x_i-\media{X})^2\sumi (y_i-\media{Y})^2}}
  = \frac{SS_{XY}}{\sqrt{\SSX\,SS_Y}} = \est{\beta}_1\sqrt{\frac{\SSX}{SS_Y}} .
\]

\begin{proposicion}[Relación entre $R^2$ y $r_{XY}$][prop:R2-r2]
  En la regresión lineal simple, $R^2=r_{XY}^2$.
\end{proposicion}

\begin{demostracion}
  Usando $SSR=\est{\beta}_1^2\SSX$ (\cref{sec:t3-anova}) y $SS_Y=SST$,
  \[
    r_{XY}^2 = \est{\beta}_1^2\,\frac{\SSX}{SS_Y} = \frac{SSR}{SST} = 1-\frac{SSE}{SST} = R^2 .
  \]
\end{demostracion}

\paragraph{Propiedades de $r_{XY}$.}
\begin{itemize}
  \item $-1\le r_{XY}\le1$.
  \item $|r_{XY}|$ mide la fuerza de la asociación lineal: cuanto mayor es, más fuerte es la
    asociación lineal entre $X$ e $Y$.
  \item $r_{XY}=0$ indica que no hay asociación lineal; $r_{XY}=1$ ($-1$), asociación lineal
    perfecta positiva (negativa): los puntos están sobre una recta de pendiente positiva
    (negativa).
  \item El signo indica asociación positiva ($r_{XY}>0$) o negativa ($r_{XY}<0$).
  \item No depende de las unidades de medida de las variables.
  \item Los papeles de $X$ e $Y$ son simétricos en su cálculo.
  \item Es muy sensible a observaciones atípicas.
\end{itemize}

\begin{nota}[Regresión y correlación]
  Son conceptos relacionados que no deben confundirse. El modelo de regresión describe la
  relación lineal entre una variable independiente $X$, fija, y una respuesta $Y$, aleatoria,
  cuyo comportamiento se quiere explicar o estimar. La correlación describe la fuerza de la
  relación lineal entre dos variables, ambas aleatorias.
\end{nota}

\subsection{Inferencias sobre el coeficiente de correlación poblacional}

Las inferencias sobre el coeficiente de correlación poblacional $\rho_{XY}$ se basan en la
distribución en el muestreo de $r_{XY}$.
\begin{itemize}
  \item Para $H_0\colon\rho_{XY}=0$ frente a $H_1\colon\rho_{XY}\neq0$ se utiliza que
    \[
      r_{XY}\,\frac{\sqrt{n-2}}{\sqrt{1-r_{XY}^2}} \bajoHo \tdist{n-2} .
    \]
  \item Para $\rho_{XY}\neq0$ la distribución de $r_{XY}$ es compleja y se usan aproximaciones.
    La más habitual es la \textbf{transformación $Z$ de Fisher}: para $n$ grande,
    \[
      z' = \frac12\ln\left(\frac{1+r_{XY}}{1-r_{XY}}\right)
      \ \approx\ \Normal\!\left(\frac12\ln\left(\frac{1+\rho_{XY}}{1-\rho_{XY}}\right),\ \frac{1}{n-3}\right).
    \]
\end{itemize}

% ==============================================================================
%  3.8  Regresión por el origen
% ==============================================================================
\section{Regresión por el origen}
\label{sec:t3-origen}

A veces se impone la restricción de que, cuando $X=0$, el valor esperado de $Y$ sea cero. Se
estiman entonces \textbf{regresiones que pasan por el origen}, es decir, rectas sin \emph{intercept}:
\[
  Y_i = \beta_1'X_i + \varepsilon_i, \qquad \varepsilon_i\iid\Normal(0,\sigma^2),\quad i=1,\ldots,n .
\]
Hay que ser cautelosos al imponer esta restricción, sobre todo cuando en la muestra no se
observan valores de $X$ cercanos a 0.

\paragraph{Estimación por mínimos cuadrados.} Se busca $\est{\beta}_1'$ que minimice
$Q=\sumi (Y_i-\beta_1'X_i)^2$. Derivando,
\[
  \frac{\partial Q}{\partial\beta_1'} = -2\sumi X_i(Y_i-\beta_1'X_i) = 0
  \quad\Longrightarrow\quad
  \sumi X_iY_i = \est{\beta}_1'\sumi X_i^2 ,
\]
de donde
\(
  \est{\beta}_1' = \dfrac{\sumi X_iY_i}{\sumi X_i^2}.
\)

\paragraph{Distribución en el muestreo.} $\est{\beta}_1'=\sum c_iY_i$ con
$c_i=X_i/\sum X_j^2$, combinación lineal de normales independientes, con
$\E(\est{\beta}_1')=\sum c_i\beta_1'X_i=\beta_1'$ y $\Var(\est{\beta}_1')=\sigma^2\sum c_i^2$. Así,
\[
  \est{\beta}_1' \sim \Normal\!\left(\beta_1',\ \frac{\sigma^2}{\sumi X_i^2}\right).
\]
Un estimador insesgado de $\sigma^2$ es
\[
  MSE = \frac{SSE}{n-1} = \frac{\sumi e_i^2}{n-1} = \frac{\sumi \big(Y_i-\est{\beta}_1'X_i\big)^2}{n-1},
\]
donde $SSE$ tiene $n-1$ grados de libertad (y no $n-2$), ya que solo se estima un parámetro,
$\beta_1'$.

\begin{resumen}[Inferencias en la regresión por el origen]
  \[
    \frac{\est{\beta}_1'-\beta_1'}{\sqrt{MSE\big/\sum_{i=1}^n X_i^2}} \sim \tdist{n-1};
    \qquad\text{para } H_0\colon\beta_1'=0:\quad
    \frac{\est{\beta}_1'}{\sqrt{MSE\big/\sum_{i=1}^n X_i^2}} \bajoHo \tdist{n-1}.
  \]
\end{resumen}

\paragraph{Ejemplos de uso.} Hay situaciones en las que tiene sentido añadir la restricción
$\beta_0=0$: la relación entre el número de artículos producidos en una hora y el número de
defectuosos en ese periodo, entre dimensiones de un producto, entre calorías ingeridas y
consumidas, o entre la cantidad de dinero invertida y su retorno.

\paragraph{Inconvenientes.} Es una restricción muy fuerte:
\begin{itemize}
  \item no se obtiene la ``mejor'' recta que se ajusta a los datos, sino la mejor de las que
    pasan por el punto $(0,0)$;
  \item en el modelo con \emph{intercept} se cumple $\sum e_i=\sum(Y_i-\est{Y}_i)=0$, pero al forzar la
    recta por el origen esto en general deja de ser cierto; la descomposición
    $SST=SSR+SSE$ ya no se verifica y puede ocurrir que $SSE>SST$, con lo que
    $R^2=1-SSE/SST<0$. Por ello $R^2$ deja de ser útil como medida de la bondad del ajuste.
\end{itemize}


% ==============================================================================
%  TEMA 4 — Validación de los modelos de regresión lineal simple
%  Fuente: diapositivas T4_Validacion_regresion.pdf
% ==============================================================================
\chapter{Validación del modelo de regresión lineal simple}
\label{cap:tema4}

% ==============================================================================
%  4.1  Objetivos de la validación
% ==============================================================================
\section{Objetivos de la validación}
\label{sec:t4-objetivos}

Validar el modelo consiste en comprobar que se cumplen los supuestos necesarios para que sea
apropiado. En el modelo de regresión lineal simple (\cref{def:rls}) hay que comprobar que
\[
  \varepsilon_i \iid \Normal(0,\sigma^2), \qquad i=1,\ldots,n,
\]
o, de forma equivalente (\cref{prop:t2-dist-Y}), que
\[
  Y_i \sim \Normal(\beta_0+\beta_1X_i,\ \sigma^2) \ \text{independientes}, \qquad i=1,\ldots,n .
\]
Esto supone verificar las cuatro hipótesis del modelo (\cref{sec:t2-hipotesis}):
\begin{itemize}
  \item \textbf{Linealidad}: si la relación es lineal, es decir, si $\E(Y)=\beta_0+\beta_1X$
    para todo $X$ o, equivalentemente, $\E(\varepsilon)=0$.
  \item \textbf{Homogeneidad de varianzas}: si la varianza depende de $X$; el modelo supone
    $\Var(Y)=\sigma^2$ o, equivalentemente, $\Var(\varepsilon)=\sigma^2$, para todo $X$.
  \item \textbf{Normalidad}: si $Y$ es normal o, equivalentemente, si lo son los errores
    $\varepsilon$.
  \item \textbf{Independencia}: si las $Y_i$ o, equivalentemente, los $\varepsilon_i$,
    $i=1,\ldots,n$, son independientes.
\end{itemize}

Además, la validación debe responder a otras cuestiones:
\begin{itemize}
  \item si el \textbf{predictor} $X$ incluido en el modelo es \textbf{adecuado}
    (\cref{sec:t4-predictor});
  \item si hay algún \textbf{predictor importante} que \textbf{no se haya incluido} en el
    modelo (\cref{sec:t4-omitidas});
  \item si hay \textbf{outliers y puntos de influencia}. En ese caso, la estimación del modelo
    puede depender en gran medida de un solo dato o de un pequeño subconjunto de
    observaciones, por lo que parte de la validación consiste en identificarlos
    (\cref{sec:t4-outliers}).
\end{itemize}

La validación se basa principalmente en los residuos, $e_i=y_i-\est{y}_i$, y por tanto se lleva a
cabo después del ajuste.

% ==============================================================================
%  4.2  Diagnóstico para el predictor X
% ==============================================================================
\section{Diagnóstico para el predictor}
\label{sec:t4-predictor}

El modelo no hace ninguna suposición específica sobre $X$, cuyos valores se consideran
conocidos (\cref{def:rls}). Por tanto, evaluar la idoneidad de $X$ como predictor no forma
parte de la validación del modelo propiamente dicha. Sin embargo, es necesario disponer de
información sobre $X$ para interpretar lo que ocurre con la respuesta. En concreto, interesa:
\begin{itemize}
  \item comprobar si hay \textbf{valores atípicos de $X$} que puedan influir en la recta de
    regresión ajustada (\cref{sec:t4-outliers-x});
  \item conocer el \textbf{rango y la dispersión} de los valores de $X$, para determinar el
    rango de validez de las predicciones realizadas con el modelo ajustado (el \emph{scope} del
    modelo, \cref{sec:t2-analisis}).
\end{itemize}
Para ello resulta útil un análisis descriptivo de $X$, con representaciones como diagramas de
cajas o histogramas.

% ==============================================================================
%  4.3  Los residuos: leverage, propiedades y tipos de residuos
% ==============================================================================
\section{Los residuos}
\label{sec:t4-residuos}

No es habitual verificar las hipótesis del modelo directamente sobre la respuesta $Y$; en su
lugar se utilizan los residuos (\cref{def:t2-ajuste}),
\[
  e_i = Y_i-\est{Y}_i ,
\]
diferencia entre el valor observado y el valor predicho por el modelo. Los residuos pueden
entenderse como el \textbf{error observado} y se utilizan para estimar el verdadero error
$\varepsilon_i=Y_i-\E(Y_i)$, que es desconocido. Cuando es necesario distinguirlos de sus
versiones tipificadas, que se definen más adelante, se denominan \textbf{residuos regulares}.

Como el modelo supone $\varepsilon_i\iid\Normal(0,\sigma^2)$, la idea básica del análisis de los
residuos es que, \textbf{si el modelo es apropiado para los datos, los residuos observados
deben reflejar las propiedades supuestas para los $\varepsilon_i$}.

\subsection{Los valores ajustados como combinación lineal de las observaciones: el leverage}

\begin{proposicion}[Valores ajustados como combinación lineal][prop:t4-hik]
  Cada valor ajustado es una combinación lineal de las observaciones:
  \[
    \est{Y}_i = \sumk h_{ik}Y_k,
    \qquad
    h_{ik} = \frac1n + \frac{(x_i-\media{X})(x_k-\media{X})}{\SSX}.
  \]
\end{proposicion}

\begin{demostracion}
  Usando $\est{\beta}_0=\media{Y}-\est{\beta}_1\media{X}$ y la expresión de $\est{\beta}_1$
  (\cref{teo:t2-estimadores-mc}),
  \begin{align*}
    \est{Y}_i &= \est{\beta}_0 + \est{\beta}_1x_i \\
    &= \media{Y} - \est{\beta}_1\media{X} + \est{\beta}_1x_i \\
    &= \media{Y} + \est{\beta}_1(x_i-\media{X}) \\
    &= \media{Y} + \frac{\sumk (x_k-\media{X})(Y_k-\media{Y})}{\SSX}\,(x_i-\media{X}) \\
    &= \media{Y} + \frac{\sumk Y_k(x_k-\media{X})(x_i-\media{X})}{\SSX}
      - \media{Y}\,\frac{\sumk (x_k-\media{X})(x_i-\media{X})}{\SSX} .
  \end{align*}
  El último término es nulo, ya que
  \begin{align*}
    \sumk (x_k-\media{X})(x_i-\media{X}) &= (x_i-\media{X})\sumk (x_k-\media{X}) \\
    &= (x_i-\media{X})\left(\sumk x_k - n\media{X}\right) \\
    &= (x_i-\media{X})\cdot0 = 0 .
  \end{align*}
  Por tanto, escribiendo también $\media{Y}$ como suma,
  \begin{align*}
    \est{Y}_i &= \media{Y} + \frac{\sumk (x_k-\media{X})(x_i-\media{X})Y_k}{\SSX} \\
    &= \sumk \frac1n\,Y_k + \sumk \frac{(x_i-\media{X})(x_k-\media{X})}{\SSX}\,Y_k \\
    &= \sumk \left(\frac1n + \frac{(x_i-\media{X})(x_k-\media{X})}{\SSX}\right)Y_k \\
    &= \sumk h_{ik}Y_k .
  \end{align*}
\end{demostracion}

\begin{definicion}[Leverage][def:t4-leverage]
  El \textbf{leverage} de la observación $i$-ésima es el coeficiente de $Y_i$ en su propio valor
  ajustado:
  \[
    h_{ii} = \frac1n + \frac{(x_i-\media{X})^2}{\SSX}.
  \]
  En castellano suele traducirse por \textbf{palanca} o \textbf{influencia}.
\end{definicion}

El leverage mide lo lejos que está el valor $x_i$ de la media muestral de las $X$ y, en cierto
modo, cuánto aporta la observación $i$-ésima a la varianza muestral de las $X$,
$\SSX/(n-1)$.

\begin{proposicion}[Propiedades del leverage][prop:t4-prop-leverage]
  \begin{enumerate}
    \item $h_{ii}$ no depende del valor $Y_i$ observado, sino solo de los valores de $X$.
    \item $\sum_{i=1}^n h_{ik} = \sumk h_{ik} = 1$.
    \item $\sum_{i=1}^n h_{ii} = 2$, el número de coeficientes del modelo.
    \item $\frac1n \le h_{ii} \le \frac1s \le 1$, donde $s$ es el número de observaciones de la
      muestra cuyo valor de $X$ es igual a $x_i$.
  \end{enumerate}
\end{proposicion}

\subsection{Propiedades de los residuos}

\begin{proposicion}[Propiedades de los residuos][prop:t4-prop-residuos]
  \begin{itemize}
    \item \emph{Esperanza}: $\E(e_i)=\E(Y_i-\est{Y}_i)=0$.
    \item \emph{Varianza}:
      \[
        \Var(e_i) = \sigma^2\left(1-\frac1n-\frac{(X_i-\media{X})^2}{\SSX}\right)
        = \sigma^2(1-h_{ii}).
      \]
    \item \emph{Covarianza}: para $i\neq j$,
      \[
        \Cov(e_i,e_j) = -\sigma^2\left(\frac1n+\frac{(X_i-\media{X})(X_j-\media{X})}{\SSX}\right)
        = -\sigma^2h_{ij}.
      \]
  \end{itemize}
\end{proposicion}

La varianza del residuo de la observación $i$-ésima depende de $h_{ii}$, y por tanto del valor
$x_i$, de modo que \textbf{no es homogénea}. Además, los residuos \textbf{no son variables
aleatorias independientes}, puesto que todos se calculan a partir de la misma recta de
regresión. En resumen, los residuos no cumplen ni la homogeneidad de varianzas ni la
independencia. Aun así, los $e_i$ se utilizan para validar el modelo, ya que el efecto de la
dependencia entre residuos y de la falta de homogeneidad de sus varianzas es de poca
importancia cuando:
\begin{itemize}
  \item el \textbf{tamaño muestral es grande} en relación con el número de parámetros del modelo
    de regresión, y
  \item \textbf{no hay puntos} que ejerzan una gran \textbf{influencia}
    (\cref{sec:t4-outliers}).
\end{itemize}

\subsection{Residuos semiestudentizados}

Los residuos $e_i$ se miden en las unidades de $Y$ y no tienen una escala conocida, por lo que
es difícil establecer qué es un valor ``extremo''. Por ello interesa tipificarlos: los residuos
tipificados no dependen de las unidades de medida, pueden compararse entre sí y con los valores
de referencia de una distribución conocida, y permiten identificar los residuos demasiado
grandes (o pequeños), que corresponden a observaciones mal ajustadas por el modelo
(\cref{sec:t4-outliers-y}).

Como $\varepsilon_i\iid\Normal(0,\sigma^2)$, los \textbf{errores estandarizados}
$(\varepsilon_i-0)/\sigma$ siguen una $\Normal(0,1)$. Puesto que $\sigma$ es desconocida, se estima
mediante $\sqrt{MSE}$.

\begin{definicion}[Residuos semiestudentizados][def:t4-semiestudentizados]
  Los \textbf{residuos estandarizados} o \textbf{semiestudentizados} son
  \[
    e_i^* = \frac{e_i-\media{e}}{\sqrt{MSE}} = \frac{e_i}{\sqrt{MSE}}, \qquad i=1,\ldots,n,
  \]
  donde se ha usado que $\media{e}=0$, porque $\sumi e_i=0$ (\cref{prop:t2-propiedades-recta}).
  Tienen media 0 y varianza aproximadamente 1.
\end{definicion}

Si $\sqrt{MSE}$ fuera una estimación de la desviación típica de los residuos, los $e_i^*$ estarían
estudentizados (seguirían una distribución $t$ de Student). Sin embargo, la varianza estimada de
cada residuo es (\cref{prop:t4-prop-residuos})
\[
  \widehat{\Var}(e_i) = MSE\left(1-\frac1n-\frac{(X_i-\media{X})^2}{\SSX}\right) = MSE(1-h_{ii}),
\]
por lo que la tipificación con $\sqrt{MSE}$ es solo una aproximación.

\subsection{Residuos estudentizados}

\begin{definicion}[Residuos estudentizados][def:t4-estudentizados]
  Los \textbf{residuos estudentizados} son
  \[
    r_i = \frac{e_i}{\sqrt{\est{\sigma}^2\left(1-\frac1n-\frac{(X_i-\media{X})^2}{\SSX}\right)}}
    = \frac{e_i}{\sqrt{MSE(1-h_{ii})}}, \qquad i=1,\ldots,n,
  \]
  donde $\est{\sigma}^2=MSE=SSE/(n-2)$ y $h_{ii}$ es el leverage de la observación $i$-ésima.
  Verifican $\E(r_i)=0$ y $\Var(r_i)=1$.
\end{definicion}

Bajo la hipótesis de que la observación $i$-ésima sigue el modelo ajustado, y suponiendo errores
normales, $r_i$ sigue aproximadamente una distribución $t$ de Student con $n-2$ grados de
libertad:
\[
  r_i \approx \tdist{n-2}.
\]

\begin{nota}
  La distribución $\tdist{n-2}$ de $r_i$ no es exacta, porque el numerador $e_i$ y el
  denominador $\sqrt{MSE}$ no son independientes: $e_i^2$ forma parte de $SSE$. De hecho, como
  $e_i^2/(1-h_{ii})\le SSE$, se tiene $|r_i|\le\sqrt{n-2}$. La distribución $t$ exacta la tiene el
  residuo estudentizado eliminado, $t_{(i)}\sim\tdist{n-3}$ (\cref{def:t4-eliminado-estudentizado}).
\end{nota}

Conocer esta distribución es muy útil, porque proporciona una referencia para decidir qué
residuos son demasiado grandes (o pequeños) y corresponden, por tanto, a observaciones que no
están bien ajustadas por el modelo: aproximadamente el 95\,\% de los $r_i$ deben estar entre
$-2$ y $2$, y casi todos (en torno al 99\,\%) entre $-3$ y $3$.

En muestras grandes los leverages son pequeños (su media es $2/n$) y $e_i^*\approx r_i$.

\subsection{Residuos eliminados}
\label{sec:t4-eliminados}

Para cada observación puede obtenerse su residuo en el modelo ajustado \emph{sin} considerar
dicha observación:
\begin{enumerate}
  \item Se elimina la observación $(x_i,y_i)$ de la muestra, obteniendo una muestra reducida con
    $n-1$ observaciones.
  \item Con la muestra reducida se estima la recta de regresión, $\est{\beta}_{0(i)}$ y
    $\est{\beta}_{1(i)}$, y la varianza, $\est{\sigma}^2_{(i)}=MSE_{(i)}$. El subíndice $(i)$ indica
    que los parámetros se han estimado sin usar la $i$-ésima observación.
  \item Para la observación omitida se calcula el valor ajustado
    $\est{Y}_{i(i)}=\est{\beta}_{0(i)}+\est{\beta}_{1(i)}x_i$. Como la observación $i$-ésima no se
    ha usado en la estimación de los parámetros, $y_i$ e $\est{Y}_{i(i)}$ son independientes.
  \item El \textbf{residuo eliminado} es
    \[
      e_{(i)} = y_i - \est{Y}_{i(i)} .
    \]
\end{enumerate}

\begin{definicion}[Residuos estudentizados eliminados][def:t4-eliminado-estudentizado]
  La varianza estimada del residuo eliminado $e_{(i)}$ es
  \[
    \widehat{\Var}\big(e_{(i)}\big) = \frac{MSE_{(i)}}{1-h_{ii}},
  \]
  y el \textbf{residuo estudentizado eliminado} (\emph{studentized deleted residual}) es
  \[
    t_{(i)} = \frac{e_{(i)}}{\sqrt{\dfrac{MSE_{(i)}}{1-h_{ii}}}}, \qquad i=1,\ldots,n .
  \]
  Bajo la hipótesis de que la $i$-ésima observación sigue el modelo ajustado, y suponiendo
  errores normales,
  \[
    t_{(i)} \sim \tdist{n-3} .
  \]
\end{definicion}

Los grados de libertad son $n-3$ porque el modelo sin la observación $i$-ésima se ajusta con una
observación menos: $(n-1)-2=n-3$.

Para calcular cada $t_{(i)}$, $i=1,\ldots,n$, habría que ajustar de nuevo el modelo sin la
observación $i$-ésima, de modo que para evaluar todas las observaciones serían necesarios $n$
ajustes. Existe una fórmula computacionalmente más sencilla, que no requiere reajustar el
modelo:
\[
  t_{(i)} = e_i\sqrt{\frac{n-3}{SSE(1-h_{ii})-e_i^2}} .
\]

% ==============================================================================
%  4.4  El análisis de los residuos: gráficos y contrastes
% ==============================================================================
\section{El análisis de los residuos}
\label{sec:t4-analisis}

Las \textbf{representaciones gráficas} de los residuos son la principal herramienta de
validación. Las más utilizadas son:
\begin{itemize}
  \item residuos frente a la variable $X$;
  \item valor absoluto o cuadrado de los residuos frente a la variable $X$;
  \item residuos frente a los valores predichos;
  \item residuos frente al orden de las observaciones;
  \item residuos frente a otras posibles variables $X$;
  \item diagrama de cajas de los residuos;
  \item gráfico de normalidad de los residuos.
\end{itemize}
A continuación se indica qué gráficos permiten estudiar cada hipótesis del modelo.

\subsection{Linealidad}

La hipótesis de linealidad se estudia en el \textbf{gráfico de residuos}: residuos frente a $X$
o, equivalentemente, frente a los valores predichos. En el modelo de regresión lineal simple
ambos gráficos son totalmente equivalentes, ya que $\est{Y}_i=\est{\beta}_0+\est{\beta}_1X_i$ es una
transformación lineal de $X_i$ (con el eje horizontal invertido si $\est{\beta}_1<0$), por lo
que lo habitual es utilizar el de los valores predichos.

Cuando el modelo lineal es adecuado, los residuos se sitúan centrados en 0 de forma aleatoria,
sin que se observen patrones sistemáticos (\emph{gráfico nulo}). Si la función de regresión no es
lineal, los residuos presentan un patrón curvo.

\begin{figure}[H]
  \centering
  \begin{subfigure}[t]{0.32\linewidth}
    \centering
    \begin{tikzpicture}
      \begin{axis}[residuos, xlabel={$\est{y}_i$ o $x_i$}]
        \addplot[cero] {0};
        \addplot[puntos] table[x=x, y=r] {datos/tema4/nulo.dat};
      \end{axis}
    \end{tikzpicture}
    \caption{Gráfico nulo: modelo adecuado.}
  \end{subfigure}\hfill
  \begin{subfigure}[t]{0.32\linewidth}
    \centering
    \begin{tikzpicture}
      \begin{axis}[residuos]
        \addplot[cero] {0};
        \addplot[puntos] table[x=x, y=r] {datos/tema4/curva1.dat};
      \end{axis}
    \end{tikzpicture}
    \caption{Patrón curvo: relación no lineal.}
  \end{subfigure}\hfill
  \begin{subfigure}[t]{0.32\linewidth}
    \centering
    \begin{tikzpicture}
      \begin{axis}[residuos]
        \addplot[cero] {0};
        \addplot[puntos] table[x=x, y=r] {datos/tema4/curva2.dat};
      \end{axis}
    \end{tikzpicture}
    \caption{Patrón curvo: relación no lineal.}
  \end{subfigure}
  \caption{Gráficos de residuos y linealidad (datos simulados).}
\end{figure}

\subsection{Homogeneidad de varianzas}

La hipótesis de homogeneidad de varianzas también se estudia en el gráfico de residuos. Si la
varianza es homogénea, la dispersión de los residuos es aproximadamente constante a lo largo de
todo el eje horizontal; en caso contrario, la nube de puntos se abre o se cierra.

\begin{figure}[H]
  \centering
  \begin{subfigure}[t]{0.45\linewidth}
    \centering
    \begin{tikzpicture}
      \begin{axis}[residuos, xlabel={$\est{y}_i$ o $x_i$}, ymin=-3, ymax=3]
        \addplot[cero] {0};
        \addplot[puntos] table[x=x, y=r] {datos/tema4/nulo.dat};
      \end{axis}
    \end{tikzpicture}
    \caption{Varianza constante.}
  \end{subfigure}\hfill
  \begin{subfigure}[t]{0.45\linewidth}
    \centering
    \begin{tikzpicture}
      \begin{axis}[residuos, xlabel={$\est{y}_i$ o $x_i$}]
        \addplot[cero] {0};
        \addplot[puntos] table[x=x, y=r] {datos/tema4/embudo_crec.dat};
      \end{axis}
    \end{tikzpicture}
    \caption{La varianza aumenta con $\est{y}_i$.}
  \end{subfigure}

  \medskip
  \begin{subfigure}[t]{0.45\linewidth}
    \centering
    \begin{tikzpicture}
      \begin{axis}[residuos, xlabel={$\est{y}_i$ o $x_i$}]
        \addplot[cero] {0};
        \addplot[puntos] table[x=x, y=r] {datos/tema4/embudo_decr.dat};
      \end{axis}
    \end{tikzpicture}
    \caption{La varianza disminuye con $\est{y}_i$.}
  \end{subfigure}\hfill
  \begin{subfigure}[t]{0.45\linewidth}
    \centering
    \begin{tikzpicture}
      \begin{axis}[residuos, xlabel={$\est{y}_i$ o $x_i$}]
        \addplot[cero] {0};
        \addplot[puntos] table[x=x, y=r] {datos/tema4/diamante.dat};
      \end{axis}
    \end{tikzpicture}
    \caption{Mayor varianza en los valores centrales.}
  \end{subfigure}
  \caption{Gráficos de residuos y homogeneidad de varianzas (datos simulados).}
\end{figure}

También son útiles los gráficos del valor absoluto o del cuadrado de los residuos frente a
$\est{Y}_i$ (o $X_i$), especialmente cuando no hay muchas observaciones: si la varianza no es
homogénea, la magnitud de los residuos, independientemente de su signo, cambia con el valor de
$\est{Y}_i$ (o $X_i$).

\begin{nota}[Residuos regulares frente a estudentizados]
  Como $\Var(e_i)=\sigma^2(1-h_{ii})$ (\cref{prop:t4-prop-residuos}) y $h_{ii}$ es mínimo cuando
  $X_i=\media{X}$, la varianza de los residuos regulares es \emph{mayor} cuanto más cerca está
  $X_i$ de su media muestral.%
   Este patrón podría verse en el gráfico de residuos aunque los errores
  tengan varianza constante. Por ese motivo es mucho más frecuente utilizar los residuos
  estudentizados, que tienen varianza 1 para todas las observaciones.
\end{nota}

\subsection{Normalidad}

Para estudiar la normalidad de los errores se utilizan:
\begin{itemize}
  \item el \textbf{gráfico cuantil-cuantil} (\emph{QQ-plot}), en el que los residuos ordenados se
    comparan con sus valores esperados bajo normalidad, es decir, con los cuantiles de la normal.
    Si los residuos son normales, los puntos se sitúan aproximadamente sobre una recta;
  \item el \textbf{histograma} o el \textbf{diagrama de cajas} de los residuos.
\end{itemize}

\begin{figure}[H]
  \centering
  \begin{subfigure}[t]{0.32\linewidth}
    \centering
    \begin{tikzpicture}
      \begin{axis}[mini, xlabel={$z$}, ylabel={$r_i$}, xmin=-3, xmax=3, ymin=-4, ymax=6]
        \addplot[cero, domain=-3:3] {x};
        \addplot[puntos] table[x=q, y=der] {datos/tema4/qq.dat};
      \end{axis}
    \end{tikzpicture}
    \caption{Distribución asimétrica a la derecha.}
  \end{subfigure}\hfill
  \begin{subfigure}[t]{0.32\linewidth}
    \centering
    \begin{tikzpicture}
      \begin{axis}[mini, xlabel={$z$}, ylabel={$r_i$}, xmin=-3, xmax=3, ymin=-6, ymax=4]
        \addplot[cero, domain=-3:3] {x};
        \addplot[puntos] table[x=q, y=izq] {datos/tema4/qq.dat};
      \end{axis}
    \end{tikzpicture}
    \caption{Distribución asimétrica a la izquierda.}
  \end{subfigure}\hfill
  \begin{subfigure}[t]{0.32\linewidth}
    \centering
    \begin{tikzpicture}
      \begin{axis}[mini, xlabel={$z$}, ylabel={$r_i$}, xmin=-3, xmax=3, ymin=-5, ymax=5]
        \addplot[cero, domain=-3:3] {x};
        \addplot[puntos] table[x=q, y=colas] {datos/tema4/qq.dat};
      \end{axis}
    \end{tikzpicture}
    \caption{Distribución simétrica con colas pesadas.}
  \end{subfigure}
  \caption{QQ-plots de residuos no normales: residuos ordenados ($r_i$) frente a los cuantiles
    esperados de la normal ($z$). La línea discontinua es la recta que seguirían los puntos bajo
    normalidad (datos simulados).}
\end{figure}

Con un tamaño muestral adecuado también es útil comparar las frecuencias observadas con las
esperadas bajo normalidad: aproximadamente el 68\,\% de los residuos deben quedar en la banda
$\pm\sqrt{MSE}$, el 95\,\% en $\pm2\sqrt{MSE}$ y el 99.7\,\% en $\pm3\sqrt{MSE}$.

Para muestras grandes, el supuesto de normalidad no es crucial: \textbf{las inferencias sobre los
parámetros de la regresión $\beta_0$ y $\beta_1$ son robustas ante la falta de normalidad de los
errores}, ya que con tamaños muestrales suficientes la aproximación a la normal es razonable por
el teorema central del límite (\cref{sec:t3-coeficientes}).

\subsection{Independencia}

Si las observaciones proceden de una muestra aleatoria de sujetos de alguna población, en
principio serán independientes. Para detectar correlación entre los errores se representan los
residuos frente al tiempo o frente a algún otro tipo de secuencia (\textbf{gráficos
secuenciales}). La presencia de \textbf{rachas} en los residuos sugiere falta de independencia.

\begin{figure}[H]
  \centering
  \begin{subfigure}[t]{0.48\linewidth}
    \centering
    \begin{tikzpicture}
      \begin{axis}[width=6.2cm, height=4.4cm, xlabel={Orden temporal}, ylabel={$e_i$},
          xmin=0, xmax=14, ymin=-2.6, ymax=2.4, ytick={0}, xtick={1,3,...,13},
          label style={font=\small}, tick label style={font=\footnotesize},
          ylabel style={rotate=-90}]
        \addplot[black, no marks, domain=0:14] {0};
        \addplot[puntos, mark size=1.3pt] coordinates {(1,-2.2) (2,-1.4) (3,-1.6) (4,-1.2)
          (5,-1.0) (6,-0.7) (7,0.3) (8,-0.1) (9,0.9) (10,1.3) (11,1.6) (12,1.3) (13,1.8)};
      \end{axis}
    \end{tikzpicture}
    \caption{Efecto de tendencia.}
  \end{subfigure}\hfill
  \begin{subfigure}[t]{0.48\linewidth}
    \centering
    \begin{tikzpicture}
      \begin{axis}[width=6.2cm, height=4.4cm, xlabel={Orden temporal}, ylabel={$e_i$},
          xmin=0, xmax=15, ymin=-2.2, ymax=2.4, ytick={0}, xtick={1,3,...,13},
          label style={font=\small}, tick label style={font=\footnotesize},
          ylabel style={rotate=-90}]
        \addplot[black, no marks, domain=0:15] {0};
        \addplot[puntos, mark size=1.3pt] coordinates {(1,-1.5) (2,-1.3) (3,-0.6) (4,0.2)
          (5,0.15) (6,1.0) (7,1.8) (8,1.2) (9,0.2) (10,0.15) (11,-1.1) (12,-1.4) (13,-1.2)
          (14,-0.3)};
      \end{axis}
    \end{tikzpicture}
    \caption{Falta de independencia cíclica.}
  \end{subfigure}
  \caption{Gráficos secuenciales de residuos con rachas (datos reconstruidos a partir de las
    diapositivas).}
\end{figure}

En el gráfico de residuos frente a $X$ la nube de puntos puede parecer no aleatoria sin que
exista falta de independencia. Normalmente, en ese caso, el problema es que $X$ no es un buen
predictor de $Y$.

\subsection{Identificación de variables \texorpdfstring{$X$}{X} importantes}
\label{sec:t4-omitidas}

Representar los residuos de un modelo frente a variables $X$ que no se han incluido en él puede
indicar que dicha variable tiene sobre la respuesta un efecto importante que el modelo no
recoge.

\begin{ejemplo}[Masa muscular y edad][ej:t4-masa]
  Se espera que la masa muscular de una persona decrezca con la edad. Con $Y$ = masa muscular y
  $X$ = edad se ajusta el modelo de regresión lineal simple
  \[
    Y_i = \beta_0+\beta_1X_i+\varepsilon_i, \qquad \varepsilon_i\iid\Normal(0,\sigma^2), \quad i=1,\ldots,n,
  \]
  a una muestra de hombres y mujeres, y se representan los residuos estudentizados eliminados
  frente a la edad, distinguiendo el sexo de cada individuo.
  \begin{center}
    \begin{tikzpicture}
      \begin{axis}[width=9cm, height=5.5cm, xlabel={Edad (años)}, ylabel={$t_{(i)}$},
          xmin=40, xmax=80, ymin=-3.5, ymax=3.5, ytick={-3,-2,...,3},
          label style={font=\small}, tick label style={font=\footnotesize},
          ylabel style={rotate=-90},
          legend style={font=\scriptsize, at={(0.5,1.03)}, anchor=south, draw=none,
            legend columns=2}]
        \addplot[only marks, mark=*, mark size=1.1pt, NavyBlue]
          table[x=edad, y=t, restrict expr to domain={\thisrow{sexo}}{0:0}] {datos/tema4/masa.dat};
        \addlegendentry{Mujeres}
        \addplot[only marks, mark=triangle*, mark size=1.4pt, Maroon]
          table[x=edad, y=t, restrict expr to domain={\thisrow{sexo}}{1:1}] {datos/tema4/masa.dat};
        \addlegendentry{Hombres}
        \addplot[gray, no marks, domain=40:80, forget plot] {0};
        \addplot[gray, dashed, no marks, domain=40:80, forget plot] {3};
        \addplot[gray, dashed, no marks, domain=40:80, forget plot] {-3};
      \end{axis}
    \end{tikzpicture}\\
    {\footnotesize (Datos simulados con el patrón del gráfico de las diapositivas.)}
  \end{center}
  Los residuos de los hombres son mayoritariamente positivos y los de las mujeres
  mayoritariamente negativos: el modelo subestima la masa muscular de los hombres y la
  sobrestima en las mujeres. El sexo tiene, por tanto, un efecto importante sobre la masa
  muscular que el modelo no recoge, y convendría incluirlo como variable explicativa.
\end{ejemplo}

\subsection{Contrastes de hipótesis sobre los residuos}

Si a la vista de los gráficos no está claro qué concluir, pueden utilizarse ciertos
\textbf{contrastes de hipótesis}. Estos contrastes \textbf{no reemplazan a las herramientas
gráficas}, ya que dependen mucho del tamaño muestral, sino que se utilizan como complemento. Los
más habituales son:
\begin{itemize}
  \item \emph{Homogeneidad de varianzas}: test de Brown--Forsythe y test de Breusch--Pagan.
  \item \emph{Normalidad}: test de Shapiro--Wilk y test de Kolmogorov--Smirnov, entre otros.
  \item \emph{Independencia}: test de Durbin--Watson.
\end{itemize}

\subsubsection{Contraste de homogeneidad de varianzas: test de Brown--Forsythe}

El más habitual es el test de Brown--Forsythe, una modificación del test de Levene que es robusta
frente a la falta de normalidad.
\begin{itemize}
  \item La muestra se divide en dos grupos: valores bajos y valores altos de $X$.
  \item Si la varianza no es constante, los residuos de un grupo serán más variables que los del
    otro, y las desviaciones absolutas de los residuos alrededor del centro de su grupo serán
    distintas en ambos grupos.
  \item El test se basa en el contraste $t$ de Student para dos muestras aplicado a esas
    desviaciones.
\end{itemize}
Sean $e_{i1}$, $i=1,\ldots,n_1$, y $e_{i2}$, $i=1,\ldots,n_2$, los residuos de los grupos 1 y 2
($n_1+n_2=n$), y $\tilde e_1$, $\tilde e_2$ sus medianas. Se definen
\[
  d_{i1} = |e_{i1}-\tilde e_1|, \qquad d_{i2} = |e_{i2}-\tilde e_2| ,
\]
con medias $\media{d}_1$ y $\media{d}_2$. El estadístico es
\[
  T_{BF} = \frac{\media{d}_1-\media{d}_2}{S\sqrt{\dfrac{1}{n_1}+\dfrac{1}{n_2}}},
  \qquad
  S^2 = \frac{\sum_{i=1}^{n_1}\big(d_{i1}-\media{d}_1\big)^2+\sum_{i=1}^{n_2}\big(d_{i2}-\media{d}_2\big)^2}{n_1+n_2-2},
\]
que, bajo la hipótesis de varianza constante, sigue aproximadamente una $\tdist{n_1+n_2-2}$. Se
rechaza la homogeneidad de varianzas si $|T_{BF}|$ es grande comparado con los cuantiles de esa
distribución.

\subsubsection{Contrastes de normalidad}

En general, estos contrastes sirven para contrastar hipótesis sobre la distribución de una
variable, no solo la normalidad. Se aplican a los residuos; a continuación se describen para
una muestra genérica $x_1,\ldots,x_n$.

\paragraph{Test $\chi^2$ de Pearson.} Es un contraste de bondad de ajuste a una distribución $F$:
\[
  H_0\colon \text{la distribución poblacional es } F
  \quad\text{frente a}\quad
  H_1\colon \text{la distribución poblacional no es } F .
\]
Se basa en el estadístico
\[
  Q = \sum_{j=1}^k \frac{(O_j-E_j)^2}{E_j} \bajoHo \chisq{k-1} \ \text{(aproximadamente)},
\]
donde $k$ es el número de niveles de una variable discreta o de clases de una continua, y $O_j$ y
$E_j$ son las frecuencias observadas y esperadas bajo $F$ en el nivel o clase $j$. Los $k-1$
grados de libertad corresponden a una $F$ completamente especificada; si se estiman $m$
parámetros de $F$ a partir de la muestra (por ejemplo, la varianza de los errores), son $k-1-m$.

\paragraph{Test de Kolmogorov--Smirnov.} Es un contraste de bondad de ajuste a distribuciones
continuas. Se basa en la máxima distancia entre la función de distribución empírica de la
muestra, $F_n$, y la función de distribución teórica $F$:
\[
  D_n = \sup_x\big|F_n(x)-F(x)\big| .
\]
Valores grandes de $D_n$ llevan a rechazar $H_0$.

\begin{definicion}[Función de distribución empírica][def:t4-empirica]
  Sean $x_{(1)}\le\cdots\le x_{(n)}$ los datos ordenados de menor a mayor. La \textbf{función de
  distribución empírica} de la muestra es
  \[
    F_n(x) =
    \begin{cases}
      0 & \text{si } x < x_{(1)}, \\[2pt]
      \dfrac{\operatorname{card}\{j : x_j\le x\}}{n} & \text{si } x_{(i)}\le x<x_{(i+1)},\quad i=1,\ldots,n-1, \\[8pt]
      1 & \text{si } x\ge x_{(n)},
    \end{cases}
  \]
  donde $\operatorname{card}\{j : x_j\le x\}$ es el número de observaciones menores o iguales
  que $x$. Es una función escalonada y no decreciente, que en cada observación da un salto igual
  a la proporción de observaciones de la muestra que toman ese valor.
\end{definicion}

\paragraph{Test de Shapiro--Wilk.} Es útil con muestras pequeñas, ya que es más potente que el de
Kolmogorov--Smirnov. Se basa en un estadístico que cuantifica las diferencias entre las
observaciones y los valores esperados bajo la distribución normal. Toma valores entre 0 y 1, con
valores pequeños para muestras alejadas de la normalidad.

\paragraph{Test de asimetría.}
En las distribuciones simétricas, como la normal, el coeficiente de
asimetría poblacional es 0. Con los momentos centrales muestrales
$m_k=\frac1n\sumi (x_i-\media{x})^k$ y $S_X^2=m_2$, el coeficiente de asimetría se estima por
\[
  CA = \frac{m_3}{S_X^3} = \frac{\sumi (x_i-\media{x})^3}{nS_X^3}
  \ \bajoHo\ \Normal\!\left(0,\ \frac6n\right) \ \text{(aproximadamente)},
\]
o, equivalentemente, estandarizando, $CAS=\sqrt{\frac n6}\,CA\bajoHo\Normal(0,1)$.

\paragraph{Test de apuntamiento.} La distribución normal se caracteriza porque su coeficiente de
apuntamiento es 0. Este se estima por
\[
  CAp = \frac{m_4}{S_X^4}-3 = \frac{\sumi (x_i-\media{x})^4}{nS_X^4}-3
  \ \bajoHo\ \Normal\!\left(0,\ \frac{24}{n}\right) \ \text{(aproximadamente)},
\]
o, equivalentemente, estandarizando, $CApS=\sqrt{\frac{n}{24}}\,CAp\bajoHo\Normal(0,1)$.

En ambos casos se rechaza la normalidad si el estadístico estandarizado es grande en valor
absoluto comparado con los cuantiles de la $\Normal(0,1)$.

% ==============================================================================
%  4.5  Outliers y puntos de influencia
% ==============================================================================
\section{Outliers y puntos de influencia}
\label{sec:t4-outliers}

En algunos problemas, la respuesta observada en unos pocos casos no parece seguir el modelo que
sí se ajusta bien a la gran mayoría de los datos.

\begin{definicion}[Outlier en regresión][def:t4-outlier]
  Un \textbf{outlier en regresión} es un punto que no sigue el patrón lineal general del resto
  de los datos, es decir, que no se ajusta bien al modelo propuesto: su valor ajustado,
  $\est{Y}_i$, está lejos del verdadero valor observado, $Y_i$. Por tanto, estos puntos se
  caracterizan por tener residuos $e_i=Y_i-\est{Y}_i$ altos.
\end{definicion}

\begin{itemize}
  \item El método de mínimos cuadrados es muy sensible a la presencia de outliers, que pueden
    modificar sustancialmente la estimación.
  \item Puede haber outliers tanto en la variable $X$ como en la $Y$.
  \item También puede ocurrir que una observación no sea un outlier en ninguna de las dos
    variables, pero sí en la relación entre $X$ e $Y$ (\cref{sec:t1-descriptivo}).
\end{itemize}
En los modelos de regresión simple es relativamente fácil identificar este tipo de
observaciones mediante diagramas de dispersión.

\begin{figure}[H]
  \centering
  \pgfplotsset{outl/.style={mini, width=3.5cm, height=3.4cm, xmin=0, xmax=15, ymin=0, ymax=15}}
  \def\nube{(1,1) (1,2) (2,2) (2,3) (2,4) (3,3) (3,4) (4,3) (4,4) (4,5) (5,5) (5,6) (6,5)
    (6,6) (6,7) (7,7) (7,8) (8,8) (8,9) (8,10)}
  \begin{subfigure}[t]{0.24\linewidth}
    \centering
    \begin{tikzpicture}
      \begin{axis}[outl]
        \addplot[puntos, mark size=1.2pt] coordinates {\nube};
        \addplot[only marks, mark=*, mark size=1.8pt, colRes] coordinates {(13,5)};
      \end{axis}
    \end{tikzpicture}
    \caption{Valor extremo de $X$.}
  \end{subfigure}\hfill
  \begin{subfigure}[t]{0.24\linewidth}
    \centering
    \begin{tikzpicture}
      \begin{axis}[outl]
        \addplot[puntos, mark size=1.2pt] coordinates {\nube};
        \addplot[only marks, mark=*, mark size=1.8pt, colRes] coordinates {(5,14)};
      \end{axis}
    \end{tikzpicture}
    \caption{Valor extremo de $Y$.}
  \end{subfigure}\hfill
  \begin{subfigure}[t]{0.24\linewidth}
    \centering
    \begin{tikzpicture}
      \begin{axis}[outl]
        \addplot[puntos, mark size=1.2pt] coordinates {\nube};
        \addplot[only marks, mark=*, mark size=1.8pt, colRes] coordinates {(13,14)};
      \end{axis}
    \end{tikzpicture}
    \caption{Extremo en $X$ y en $Y$.}
  \end{subfigure}\hfill
  \begin{subfigure}[t]{0.24\linewidth}
    \centering
    \begin{tikzpicture}
      \begin{axis}[outl]
        \addplot[puntos, mark size=1.2pt] coordinates {\nube};
        \addplot[only marks, mark=*, mark size=1.8pt, colRes] coordinates {(1,10)};
      \end{axis}
    \end{tikzpicture}
    \caption{Alejado de la relación.}
  \end{subfigure}
  \caption{Tipos de outliers (en rojo) en un diagrama de dispersión.}
\end{figure}

Hay, por tanto, dos tipos de atipicidad:
\begin{itemize}
  \item \textbf{outliers univariantes}, atípicos en $X$ o en $Y$;
  \item \textbf{outliers en la relación}, que no tienen por qué ser atípicos en ninguna de las
    dos variables.
\end{itemize}

\begin{definicion}[Punto de influencia][def:t4-influencia]
  Un \textbf{punto de influencia} es un outlier que tiene un efecto muy importante en la recta de
  regresión ajustada.
\end{definicion}

\begin{figure}[H]
  \centering
  \begin{subfigure}[t]{0.45\linewidth}
    \centering
    \begin{tikzpicture}
      \begin{axis}[mini, xmin=-0.1, xmax=2.1, ymin=-0.8, ymax=3.6]
        \addplot[puntos, mark size=1.2pt] table {datos/tema4/infl_sin.dat};
        \addplot[only marks, mark=triangle*, mark size=2.5pt, colRes] coordinates {(2,0)};
        \addplot[recta] table[y={create col/linear regression={y=y}}] {datos/tema4/infl_con.dat};
      \end{axis}
    \end{tikzpicture}
    \caption{Con el punto de influencia (triángulo rojo).}
  \end{subfigure}\hfill
  \begin{subfigure}[t]{0.45\linewidth}
    \centering
    \begin{tikzpicture}
      \begin{axis}[mini, xmin=-0.1, xmax=2.1, ymin=-0.8, ymax=3.6]
        \addplot[puntos, mark size=1.2pt] table {datos/tema4/infl_sin.dat};
        \addplot[recta] table[y={create col/linear regression={y=y}}] {datos/tema4/infl_sin.dat};
      \end{axis}
    \end{tikzpicture}
    \caption{Sin el punto de influencia.}
  \end{subfigure}
  \caption{Recta ajustada con y sin un punto de influencia (datos simulados).}
\end{figure}

Si se encuentran puntos de influencia en un conjunto de datos, hay que tener en cuenta que pueden
ser observaciones que no representan el comportamiento de interés, y comparar las decisiones que
se toman cuando se incluyen y cuando no.

\subsection{Outliers en \texorpdfstring{$X$}{X}: el leverage}
\label{sec:t4-outliers-x}

Los outliers en $X$ pueden identificarse por valores grandes del leverage $h_{ii}$
(\cref{def:t4-leverage}).
\begin{itemize}
  \item En general, $0\le h_{ii}\le1$ y $\sumi h_{ii}=2$ (\cref{prop:t4-prop-leverage}).
  \item Valores grandes del leverage indican que el valor $x_i$ está alejado del centro de todas
    las observaciones: $(x_i-\media{X})^2$ grande.
  \item Se considera que un leverage es grande si es mayor que dos veces la media de los
    leverages:
    \[
      h_{ii} > 2\media{h} = \frac2n\sumi h_{ii} = \frac4n .
    \]
\end{itemize}

Los puntos con leverage grande son \textbf{potenciales} puntos de influencia, es decir, pueden
condicionar de manera importante el ajuste del modelo:
\begin{itemize}
  \item como $\Var(e_i)=\sigma^2(1-h_{ii})$, cuanto mayor es $h_{ii}$ menor es la varianza de
    $e_i$;
  \item puesto que $0\le h_{ii}\le1$, cuanto más próximo a 1 es $h_{ii}$, más próxima a cero es la
    varianza del residuo de la observación $i$-ésima;
  \item por tanto, en las observaciones con $h_{ii}$ grande, $\est{Y}_i$ tiende a estar cerca del
    valor observado $Y_i$;
  \item en el caso extremo e hipotético $h_{ii}=1$, la recta ajustada estaría forzada a pasar por
    el valor observado $(x_i,Y_i)$.
\end{itemize}

\subsection{Outliers en \texorpdfstring{$Y$}{Y}: los residuos}
\label{sec:t4-outliers-y}

Los residuos permiten identificar los outliers en la variable respuesta: valores altos o bajos
indican observaciones que se ajustan mal al modelo. Los residuos regulares, $e_i=Y_i-\est{Y}_i$,
no tienen una escala conocida, por lo que es difícil definir qué es un ``valor extremo''; por
eso siempre resulta más útil utilizar los residuos estudentizados, $r_i$
(\cref{def:t4-estudentizados}), o los estudentizados eliminados, $t_{(i)}$
(\cref{def:t4-eliminado-estudentizado}).
\begin{itemize}
  \item Tanto $r_i$ como $t_{(i)}$ se distribuyen según una $t$ de Student (aproximadamente en el
    caso de $r_i$). Para un tamaño $n$ razonable, un valor de $|r_i|$ o de $|t_{(i)}|$
    \textbf{mayor que 3} indica que la observación es un outlier.
  \item Los valores de $|r_i|$ o $|t_{(i)}|$ \textbf{entre 2 y 3} deben considerarse
    sospechosos.
  \item Muchas observaciones con residuos extremos podrían indicar que el modelo no es adecuado.
  \item Puede ser más útil utilizar $t_{(i)}$ que $r_i$, ya que la observación no interviene en
    el cálculo de su propio residuo.
\end{itemize}

\paragraph{Contraste de outliers.} De manera más formal, para cada observación puede plantearse
el contraste
\[
  H_0\colon \text{la observación } i \text{ no es un outlier}
  \qquad\text{frente a}\qquad
  H_1\colon \text{la observación } i \text{ es un outlier},
\]
donde $H_0$ significa que la observación $i$ sigue el modelo ajustado. Bajo $H_0$ (y con errores
normales), $t_{(i)}\sim\tdist{n-3}$, de modo que se rechaza $H_0$ a nivel $\alpha$ si
$|t_{(i)}|>\tq{n-3}{1-\alpha/2}$. Como se contrastan los $n$ residuos, aparece el problema de las
comparaciones múltiples (\cref{sec:t3-bonferroni}), y se aplica el ajuste de Bonferroni con
$\alpha^*=\alpha/n$.

\begin{resumen}[Identificación de outliers en $Y$]
  Se contrasta cada uno de los $n$ residuos para determinar si corresponde a un outlier: la
  observación $i$ se considera un outlier si
  \[
    |t_{(i)}| > \tq{n-3}{1-\alpha/(2n)} .
  \]
\end{resumen}

\subsection{Medidas de influencia}

Un outlier no siempre es un punto de influencia, ni los outliers son los únicos puntos
sospechosos de influir demasiado en el ajuste.
\begin{itemize}
  \item Un outlier en la variable $Y$ puede ser influyente. Serán potencialmente más influyentes
    los puntos con residuos mucho mayores (o menores) que el resto de residuos de la muestra.
  \item Los puntos con un valor particularmente alto o bajo de la variable $X$ pueden tener una
    gran influencia en la estimación de los dos coeficientes. Estos puntos son potencialmente
    influyentes.
\end{itemize}
Así, con los residuos estudentizados (eliminados) y los leverages se detectan los puntos
potencialmente influyentes, pero es necesario evaluar su \textbf{influencia efectiva}. Para ello
se utilizan las siguientes medidas:
\begin{itemize}
  \item la \textbf{distancia de Cook}, que mide la influencia de una observación en todos los
    valores ajustados;
  \item los \textbf{DFFits}, que miden la influencia de una observación en su propio valor
    ajustado;
  \item los \textbf{DFBeta}, que miden la influencia de una observación en un coeficiente de
    regresión particular.
\end{itemize}
Las tres se basan en comparar el ajuste con todas las observaciones y el ajuste sin la
observación $i$-ésima (\cref{sec:t4-eliminados}).

\paragraph{Distancia de Cook.} Determina la influencia de una observación en todos los valores
predichos; valores grandes indican que la observación tiene mucha influencia. Mide cómo cambian
los estimadores de los coeficientes de regresión cuando se elimina una observación, utilizando la
distancia de Mahalanobis para cuantificar el cambio:
\[
  D_i = \frac{\sum_{j=1}^n\big(\est{Y}_j-\est{Y}_{j(i)}\big)^2}{2\cdot MSE}
  = \frac{e_i^2}{2\cdot MSE}\left(\frac{h_{ii}}{(1-h_{ii})^2}\right).
\]
Su valor se compara con los de la distribución $\Fdist{2}{n-2}$: si $D_i$ se sitúa en torno a la
mediana de la $\Fdist{2}{n-2}$ o por encima de ella, la observación tiene una influencia
importante en el ajuste; si queda por debajo de su percentil 10 o 20, apenas influye.

\paragraph{DFFits.} Determina la influencia de una observación en su propia predicción. Es una
medida estandarizada del cambio que sufre $\est{Y}_i$ al eliminar la observación $i$:
\[
  \DFFits_i = \frac{\est{Y}_i-\est{Y}_{i(i)}}{\sqrt{MSE_{(i)}\cdot h_{ii}}} .
\]
En general, una observación se considera influyente si
\(
  |\DFFits_i| > 2\sqrt{2/n} .
\)

\paragraph{DFBeta.} Hay un DFBeta por cada parámetro y cada observación, y determina la influencia
de cada observación en la estimación de cada parámetro. Es una medida estandarizada del cambio que
sufre $\est{\beta}_j$, $j=0,1$, al eliminar la observación $i$:
\[
  \DFBeta_{ji} = \frac{\est{\beta}_j-\est{\beta}_{j(i)}}{\sqrt{\widehat{\Var}\big(\est{\beta}_{j(i)}\big)}},
\]
donde el error estándar del denominador se calcula con $MSE_{(i)}$ en lugar de $MSE$
(\cref{teo:t3-dist-b0,teo:t3-dist-b1}); por ejemplo, para la pendiente es
$\sqrt{MSE_{(i)}/\SSX}$, con $\SSX$ calculado con todas las observaciones. Se consideran altos los valores absolutos mayores que 1 (en muestras
pequeñas o medianas) o que $2/\sqrt{n}$ (en muestras grandes).

\subsection{Tipos de outliers y puntos de influencia}

\begin{center}
  \small
  \renewcommand{\arraystretch}{1.3}
  \begin{tabular}{l>{\raggedright\arraybackslash}p{4.6cm}>{\raggedright\arraybackslash}p{4.6cm}}
    \toprule
    & \textbf{No outlier en $Y$} & \textbf{Outlier en $Y$} \\
    \midrule
    \textbf{No outlier en $X$} &
      Puntos normales (A): residuos pequeños y leverages pequeños. &
      Outlier, problema leve (B): residuos grandes y leverages pequeños. \\
    \addlinespace
    \textbf{Outlier en $X$} &
      Punto de influencia ``bueno'' (C): leverage grande y medidas de influencia efectiva
      pequeñas. &
      Punto de influencia ``malo'' (D): leverage grande y medidas de influencia efectiva
      grandes. \\
    \bottomrule
  \end{tabular}
\end{center}

\begin{figure}[H]
  \centering
  \begin{tikzpicture}
    \begin{axis}[width=7.5cm, height=5.5cm, ticks=none, axis lines=box, xmin=-0.5, xmax=11,
        ymin=-0.5, ymax=11]
      \addplot[black, no marks, domain=0:10.6] {0.15+0.93*x};
      \addplot[puntos, mark size=1.4pt] coordinates {(1,1.2) (1.5,0.7) (1.5,1.7) (2,2.1)
        (2.3,1.4) (2.7,2.6) (3.2,2.1) (3.6,2.8) (3.7,3.6) (4.1,4.3) (4.4,3.2) (4.6,4.1)
        (5,4.9) (5.3,4.3)};
      \addplot[only marks, mark=*, mark size=1.8pt, colRes] coordinates {(2.2,6.6) (10,9.6) (10,3.8)};
      \node[font=\footnotesize, anchor=north] at (axis cs:2.9,0.9) {A};
      \node[font=\footnotesize, anchor=south east] at (axis cs:2.1,6.7) {B};
      \node[font=\footnotesize, anchor=south east] at (axis cs:9.9,9.7) {C};
      \node[font=\footnotesize, anchor=south east] at (axis cs:9.9,3.9) {D};
    \end{axis}
  \end{tikzpicture}
  \caption{A: puntos normales; B: outlier; C: punto de influencia bueno; D: punto de influencia
    malo.}
\end{figure}

\paragraph{Tratamiento.} Se eliminan los casos que condicionan de manera importante el análisis y
se estudian las causas de la aparición de dichas observaciones.

% ==============================================================================
%  4.6  Soluciones cuando el modelo no es válido
% ==============================================================================
\section{Soluciones cuando el modelo no es válido}
\label{sec:t4-soluciones}

Cuando el modelo no es válido hay, básicamente, dos opciones:
\begin{itemize}
  \item utilizar una metodología más apropiada, que con frecuencia será más complicada;
  \item encontrar una transformación de los datos para la que el modelo de regresión lineal sea
    adecuado.
\end{itemize}
Un modelo ajustado con alguna de las variables transformadas es un modelo nuevo, que requiere
el cumplimiento de las hipótesis: también es necesario validar los modelos transformados.

Según la hipótesis que falle, las soluciones habituales son las siguientes.
\begin{itemize}
  \item \textbf{Relación no lineal}:
    \begin{itemize}
      \item transformar $Y$, $X$ o ambas, de forma que se linealice la relación entre ellas;
      \item introducir términos polinómicos de orden mayor o nuevas variables independientes.
        Por ejemplo, podría ajustarse una función de regresión cuadrática,
        \[
          Y_i = \beta_0+\beta_1X_i+\beta_2X_i^2+\varepsilon_i ,
        \]
        o exponencial,
        \[
          Y_i = \beta_0\beta_1^{X_i}+\varepsilon_i .
        \]
    \end{itemize}
  \item \textbf{Varianza no constante}:
    \begin{itemize}
      \item transformar $Y$;
      \item utilizar otros métodos de ajuste, como el de mínimos cuadrados ponderados.
    \end{itemize}
  \item \textbf{Errores no normales}: transformar $Y$. Con frecuencia la falta de normalidad y la
    heterogeneidad de varianzas se presentan juntas; lo habitual es buscar primero una
    transformación que estabilice la varianza.
\end{itemize}

Muchas veces, transformaciones sencillas de $X$, de $Y$ o de ambas permiten que los modelos
lineales resulten adecuados con las variables transformadas.

\subsection{Transformaciones para linealizar relaciones no lineales}

Supongamos que es razonable considerar que los errores son normales y con varianza constante.
\begin{itemize}
  \item Las transformaciones de $Y$ podrían modificar la forma de su distribución y, por tanto,
    la de los errores, incluidas la normalidad y la homogeneidad de varianzas.
  \item Las transformaciones de $X$ pueden ayudar a linealizar la relación sin afectar a la
    distribución de $Y$.
\end{itemize}
En este caso conviene, por tanto, transformar $X$. Las transformaciones adecuadas dependen de la
forma de la relación observada en el diagrama de dispersión:
\begin{center}
  \renewcommand{\arraystretch}{1.3}
  \begin{tabular}{ll}
    \toprule
    Patrón de la relación & Transformaciones de $X$ \\
    \midrule
    Creciente, cada vez más lentamente (cóncava) & $X'=\log_{10}X$, \quad $X'=\sqrt{X}$ \\
    Creciente, cada vez más rápidamente (convexa) & $X'=X^2$, \quad $X'=\exp(X)$ \\
    Decreciente, cada vez más lentamente (convexa) & $X'=1/X$, \quad $X'=\exp(-X)$ \\
    \bottomrule
  \end{tabular}
\end{center}

\begin{nota}
  Puede ser necesario incluir alguna constante en la transformación para que tenga sentido en
  todos los valores observados. Por ejemplo, $X'=\sqrt{X+k}$, con $k$ tal que $X+k\ge0$.
\end{nota}

\subsection{Transformaciones para la heterogeneidad de varianzas y la falta de normalidad}

Con frecuencia la falta de normalidad y la heterogeneidad de varianzas aparecen juntas. En ese
caso se transforma $Y$ para modificar la forma de su distribución, y a veces también se consigue
linealizar la relación. Una transformación muy utilizada es el \textbf{logaritmo}.

\begin{ejemplo}[Transformación logarítmica][ej:t4-log]
  \begin{itemize}
    \item El histograma de una variable de gastos es muy asimétrico a la derecha: la mayoría de
      los valores se concentran cerca de 0 y unos pocos son muy grandes. El histograma del
      logaritmo de los gastos es, en cambio, aproximadamente simétrico y con forma de campana.
    \item Al representar el salario de los empleados de una empresa frente a su antigüedad (en
      años), la relación es creciente y aproximadamente lineal, pero la dispersión de los
      salarios aumenta con la antigüedad. En el gráfico de residuos frente a valores ajustados se
      observa que la dispersión de los residuos crece con $\est{y}_i$ (forma de embudo). Si se
      ajusta el modelo con $\log(\text{salario})$ como respuesta, la relación con la antigüedad
      sigue siendo lineal y los residuos presentan una dispersión aproximadamente constante.
  \end{itemize}
\end{ejemplo}

Además del logaritmo, hay otras transformaciones útiles de $Y$, como la raíz cuadrada o la
inversa, que forman parte de la familia de transformaciones potencia. Para elegir la
transformación que mejor soluciona el problema se utiliza la transformación de Box--Cox, que
selecciona dentro de esa familia la más adecuada a los datos.

\begin{definicion}[Transformación de Box--Cox][def:t4-boxcox]
  La \textbf{transformación de Box--Cox} es la familia de transformaciones no lineales, para
  $y>0$,
  \[
    g_\lambda(y) =
    \begin{cases}
      \dfrac{y^\lambda-1}{\lambda} & \text{si } \lambda\neq0, \\[8pt]
      \log(y) & \text{si } \lambda=0 .
    \end{cases}
  \]
  Cuando $Y$ puede tomar valores negativos, se utiliza
  \[
    g_\lambda(y) =
    \begin{cases}
      \dfrac{(y+k)^\lambda-1}{\lambda} & \text{si } \lambda\neq0, \\[8pt]
      \log(y+k) & \text{si } \lambda=0 ,
    \end{cases}
  \]
  con $k$ tal que $y+k>0$ para todos los valores observados.
\end{definicion}

\begin{itemize}
  \item La transformación depende del parámetro $\lambda$, que puede elegirse maximizando la
    verosimilitud, suponiendo que para algún $\lambda$ dado las variables transformadas siguen una
    distribución $\Normal(\mu,\sigma^2)$. En la práctica se representa la log-verosimilitud en
    función de $\lambda$ y se toma el valor que la maximiza, junto con un intervalo de confianza
    del 95\,\% para $\lambda$. Por ejemplo, si el máximo se alcanza en $\lambda\approx0$ y el
    intervalo es aproximadamente $(-0.5,\,0.5)$, se elige la transformación logarítmica.
  \item $\lambda=1$ corresponde a la transformación $y'=y-1$, que solo supone un cambio en el
    \emph{intercept} del modelo y no en la forma de la distribución; es decir, equivale a no
    transformar.
\end{itemize}

% ==============================================================================
%  4.7  Validación de la capacidad predictiva
% ==============================================================================
\section{Validación de la capacidad predictiva}
\label{sec:t4-capacidad-predictiva}

La \textbf{capacidad predictiva} de un modelo es su capacidad para realizar predicciones válidas
en individuos distintos de aquellos con los que se ha construido.

\paragraph{Validación externa.} La mejor validación de un modelo es la que utiliza una
\textbf{muestra externa}, es decir, la que comprueba cómo funciona el modelo con datos nuevos que
no se han utilizado para ajustarlo. Con frecuencia esto no es posible.

\paragraph{Validación interna.} Si no se dispone de una muestra externa, se utilizan las
observaciones disponibles. Cuando el conjunto de datos es grande, una solución intermedia es
repartirlo aleatoriamente en dos grupos: una \textbf{muestra de ajuste} o de entrenamiento
(\emph{training data set}), con la que se ajusta el modelo, y una \textbf{muestra de validación}
(\emph{test data set}), con la que se evalúa. Cuando no es posible dividir la muestra de esta
forma, se recurre a métodos iterativos, que se describen a continuación.

\subsection{Validación cruzada dejando uno fuera (LOOCV)}

La validación cruzada dejando uno fuera (\emph{Leave One Out Cross-Validation}, LOOCV) es un
método iterativo:
\begin{itemize}
  \item el conjunto de entrenamiento lo forman todas las observaciones disponibles excepto una;
  \item la observación excluida es la muestra de validación;
  \item el proceso se repite tantas veces como observaciones haya, excluyendo en cada iteración
    una observación distinta, ajustando el modelo con el resto y calculando el error de
    predicción en la observación excluida.
\end{itemize}
En la iteración $i$-ésima, el error de predicción es el residuo eliminado
$e_{(i)}=y_i-\est{Y}_{i(i)}$ (\cref{sec:t4-eliminados}).

\subsection{Validación cruzada en \texorpdfstring{$k$}{k} grupos (\emph{k-fold})}

La validación cruzada en $k$ grupos (\emph{K-Fold Cross-Validation}) también es un método
iterativo:
\begin{itemize}
  \item los datos se dividen aleatoriamente en $k$ grupos de aproximadamente el mismo tamaño;
  \item $k-1$ grupos se emplean para entrenar el modelo, y el grupo restante como muestra de
    validación;
  \item el proceso se repite $k$ veces, utilizando en cada iteración un grupo distinto como
    muestra de validación.
\end{itemize}
El método LOOCV es el caso particular $k=n$.

\subsection{Bootstrap}

\begin{itemize}
  \item Es un método iterativo.
  \item Una \textbf{muestra bootstrap} es una muestra obtenida a partir de la muestra original
    mediante muestreo aleatorio con reposición, y del mismo tamaño que la muestra original.
  \item Las observaciones de la muestra original no seleccionadas en una muestra bootstrap se
    denominan \emph{out-of-bag} (OOB).
  \item En cada iteración se genera una nueva muestra bootstrap, se ajusta el modelo con ella y se
    evalúa con las observaciones \emph{out-of-bag}. La estimación del error de predicción en esa
    iteración es la media de los errores de predicción en las observaciones \emph{out-of-bag}.
  \item El proceso se repite $B$ veces y se calcula la media de las $B$ estimaciones del error de
    predicción.
\end{itemize}

\begin{ejemplo}[Muestras bootstrap]
  A partir de una muestra original $\{x_1,\ldots,x_{10}\}$ se generan tres muestras bootstrap
  (conjuntos de entrenamiento), cada una con sus observaciones \emph{out-of-bag} (conjuntos de
  validación):
  \begin{center}
    \renewcommand{\arraystretch}{1.2}
    \begin{tabular}{lll}
      \toprule
      & Muestra bootstrap (entrenamiento) & \emph{Out-of-bag} (validación) \\
      \midrule
      1 & $x_8,x_6,x_2,x_9,x_5,x_8,x_1,x_4,x_8,x_2$ & $x_3,x_7,x_{10}$ \\
      2 & $x_{10},x_1,x_3,x_5,x_1,x_7,x_4,x_2,x_1,x_8$ & $x_6,x_9$ \\
      3 & $x_6,x_5,x_4,x_1,x_2,x_4,x_2,x_6,x_9,x_2$ & $x_3,x_7,x_8,x_{10}$ \\
      \bottomrule
    \end{tabular}
  \end{center}
\end{ejemplo}


% ==============================================================================
%  TEMA 5 — El modelo de regresión lineal en notación matricial
%  Fuente: diapositivas T5_Matricial.pdf
% ==============================================================================
\chapter{El modelo de regresión lineal en notación matricial}
\label{cap:tema5}

% ==============================================================================
%  5.1  El modelo de regresión lineal en notación matricial
% ==============================================================================
\section{El modelo en notación matricial}
\label{sec:t5-modelo}

El modelo de regresión lineal simple (\cref{def:rls}) se formula en forma escalar como
\[
  Y_i = \beta_0+\beta_1X_i+\varepsilon_i, \qquad \varepsilon_i\iid\Normal(0,\sigma^2), \quad i=1,\ldots,n,
\]
es decir, como un sistema de $n$ ecuaciones, una por observación:
\[
  \begin{aligned}
    Y_1 &= \beta_0+\beta_1X_1+\varepsilon_1, \\
    Y_2 &= \beta_0+\beta_1X_2+\varepsilon_2, \\
        &\;\;\vdots \\
    Y_n &= \beta_0+\beta_1X_n+\varepsilon_n .
  \end{aligned}
\]
En \textbf{forma matricial},
\[
  \underbrace{\begin{pmatrix} Y_1\\ Y_2\\ \vdots\\ Y_n \end{pmatrix}}_{\vect{Y}_{n\times1}}
  =
  \underbrace{\begin{pmatrix} 1 & X_1\\ 1 & X_2\\ \vdots & \vdots\\ 1 & X_n \end{pmatrix}}_{\vect{X}_{n\times2}}
  \underbrace{\begin{pmatrix} \beta_0\\ \beta_1 \end{pmatrix}}_{\vect{\beta}_{2\times1}}
  +
  \underbrace{\begin{pmatrix} \varepsilon_1\\ \varepsilon_2\\ \vdots\\ \varepsilon_n \end{pmatrix}}_{\vect{\varepsilon}_{n\times1}},
\]
o, abreviadamente, $\vect{Y}=\vect{X}\vect{\beta}+\vect{\varepsilon}$.

\begin{definicion}[Modelo de regresión lineal en notación matricial][def:t5-modelo]
  \[
    \vect{Y} = \vect{X}\vect{\beta}+\vect{\varepsilon},
    \qquad \vect{\varepsilon}\sim\Normal(\vect{0},\ \sigma^2\vect{I}_n),
  \]
  donde
  \begin{itemize}
    \item $\vect{Y}_{n\times1}$ es el \textbf{vector de respuestas};
    \item $\vect{\beta}_{2\times1}$ es el \textbf{vector de parámetros};
    \item $\vect{X}_{n\times2}$ es la \textbf{matriz de diseño}, conocida;
    \item $\vect{\varepsilon}_{n\times1}$ es el \textbf{vector de errores}, normal multivariante con
      media $\vect{0}$ y matriz de varianzas-covarianzas
      \[
        \Var(\vect{\varepsilon}) =
        \begin{pmatrix}
          \sigma^2 & 0 & \cdots & 0\\
          0 & \sigma^2 & \cdots & 0\\
          \vdots & \vdots & \ddots & \vdots\\
          0 & 0 & \cdots & \sigma^2
        \end{pmatrix}
        = \sigma^2\vect{I}_n .
      \]
  \end{itemize}
\end{definicion}

La forma diagonal de $\Var(\vect{\varepsilon})$ recoge a la vez la homogeneidad de varianzas
($\Var(\varepsilon_i)=\sigma^2$ para todo $i$) y la independencia de los errores
($\Cov(\varepsilon_i,\varepsilon_j)=0$ para $i\neq j$, que en el caso normal equivale a
independencia). Equivalentemente, $\vect{Y}\sim\Normal(\vect{X}\vect{\beta},\ \sigma^2\vect{I}_n)$,
que es la versión matricial de la \cref{prop:t2-dist-Y}.

El modelo es un caso particular del modelo lineal general $\vect{Y}=Z\vect{\theta}+\vect{U}$ del
Tema 1 (\cref{sec:t1-modelos}), con $Z=\vect{X}$, $\vect{\theta}=\vect{\beta}$ y
$\vect{U}=\vect{\varepsilon}$. Esta formulación permite extender directamente los resultados al
caso de varias variables explicativas, añadiendo columnas a $\vect{X}$ y componentes a
$\vect{\beta}$.

% ==============================================================================
%  5.2  Estimación de los coeficientes de regresión. La matriz hat
% ==============================================================================
\section{Estimación de los coeficientes de regresión}
\label{sec:t5-estimacion}

El modelo se ajusta por el método de mínimos cuadrados (\cref{sec:t2-mc}). Para un valor dado de
$\vect{\beta}$, los errores son $\vect{\varepsilon}=\vect{Y}-\vect{X}\vect{\beta}$, y la suma de
sus cuadrados se escribe como un producto escalar:
\[
  \sumi \varepsilon_i^2 =
  \begin{pmatrix} \varepsilon_1 & \varepsilon_2 & \cdots & \varepsilon_n \end{pmatrix}
  \begin{pmatrix} \varepsilon_1\\ \varepsilon_2\\ \vdots\\ \varepsilon_n \end{pmatrix}
  = \vect{\varepsilon}\T\vect{\varepsilon} .
\]
Se trata, por tanto, de minimizar en $\vect{\beta}$
\[
  \vect{\varepsilon}\T\vect{\varepsilon} = (\vect{Y}-\vect{X}\vect{\beta})\T(\vect{Y}-\vect{X}\vect{\beta}) .
\]
Derivando respecto de $\vect{\beta}$,
\[
  \frac{\partial}{\partial\vect{\beta}}
  \Big((\vect{Y}-\vect{X}\vect{\beta})\T(\vect{Y}-\vect{X}\vect{\beta})\Big)
  = -2\vect{X}\T(\vect{Y}-\vect{X}\vect{\beta}),
\]
e igualando a $\vect{0}$ se obtienen las \textbf{ecuaciones normales}.

\begin{teorema}[Estimador de mínimos cuadrados][teo:t5-mc]
  Las ecuaciones normales son
  \[
    (\XtX)\,\bhat = \vect{X}\T\vect{Y},
  \]
  y, si $\XtX$ es invertible, su solución es
  \[
    \boxed{\ \bhat = \XtXi\vect{X}\T\vect{Y}\ } .
  \]
\end{teorema}

\begin{nota}
  $\XtX$ es invertible si las columnas de $\vect{X}$ son linealmente independientes; en la
  regresión simple, si no todos los $X_i$ son iguales ($\SSX>0$).
\end{nota}

\subsection{Valores ajustados, residuos y matriz hat}

Una vez estimados los coeficientes se obtienen el vector de \textbf{valores ajustados} y el
vector de \textbf{residuos} (\cref{def:t2-ajuste}):
\[
  \est{\vect{Y}}_{n\times1} = \begin{pmatrix} \est{Y}_1\\ \vdots\\ \est{Y}_n \end{pmatrix},
  \qquad
  \vect{e}_{n\times1} = \begin{pmatrix} e_1\\ \vdots\\ e_n \end{pmatrix}.
\]
Los valores ajustados se calculan como
\[
  \est{\vect{Y}} = \vect{X}\bhat = \vect{X}\XtXi\vect{X}\T\vect{Y},
\]
y los residuos $e_i=Y_i-\est{Y}_i$ se escriben matricialmente como
\[
  \vect{e} = \vect{Y}-\est{\vect{Y}} = \vect{Y}-\vect{X}\bhat
  = \vect{Y}-\vect{X}\XtXi\vect{X}\T\vect{Y}
  = \big(\vect{I}-\vect{X}\XtXi\vect{X}\T\big)\vect{Y} .
\]

\begin{definicion}[Matriz hat][def:t5-hat]
  La \textbf{matriz hat} es
  \[
    \vect{H}_{n\times n} = \vect{X}\XtXi\vect{X}\T \in\R^{n\times n} .
  \]
  Con ella,
  \[
    \est{\vect{Y}} = \vect{H}\vect{Y}, \qquad \vect{e} = (\vect{I}-\vect{H})\vect{Y} .
  \]
\end{definicion}

La matriz hat desempeña un papel muy importante en la identificación de puntos influyentes: su
elemento $(i,k)$ es el coeficiente $h_{ik}$ de $Y_k$ en el valor ajustado $\est{Y}_i$
(\cref{prop:t4-hik}), y su diagonal contiene los leverages $h_{ii}$ (\cref{def:t4-leverage}).

\begin{proposicion}[Propiedades de la matriz hat][prop:t5-hat]
  La matriz hat es \textbf{simétrica} e \textbf{idempotente}:
  \[
    \vect{H}\T=\vect{H}, \qquad \vect{H}\vect{H}=\vect{H},
  \]
  y por tanto $\vect{H}\T\vect{H}=\vect{H}\vect{H}=\vect{H}$. La matriz de varianzas-covarianzas de
  los residuos es
  \[
    \Var(\vect{e}) = \sigma^2(\vect{I}-\vect{H}) .
  \]
\end{proposicion}

\subsection{Propiedades de los residuos en forma matricial}

El modelo impone que los errores $\varepsilon_i$ sean independientes, normales y con la misma
varianza. Los errores no son observables; los residuos $e_i$ sí lo son, y se utilizan como
estimadores de los $\varepsilon_i$. Sin embargo, los residuos no son independientes entre sí y sus
varianzas no son iguales (\cref{prop:t4-prop-residuos}):
\begin{itemize}
  \item de $\Var(\vect{e})=\sigma^2(\vect{I}-\vect{H})$, la varianza de $e_i$ es el elemento $i$ de
    la diagonal,
    \[
      \Var(e_i) = \sigma^2(1-h_{ii}),
    \]
    que depende del $i$-ésimo elemento de la diagonal de la matriz hat, el leverage de la
    observación $i$, que se utiliza para detectar observaciones atípicas o potencialmente
    influyentes (\cref{sec:t4-outliers-x});
  \item la covarianza entre $e_i$ y $e_j$ es el elemento $(i,j)$ de $\sigma^2(\vect{I}-\vect{H})$,
    \[
      \Cov(e_i,e_j) = \sigma^2(-h_{ij}) = -\sigma^2h_{ij}, \qquad i\neq j .
    \]
\end{itemize}
Se tendría $\Var(e_i)=\sigma^2$ y $\Cov(e_i,e_j)=0$ solo si $\vect{H}=\vect{0}$, lo que no tiene
sentido, ya que entonces $\est{\vect{Y}}=\vect{H}\vect{Y}=\vect{0}$.

\subsection{El caso de la regresión simple}

\begin{proposicion}[Estimadores y matriz hat en la regresión simple][prop:t5-simple]
  Con $\vect{X}$ la matriz de diseño de la \cref{def:t5-modelo},
  \[
    \XtX = \begin{pmatrix} n & \sumi X_i\\ \sumi X_i & \sumi X_i^2 \end{pmatrix},
    \qquad
    \XtXi = \begin{pmatrix}
      \dfrac1n+\dfrac{\media{X}^2}{\SSX} & -\dfrac{\media{X}}{\SSX}\\[10pt]
      -\dfrac{\media{X}}{\SSX} & \dfrac{1}{\SSX}
    \end{pmatrix},
  \]
  y se obtienen:
  \begin{enumerate}
    \item los estimadores de mínimos cuadrados del Tema 2 (\cref{teo:t2-estimadores-mc}),
      $\est{\beta}_1=SS_{XY}/\SSX$ y $\est{\beta}_0=\media{Y}-\est{\beta}_1\media{X}$;
    \item los elementos de la matriz hat del Tema 4 (\cref{prop:t4-hik}),
      $h_{ik}=\frac1n+\frac{(X_i-\media{X})(X_k-\media{X})}{\SSX}$.
  \end{enumerate}
\end{proposicion}

\begin{demostracion}
  \emph{Inversa de $\XtX$.} Se usan las identidades
  \[
    \sumi X_i = n\media{X},
    \qquad
    \sumi X_i^2 = \SSX + n\media{X}^2 .
  \]
  El determinante de $\XtX$ es
  \begin{align*}
    \det(\XtX) &= n\sumi X_i^2-\left(\sumi X_i\right)^2 \\
    &= n\left(\SSX+n\media{X}^2\right) - n^2\media{X}^2 \\
    &= n\SSX ,
  \end{align*}
  así que
  \begin{align*}
    \XtXi &= \frac{1}{n\SSX}\begin{pmatrix} \sumi X_i^2 & -\sumi X_i\\[4pt] -\sumi X_i & n \end{pmatrix} \\[4pt]
    &= \frac{1}{n\SSX}\begin{pmatrix} \SSX+n\media{X}^2 & -n\media{X}\\ -n\media{X} & n \end{pmatrix} ,
  \end{align*}
  que, dividiendo cada elemento por $n\SSX$, es la expresión del enunciado.
  \begin{enumerate}
    \item Como
      \[
        \vect{X}\T\vect{Y} = \begin{pmatrix} \sumi Y_i\\[4pt] \sumi X_iY_i \end{pmatrix}
        = \begin{pmatrix} n\media{Y}\\[4pt] \sumi X_iY_i \end{pmatrix},
      \]
      el producto $\bhat=\XtXi\vect{X}\T\vect{Y}$ da, para la pendiente,
      \begin{align*}
        \est{\beta}_1 &= -\frac{\media{X}}{\SSX}\,n\media{Y}+\frac{1}{\SSX}\sumi X_iY_i \\
        &= \frac{\sumi X_iY_i-n\media{X}\,\media{Y}}{\SSX} \\
        &= \frac{SS_{XY}}{\SSX} ,
      \end{align*}
      y, para el \emph{intercept},
      \begin{align*}
        \est{\beta}_0 &= \left(\frac1n+\frac{\media{X}^2}{\SSX}\right)n\media{Y}
          - \frac{\media{X}}{\SSX}\sumi X_iY_i \\
        &= \media{Y} + \frac{n\media{X}^2\,\media{Y}}{\SSX} - \frac{\media{X}}{\SSX}\sumi X_iY_i \\
        &= \media{Y}-\media{X}\,\frac{\sumi X_iY_i-n\media{X}\,\media{Y}}{\SSX} \\
        &= \media{Y}-\est{\beta}_1\media{X} .
      \end{align*}
    \item La fila $i$ de $\vect{X}$ es $(1,\ X_i)$ y $h_{ik}$ es el elemento $(i,k)$ de
      $\vect{H}=\vect{X}\XtXi\vect{X}\T$, luego
      \begin{align*}
        h_{ik} &= (1,\ X_i)\XtXi\begin{pmatrix} 1\\ X_k \end{pmatrix} \\
        &= \frac1n+\frac{\media{X}^2}{\SSX}-\frac{\media{X}X_k}{\SSX}
           -\frac{\media{X}X_i}{\SSX}+\frac{X_iX_k}{\SSX} \\
        &= \frac1n+\frac{\media{X}^2-\media{X}X_i-\media{X}X_k+X_iX_k}{\SSX} \\
        &= \frac1n+\frac{(X_i-\media{X})(X_k-\media{X})}{\SSX} .
      \end{align*}
  \end{enumerate}
\end{demostracion}

% ==============================================================================
%  5.3  El análisis de la varianza en notación matricial
% ==============================================================================
\section{El análisis de la varianza}
\label{sec:t5-anova}

Las sumas de cuadrados de la descomposición de la variabilidad (\cref{teo:t3-sst}) se escriben
matricialmente como formas cuadráticas en $\vect{Y}$. Sea $\vect{J}$ la matriz $n\times n$ con
todos sus elementos iguales a 1, de modo que $\vect{Y}\T\vect{J}\vect{Y}=\big(\sum Y_i\big)^2=n^2\media{Y}^2$.

\begin{proposicion}[Sumas de cuadrados en forma matricial][prop:t5-sumas]
  \begin{itemize}
    \item Variabilidad total:
      \[
        SST = \sumi (Y_i-\media{Y})^2 = \sumi Y_i^2-n\media{Y}^2
        = \vect{Y}\T\vect{Y}-\frac1n\vect{Y}\T\vect{J}\vect{Y}
        = \vect{Y}\T\left(\vect{I}-\frac1n\vect{J}\right)\vect{Y} .
      \]
    \item Variabilidad explicada por la regresión:
      \[
        SSR = \sumi (\est{Y}_i-\media{Y})^2
        = \bhat\T\vect{X}\T\vect{Y}-\frac1n\vect{Y}\T\vect{J}\vect{Y}
        = \vect{Y}\T\left(\vect{H}-\frac1n\vect{J}\right)\vect{Y} .
      \]
    \item Variabilidad residual:
      \begin{align*}
        SSE &= \sumi (Y_i-\est{Y}_i)^2 = \sumi e_i^2 = \vect{e}\T\vect{e}
        = (\vect{Y}-\vect{X}\bhat)\T(\vect{Y}-\vect{X}\bhat) \\
        &= \vect{Y}\T\vect{Y}-\bhat\T\vect{X}\T\vect{Y}
        = \vect{Y}\T(\vect{I}-\vect{H})\vect{Y} .
      \end{align*}
  \end{itemize}
  Se cumple que $SST=SSR+SSE$.
\end{proposicion}

En forma matricial, la descomposición es inmediata: las matrices de las tres formas cuadráticas
verifican $\vect{I}-\frac1n\vect{J}=\big(\vect{H}-\frac1n\vect{J}\big)+(\vect{I}-\vect{H})$.

Si $SSR$ es grande, usar el modelo de regresión lineal es muy distinto de usar la media para
predecir la respuesta: el modelo de regresión mejora mucho la calidad de la predicción. Si
$SSR$ es pequeño, usar el modelo de regresión es solo un poco mejor que usar la media. Esta idea
se formaliza con el contraste de la regresión y con el coeficiente de determinación
(\cref{sec:t3-anova,sec:t3-correlacion}), que se calculan igual a partir de estas sumas de
cuadrados; en particular, $MSE=SSE/(n-2)$ es el estimador insesgado de $\sigma^2$.

% ==============================================================================
%  5.4  Inferencias sobre los parámetros
%  5.5  Inferencias sobre la respuesta media y predicción
% ==============================================================================
\section{Inferencias sobre los parámetros de la regresión}
\label{sec:t5-parametros}

La distribución de $\bhat=\XtXi\vect{X}\T\vect{Y}$ se obtiene a partir del siguiente resultado
sobre la normal multivariante.

\begin{teorema}[Transformaciones lineales de la normal multivariante][teo:t5-normal-lineal]
  Sea $\vect{U}\sim\Normal(\vect{\mu},\vect{\Sigma})$ un vector normal multivariante y
  $\vect{V}=\vect{c}+\vect{D}\vect{U}$ una transformación lineal de $\vect{U}$, donde $\vect{c}$ es
  un vector y $\vect{D}$ una matriz. Entonces
  \[
    \vect{V}\sim\Normal\big(\vect{c}+\vect{D}\vect{\mu},\ \vect{D}\vect{\Sigma}\vect{D}\T\big).
  \]
\end{teorema}

\begin{teorema}[Distribución de $\bhat$][teo:t5-dist-beta]
  \[
    \bhat\sim\Normal\big(\vect{\beta},\ \sigma^2\XtXi\big).
  \]
\end{teorema}

\begin{demostracion}
  Se tiene $\vect{Y}\sim\Normal(\vect{X}\vect{\beta},\sigma^2\vect{I})$ y, tomando
  $\vect{D}=\XtXi\vect{X}\T$, $\bhat=\vect{D}\vect{Y}$ es una transformación lineal de $\vect{Y}$.
  Por el \cref{teo:t5-normal-lineal}, $\bhat$ es normal, con
  \begin{align*}
    \E(\bhat) &= \vect{D}\vect{X}\vect{\beta}
      = \underbrace{\XtXi\vect{X}\T\vect{X}}_{\vect{I}}\vect{\beta} = \vect{\beta}, \\
    \Var(\bhat) &= \vect{D}\,\sigma^2\vect{I}\,\vect{D}\T
      = \sigma^2\big(\XtXi\vect{X}\T\big)\big(\XtXi\vect{X}\T\big)\T \\
      &= \sigma^2\XtXi\underbrace{\XtX\big(\XtXi\big)\T}_{\vect{I}} = \sigma^2\XtXi ,
  \end{align*}
  ya que $\XtX$ y su inversa son simétricas: $\big(\XtXi\big)\T=\XtXi$ y
  $\XtX\big(\XtXi\big)\T=\vect{I}$.
\end{demostracion}

Como $MSE$ es un estimador insesgado de $\sigma^2$ (\cref{prop:t2-mse}), la matriz de
varianzas-covarianzas de $\bhat$ se estima por
\[
  \widehat{\Var}(\bhat) = MSE\,\XtXi .
\]

\begin{proposicion}[Varianzas en la regresión simple][prop:t5-var-simple]
  Con la expresión de $\XtXi$ de la \cref{prop:t5-simple},
  \[
    \Var(\bhat) = \sigma^2\begin{pmatrix}
      \dfrac1n+\dfrac{\media{X}^2}{\SSX} & -\dfrac{\media{X}}{\SSX}\\[10pt]
      -\dfrac{\media{X}}{\SSX} & \dfrac{1}{\SSX}
    \end{pmatrix}:
  \]
  la diagonal contiene las varianzas de $\est{\beta}_0$ y $\est{\beta}_1$ del Tema 3
  (\cref{teo:t3-dist-b0,teo:t3-dist-b1}), y fuera de ella aparece su covarianza,
  $\Cov(\est{\beta}_0,\est{\beta}_1)=-\sigma^2\media{X}/\SSX$.
\end{proposicion}

\section{Inferencias sobre la respuesta media y predicción}
\label{sec:t5-respuesta-media}

\subsection{Respuesta media}

Dado el vector $\vect{X}_h=(1,\ x_h)\T$, la respuesta media en $\vect{X}_h$ es
\[
  \E(Y_h) = \E(Y\cond \vect{X}_h) = \vect{X}_h\T\vect{\beta},
\]
y su estimador es
\[
  \est{Y}_h = \vect{X}_h\T\bhat .
\]
$\est{Y}_h$ estima el parámetro $\E(Y_h)=\vect{X}_h\T\vect{\beta}=\beta_0+\beta_1x_h$, la media de la
distribución de $Y$ condicionada a $\vect{X}_h$, y no el valor de una observación concreta
(\cref{sec:t3-respuesta-media}). La estimación tiene sentido siempre que $\vect{X}_h$ tome valores
que podrían haberse dado en la muestra considerada, es decir, dentro del \emph{scope} del modelo:
la región de valores cubierta conjuntamente por todas las variables independientes.

\begin{proposicion}[Distribución de $\est{Y}_h$][prop:t5-Yh]
  Bajo los supuestos del modelo, el estimador $\est{Y}_h=\vect{X}_h\T\bhat$ de $\E(Y_h)$ es:
  \begin{itemize}
    \item normal, puesto que $\bhat$ lo es;
    \item insesgado: $\E(\est{Y}_h)=\E(\vect{X}_h\T\bhat)=\vect{X}_h\T\vect{\beta}=\E(Y_h)$;
    \item de varianza
      $\Var(\est{Y}_h)=\vect{X}_h\T\Var(\bhat)\vect{X}_h=\sigma^2\vect{X}_h\T\XtXi\vect{X}_h$.
  \end{itemize}
  La varianza estimada es
  \[
    \widehat{\Var}(\est{Y}_h) = \vect{X}_h\T\widehat{\Var}(\bhat)\vect{X}_h
    = MSE\cdot\vect{X}_h\T\XtXi\vect{X}_h .
  \]
\end{proposicion}

En la regresión simple, el mismo cálculo que el de $h_{ik}$ en la \cref{prop:t5-simple} da
\[
  \vect{X}_h\T\XtXi\vect{X}_h=\frac1n+\frac{(x_h-\media{X})^2}{\SSX},
\]
que es la expresión del \cref{teo:t3-dist-Yh}.

\subsection{Predicción de una nueva observación}

El objetivo es predecir una nueva observación de $Y$, $\Ynew$, cuando se conocen los valores de
la variable independiente, $\vect{X}_h$.
\begin{itemize}
  \item La nueva observación puede verse como el resultado de una nueva repetición del
    experimento, independiente de los resultados anteriores en los que se basa el análisis de
    regresión.
  \item Se asume que el modelo de regresión subyacente sigue siendo apropiado para la nueva
    observación.
\end{itemize}

Para hacer inferencias sobre $\Ynew$ es necesario conocer un estimador y la distribución del
error que se comete (\cref{sec:t3-prediccion}).
\begin{itemize}
  \item El estimador puntual sigue siendo $\est{Y}_h=\vect{X}_h\T\bhat$.
  \item Al estimar la respuesta media $\E(Y_h)$ se estima la media de la distribución de $Y$; al
    predecir una nueva respuesta $\Ynew$ se quiere predecir un resultado individual de la
    distribución de $Y$, que se situará alrededor de su media.
  \item Hay que tener en cuenta dos fuentes de variación: la variación en la posible ubicación
    del centro de la distribución de $Y$ y la variación dentro de la distribución de
    probabilidad de $Y$.
\end{itemize}
Como $\Ynew$ es independiente de la muestra, y por tanto de $\est{Y}_h$, la varianza del error de
predicción es
\begin{align*}
  \Var(\Ynew-\est{Y}_h) &= \Var(\Ynew)+\Var(\est{Y}_h)
  = \sigma^2+\sigma^2\vect{X}_h\T\XtXi\vect{X}_h \\
  &= \sigma^2\big(1+\vect{X}_h\T\XtXi\vect{X}_h\big),
\end{align*}
con estimador
\[
  \widehat{\Var}(\Ynew-\est{Y}_h) = MSE\big(1+\vect{X}_h\T\XtXi\vect{X}_h\big).
\]

\paragraph{Media de $m$ nuevas observaciones.} Para predecir la media $\Ymnew$ de $m$ nuevas
observaciones en $\vect{X}_h$: la varianza asociada a una observación es $\sigma^2$ y la asociada a
la media de $m$ observaciones es $\sigma^2/m$. La varianza del error de predicción es la suma de la
varianza de esa media y la de la estimación de la respuesta media:
\[
  \Var(\Ymnew-\est{Y}_h) = \sigma^2\left(\frac1m+\vect{X}_h\T\XtXi\vect{X}_h\right).
\]

\subsection{Intervalos de confianza y de predicción}

Sustituyendo $\sigma^2$ por $MSE$, los estadísticos tipificados siguen una $\tdist{n-2}$ y se
obtienen, en forma matricial, los mismos intervalos y contrastes del Tema 3.

\begin{resumen}[Inferencias en forma matricial (nivel $1-\alpha$)]
  Sea $c_{jj}$, $j=0,1$, el elemento $j$-ésimo de la diagonal de $\XtXi$.
  \begin{itemize}
    \item Coeficientes:
      $\displaystyle \beta_j\in\left(\est{\beta}_j\pm\tq{n-2}{1-\alpha/2}\sqrt{MSE\,c_{jj}}\right)$;
      para $H_0\colon\beta_j=b_j$,
      $\displaystyle t^*=\frac{\est{\beta}_j-b_j}{\sqrt{MSE\,c_{jj}}}\bajoHo\tdist{n-2}$.
    \item Respuesta media:
      $\displaystyle \E(Y_h)\in\left(\est{Y}_h\pm\tq{n-2}{1-\alpha/2}
      \sqrt{MSE\cdot\vect{X}_h\T\XtXi\vect{X}_h}\right)$.
    \item Nueva observación:
      $\displaystyle \Ynew\in\left(\est{Y}_h\pm\tq{n-2}{1-\alpha/2}
      \sqrt{MSE\big(1+\vect{X}_h\T\XtXi\vect{X}_h\big)}\right)$.
    \item Media de $m$ nuevas observaciones:
      $\displaystyle \Ymnew\in\left(\est{Y}_h\pm\tq{n-2}{1-\alpha/2}
      \sqrt{MSE\left(\tfrac1m+\vect{X}_h\T\XtXi\vect{X}_h\right)}\right)$.
  \end{itemize}
\end{resumen}


% ==============================================================================
%  TEMA 6 — Introducción al modelo de regresión lineal múltiple
%  Fuente: diapositivas T6_Regresion_multiple.pdf
% ==============================================================================
\chapter{Introducción al modelo de regresión lineal múltiple}
\label{cap:tema6}

% ==============================================================================
%  6.1  El modelo de regresión lineal múltiple
% ==============================================================================
\section{El modelo de regresión lineal múltiple}
\label{sec:t6-modelo}

Se quiere modelizar la esperanza de una variable continua $Y$ cuando se conocen $p-1$
variables explicativas (predictores o regresores) $X_1,X_2,\ldots,X_{p-1}$. El modelo es
``lineal'' porque la esperanza de $Y$ condicionada a las $X$ depende linealmente de
$X_1,X_2,\ldots,X_{p-1}$.

Por ejemplo, con dos predictores ($p=3$) se propone el modelo
\[
  \E(Y\cond X_1,X_2) = \beta_0+\beta_1X_1+\beta_2X_2 ,
\]
donde:
\begin{itemize}
  \item $\beta_0$, $\beta_1$ y $\beta_2$ son constantes desconocidas, llamadas
    \textbf{parámetros del modelo} o \textbf{coeficientes de la regresión};
  \item se estiman a partir de una muestra aleatoria de $n$ observaciones $(X_{i1},X_{i2},Y_i)$,
    $i=1,\ldots,n$, donde $Y_i$ es la respuesta medida en el $i$-ésimo individuo y $X_{i1}$,
    $X_{i2}$ son los valores de las variables explicativas en ese individuo.
\end{itemize}
Por simplicidad, se escribe $\E(Y)$ en lugar de $\E(Y\cond X_1,X_2)$.

\begin{definicion}[Modelo de regresión lineal múltiple][def:t6-rlm]
  \[
    Y_i = \beta_0+\beta_1X_{i1}+\beta_2X_{i2}+\cdots+\beta_{p-1}X_{i,p-1}+\varepsilon_i,
    \qquad \varepsilon_i\iid\Normal(0,\sigma^2), \quad i=1,\ldots,n .
  \]
  La función $\E(Y\cond X_1,X_2,\ldots,X_{p-1})$ se llama \textbf{función de respuesta}.
\end{definicion}

En el modelo de regresión lineal simple, $\E(Y\cond X_1)$ es una recta; con dos variables
explicativas, la función de respuesta $\E(Y\cond X_1,X_2)$ es un plano. A partir de $p>3$ ya no es
posible representarla gráficamente. En regresión múltiple, la función de respuesta también se
denomina \textbf{superficie de regresión} o \textbf{superficie de respuesta}.

\begin{figure}[H]
  \centering
  \begin{tikzpicture}
    \begin{axis}[view={50}{22}, width=10cm, height=8cm, xmin=0, xmax=10, ymin=0, ymax=10,
        zmin=0, zmax=90, xlabel={$X_1$}, ylabel={$X_2$}, zlabel={$Y$}, ticks=none,
        axis lines=left, clip=false, label style={font=\small}]
      \addplot3[surf, shader=flat, fill=NavyBlue!12, faceted color=NavyBlue!60, opacity=0.85,
        domain=0:10, y domain=0:10, samples=5, samples y=5] {10+2*x+5*y};
      % observación (X_i1, X_i2) = (6, 4)
      \draw[dashed, gray] (axis cs:6,4,0) -- (axis cs:6,4,42);
      \draw[dashed, gray] (axis cs:6,0,0) -- (axis cs:6,4,0) -- (axis cs:0,4,0);
      \fill (axis cs:6,4,0) circle (1.6pt);
      \node[font=\scriptsize, anchor=north west] at (axis cs:6,4,0) {$(X_{i1},X_{i2})$};
      \fill[NavyBlue] (axis cs:6,4,42) circle (1.8pt);
      \node[font=\scriptsize, NavyBlue, anchor=west] at (axis cs:6.3,4,40) {$\E(Y_i)$};
      \draw[colRes, thick] (axis cs:6,4,42) -- (axis cs:6,4,56);
      \fill[colRes] (axis cs:6,4,56) circle (1.8pt);
      \node[font=\scriptsize, colRes, anchor=west] at (axis cs:6.3,4,57) {$Y_i$};
      \node[font=\scriptsize, colRes, anchor=east] at (axis cs:5.8,4,49) {$\varepsilon_i$};
      \node[font=\scriptsize, anchor=east] at (axis cs:0,0,10) {$\beta_0=10$};
    \end{axis}
  \end{tikzpicture}
  \caption{Plano de respuesta $\E(Y)=10+2X_1+5X_2$. Para la observación $i$, la respuesta
    $Y_i$ se desvía de su valor esperado $\E(Y_i)$, situado sobre el plano, en el error
    $\varepsilon_i$.}
\end{figure}

\subsection{Hipótesis del modelo}

Los supuestos sobre la componente aleatoria son los mismos que en la regresión simple
(\cref{sec:t2-hipotesis}): $\varepsilon_1,\ldots,\varepsilon_n$ son variables aleatorias
i.i.d.\ con $\varepsilon_i\sim\Normal(0,\sigma^2)$. Es decir:
\begin{itemize}
  \item los $\varepsilon_i$ tienen media cero, $\E(\varepsilon_i)=0$ (\textbf{linealidad});
  \item los $\varepsilon_i$ tienen todos la misma varianza $\sigma^2$, normalmente desconocida,
    que es otro parámetro del modelo a estimar (\textbf{homogeneidad de la varianza});
  \item los $\varepsilon_i$ tienen distribución normal (\textbf{normalidad});
  \item los $\varepsilon_i$ son independientes entre sí e independientes de los predictores
    $x_{i1},x_{i2},\ldots,x_{i,p-1}$ (\textbf{independencia}).
\end{itemize}
Estas hipótesis equivalen a que, conocidos $X_{i1},X_{i2},\ldots,X_{i,p-1}$,
\[
  Y_i\sim\Normal\left(\beta_0+\sum_{j=1}^{p-1}\beta_jX_{ij},\ \sigma^2\right)\ \text{independientes}.
\]

\subsection{Interpretación de los coeficientes}

En el caso con dos predictores, $\E(Y)=\beta_0+\beta_1X_1+\beta_2X_2$:
\begin{itemize}
  \item $\beta_0$ es el \textbf{\emph{intercept}} u ordenada en el origen del plano: la respuesta media
    en el punto $(X_1,X_2)=(0,0)$;
  \item $\beta_1$ indica el cambio en la respuesta media cuando $X_1$ aumenta una unidad,
    manteniendo $X_2$ constante;
  \item análogamente, $\beta_2$ indica el cambio en la respuesta media cuando $X_2$ aumenta una
    unidad, manteniendo $X_1$ constante.
\end{itemize}

\begin{ejemplo}[Interpretación de los coeficientes][ej:t6-interpretacion]
  Sea $\E(Y)=10+2X_1+5X_2$. Si se fija $X_2=3$, se obtiene la recta $\E(Y)=25+2X_1$. Lo mismo
  ocurre para cualquier otro valor de $X_2$: solo cambia la constante de la recta, $10+5X_2$.
  Por tanto, $\beta_1=2$ indica que la respuesta media aumenta 2 unidades por cada incremento
  unitario de $X_1$ cuando $X_2$ se mantiene constante, sea cual sea el valor de $X_2$.
\end{ejemplo}

Por ello, $\beta_1$ y $\beta_2$ se llaman a veces \textbf{coeficientes de regresión parcial}:
reflejan el efecto parcial de una variable explicativa cuando el otro predictor está incluido en el
modelo y se mantiene constante.

\subsection{Consideraciones sobre el modelo}

\begin{itemize}
  \item El término \textbf{``lineal''} se refiere a los \textbf{parámetros}, no a las variables
    explicativas (\cref{def:modelo-lineal}):
    \begin{itemize}
      \item pueden utilizarse modelos de regresión lineales para manejar cualquier función de
        una variable explicativa, por ejemplo $X^2$ o $\log(X)$: $\beta_0+\beta_1X+\beta_2X^2$;
      \item no se utilizan modelos de regresión lineales para funciones no lineales de los
        parámetros, como $\beta_0+\beta_1X_1+\beta_1^2X_2$.
    \end{itemize}
  \item Los \textbf{regresores} pueden ser:
    \begin{itemize}
      \item variables continuas o cualitativas;
      \item variables transformadas;
      \item potencias de una misma variable (\textbf{regresión polinomial}):
        $Y_i=\beta_0+\beta_1X_i+\beta_2X_i^2+\cdots+\beta_{p-1}X_i^{p-1}+\varepsilon_i$;
      \item \textbf{efectos de interacción}: el producto de predictores se incluye como una
        variable adicional, $Y_i=\beta_0+\beta_1X_{i1}+\beta_2X_{i2}+\beta_3X_{i1}X_{i2}+\varepsilon_i$.
    \end{itemize}
  \item \textbf{Multicolinealidad}: si las variables explicativas están muy correlacionadas entre
    sí, es posible que no se obtengan buenos estimadores de los parámetros, ya que tratan de
    explicar la misma parte de la variabilidad de la respuesta (\cref{sec:t6-multicolinealidad}).
\end{itemize}

\subsection{El modelo en notación matricial}
\label{sec:t6-matricial}

El modelo múltiple se escribe como en el Tema 5 (\cref{def:t5-modelo}),
$\vect{Y}=\vect{X}\vect{\beta}+\vect{\varepsilon}$ con
$\vect{\varepsilon}\sim\Normal(\vect{0},\sigma^2\vect{I}_n)$, añadiendo una columna a la matriz de
diseño por cada predictor:
\[
  \vect{X}_{n\times p} =
  \begin{pmatrix}
    1 & X_{11} & X_{12} & \cdots & X_{1,p-1}\\
    1 & X_{21} & X_{22} & \cdots & X_{2,p-1}\\
    \vdots & \vdots & \vdots & & \vdots\\
    1 & X_{n1} & X_{n2} & \cdots & X_{n,p-1}
  \end{pmatrix},
  \qquad
  \vect{\beta}_{p\times1} = \begin{pmatrix} \beta_0\\ \beta_1\\ \vdots\\ \beta_{p-1} \end{pmatrix}.
\]
Todos los resultados del Tema 5 se mantienen con esta $\vect{X}$: el estimador
$\bhat=\XtXi\vect{X}\T\vect{Y}$ (si $\XtX$ es invertible), la matriz hat, las sumas de cuadrados y
$\bhat\sim\Normal\big(\vect{\beta},\sigma^2\XtXi\big)$. Lo único que cambia son los grados de
libertad: como se estiman $p$ coeficientes, $SSE$ tiene $n-p$ grados de libertad, el estimador
insesgado de $\sigma^2$ es
\[
  MSE = \frac{SSE}{n-p},
\]
y los estadísticos $t$ para los coeficientes, la respuesta media y la predicción
(\cref{sec:t5-respuesta-media}) siguen una $\tdist{n-p}$. $SSR$ tiene $p-1$ grados de libertad,
uno por predictor.

\begin{center}
  \renewcommand{\arraystretch}{1.3}
  \begin{tabular}{lccccc}
    \toprule
    Fuente & $SS$ & g.l. & $MS$ & $F$ & $p$-valor \\
    \midrule
    Modelo (regresión) & $SSR$ & $p-1$ & $MSR=SSR/(p-1)$ & $F^*=MSR/MSE$ & $p$ \\
    Error & $SSE$ & $n-p$ & $MSE=SSE/(n-p)$ & & \\
    Total & $SST$ & $n-1$ & & & \\
    \bottomrule
  \end{tabular}
\end{center}
El estadístico $F=MSR/MSE$ contrasta globalmente $H_0\colon\beta_1=\cdots=\beta_{p-1}=0$ frente a
que algún $\beta_k$ sea distinto de 0, con distribución $\Fdist{p-1}{n-p}$ bajo $H_0$; es el
contraste de la regresión (\cref{sec:t3-anova}) con $p-1$ predictores. Cada coeficiente se
contrasta individualmente con el estadístico $t$ de $H_0\colon\beta_k=0$.

\begin{nota}[Sumas de cuadrados de tipo I y de tipo III]
  Las salidas de software descomponen $SSR$ en la aportación de cada predictor, como reducción de
  $SSE$ al incluirlo en el modelo (\cref{sec:t3-ftest-general}):
  \begin{itemize}
    \item las \textbf{sumas de cuadrados de tipo I} (secuenciales) miden la reducción de $SSE$ al
      añadir cada variable al modelo que contiene las anteriores, en el orden en que se
      introducen; suman $SSR$;
    \item las \textbf{sumas de cuadrados de tipo III} (parciales) miden la reducción de $SSE$ al
      añadir cada variable al modelo que contiene todas las demás.
  \end{itemize}
\end{nota}

% ==============================================================================
%  6.2  El coeficiente de determinación
% ==============================================================================
\section{El coeficiente de determinación}
\label{sec:t6-determinacion}

El coeficiente de determinación $R^2=SSR/SST=1-SSE/SST$ (\cref{def:R2}) mide la proporción de
la variabilidad de $Y$ explicada por el modelo también en regresión múltiple. Sin embargo,
\textbf{$R^2$ no es útil para comparar modelos de regresión con distinto número de variables
explicativas}: al añadir variables explicativas al modelo, $R^2$ puede aumentar de forma
artificial, independientemente de que las $X$ añadidas sean importantes o no. La inclusión de una
nueva variable independiente nunca puede hacer que $R^2$ disminuya, aunque la nueva variable no
contribuya en nada a explicar la respuesta.

En su lugar se utiliza el \textbf{coeficiente de determinación múltiple ajustado}, que corrige
$R^2$ penalizando por el número de variables explicativas incluidas en el modelo.

\begin{definicion}[Coeficiente de determinación ajustado][def:t6-R2a]
  \[
    R^2_a = 1-\frac{MSE}{MST} = 1-\left(\frac{n-1}{n-p}\right)\frac{SSE}{SST},
  \]
  donde $MSE=SSE/(n-p)$ y $MST=SST/(n-1)$ (\cref{sec:t3-anova}).
\end{definicion}

\begin{itemize}
  \item Si al añadir una nueva variable independiente al modelo $R^2_a$ aumenta mucho, esa
    variable es importante para predecir $Y$. Si, por el contrario, no aumenta o incluso
    disminuye, es señal de que no contribuye a explicar $Y$ y no debe incluirse en el modelo.
  \item A diferencia de $R^2$, $R^2_a$ \textbf{puede disminuir al incluir una nueva variable en el
    modelo}.
\end{itemize}
En la regresión simple ($p=2$) se obtiene el $R^2$ ajustado que aparece en las salidas del
Tema 3 (\cref{ej:t3-anova-arboles}).

% ==============================================================================
%  6.3  Validación del modelo. Multicolinealidad
% ==============================================================================
\section{Validación del modelo}
\label{sec:t6-validacion}

La validación de las hipótesis sobre los errores, $\varepsilon_1,\ldots,\varepsilon_n$ i.i.d.\ con
$\varepsilon_i\sim\Normal(0,\sigma^2)$, se hace con las mismas herramientas que en la regresión
simple (\cref{sec:t4-analisis}).
\begin{itemize}
  \item \textbf{Herramientas gráficas}:
    \begin{itemize}
      \item linealidad y homogeneidad de la varianza: residuos frente a los valores predichos y
        frente a cada variable explicativa;
      \item normalidad: QQ-plot, residuos frente a los cuantiles normales;
      \item homogeneidad de la varianza: residuos frente a los valores predichos;
      \item independencia: residuos frente al orden de las observaciones.
    \end{itemize}
  \item \textbf{Contrastes de hipótesis}:
    \begin{itemize}
      \item normalidad: Kolmogorov--Smirnov, Shapiro--Wilk, etc.;
      \item homogeneidad de la varianza: Breusch--Pagan o Brown--Forsythe.
    \end{itemize}
  \item \textbf{Si no se verifican} (\cref{sec:t4-soluciones}):
    \begin{itemize}
      \item se transforma $X$ si la relación es no lineal;
      \item se transforma $Y$ si no se cumplen la homogeneidad de varianzas o la normalidad,
        buscando la transformación entre la familia de Box--Cox (\cref{def:t4-boxcox}).
    \end{itemize}
\end{itemize}

\subsection{Multicolinealidad}
\label{sec:t6-multicolinealidad}

\begin{definicion}[Multicolinealidad][def:t6-multicolinealidad]
  Hay \textbf{multicolinealidad} cuando existe correlación entre las variables independientes.
\end{definicion}

\begin{itemize}
  \item Supone un problema cuando esa correlación es muy alta (por encima de 0.8--0.9).
  \item La \textbf{situación ideal} es que no haya correlación entre los predictores: así no se
    solapan en la variabilidad de la respuesta que pueden explicar. En ese caso, las sumas de
    cuadrados y las pendientes estimadas no se modifican al añadir variables al modelo.
\end{itemize}
La multicolinealidad puede suponer un problema para:
\begin{itemize}
  \item evaluar la importancia relativa de cada uno de los predictores;
  \item determinar la magnitud del efecto que tiene un predictor sobre la variable respuesta.
\end{itemize}

\begin{ejemplo}[Predictores no correlacionados][ej:t6-productividad]
  Se evalúa la productividad ($Y$) de ocho empleados según el tamaño del turno en el que trabajan
  ($X_1$) y los incentivos que reciben ($X_2$). Los dos predictores están incorrelados: cada
  combinación de $X_1$ y $X_2$ aparece con dos empleados.
  \begin{center}
    \renewcommand{\arraystretch}{1.15}
    \begin{tabular}{ccc}
      \toprule
      Productividad ($Y$) & Tamaño del turno ($X_1$) & Incentivos ($X_2$) \\
      \midrule
      42, 39 & 4 & 2 \\
      48, 51 & 4 & 3 \\
      49, 53 & 6 & 2 \\
      61, 60 & 6 & 3 \\
      \bottomrule
    \end{tabular}
  \end{center}

  \textbf{Modelo con las dos variables independientes.}
  \begin{center}
    \small
    \renewcommand{\arraystretch}{1.15}
    \begin{tabular}{lrrrrr}
      \toprule
      Fuente & g.l. & $SS$ & $MS$ & $F$ & $p$-valor \\
      \midrule
      Modelo & 2 & 402.250 & 201.125 & 57 & 0.0004 \\
      Error & 5 & 17.625 & 3.525 & & \\
      Total & 7 & 419.875 & & & \\
      \bottomrule
    \end{tabular}
    \qquad
    \begin{tabular}{lrrr}
      \toprule
      Fuente & g.l. & $SS$ tipo I & $SS$ tipo III \\
      \midrule
      Tamaño & 1 & 231.125 & 231.125 \\
      Incentivos & 1 & 171.125 & 171.125 \\
      \bottomrule
    \end{tabular}

    \medskip
    \begin{tabular}{lrrrr}
      \toprule
      Parámetro & Estimación & Error estándar & $t$ & $p$-valor \\
      \midrule
      \emph{Intercept} & 0.375 & 4.74 & 0.08 & 0.9400 \\
      Tamaño & 5.375 & 0.66 & 8.10 & 0.0005 \\
      Incentivos & 9.250 & 1.33 & 6.97 & 0.0009 \\
      \bottomrule
    \end{tabular}
  \end{center}
  (Las medias cuadráticas de las sumas de cuadrados de tipo I y III coinciden con ellas, al tener
  cada una un grado de libertad.)

  \textbf{Modelo solo con el tamaño del turno.}
  \begin{center}
    \small
    \renewcommand{\arraystretch}{1.15}
    \begin{tabular}{lrrrrr}
      \toprule
      Fuente & g.l. & $SS$ & $MS$ & $F$ & $p$-valor \\
      \midrule
      Modelo & 1 & 231.125 & 231.125 & 7.4 & 0.035 \\
      Error & 6 & 188.750 & 31.458 & & \\
      Total & 7 & 419.875 & & & \\
      \bottomrule
    \end{tabular}
    \qquad
    \begin{tabular}{lrrrr}
      \toprule
      Parámetro & Estimación & Error estándar & $t$ & $p$-valor \\
      \midrule
      \emph{Intercept} & 23.5 & 10.1 & 2.32 & 0.0591 \\
      Tamaño & 5.375 & 1.98 & 2.71 & 0.0351 \\
      \bottomrule
    \end{tabular}
  \end{center}
  Como los predictores no están correlacionados, la suma de cuadrados asociada al tamaño del turno
  es la misma, 231.125, en el modelo simple y en el múltiple (y en este coinciden las de tipo I y
  tipo III), y la pendiente estimada del tamaño es también la misma en ambos modelos, 5.375.
\end{ejemplo}

\begin{ejemplo}[Predictores perfectamente correlacionados][ej:t6-colineales]
  Se tienen cuatro observaciones cuyos predictores están exactamente sobre la recta
  $X_2=5+0.5X_1$:
  \begin{center}
    \renewcommand{\arraystretch}{1.15}
    \begin{tabular}{ccc}
      \toprule
      $X_1$ & $X_2$ & $Y$ \\
      \midrule
      2 & 6 & 25 \\
      8 & 9 & 81 \\
      6 & 8 & 60 \\
      10 & 10 & 113 \\
      \bottomrule
    \end{tabular}
  \end{center}
  Al ajustar el modelo con $X_1$ y $X_2$, SAS avisa de que el modelo no es de rango completo: las
  soluciones de mínimos cuadrados para los parámetros no son únicas, algunos estadísticos pueden
  ser engañosos, y las estimaciones marcadas con g.l.\ 0 o B están sesgadas. Además, indica que el
  parámetro de $X_2$ se ha fijado en 0 porque la variable es combinación lineal de otras,
  $X_2=5\cdot\textit{Intercept}+0.5\,X_1$. La salida es:
  \begin{center}
    \small
    \renewcommand{\arraystretch}{1.15}
    \begin{tabular}{lrrrrr}
      \toprule
      Fuente & g.l. & $SS$ & $MS$ & $F$ & $p$-valor \\
      \midrule
      Modelo & 1 & 4007.150 & 4007.150 & 91.49 & 0.0108 \\
      Error & 2 & 87.600 & 43.800 & & \\
      Total corregido & 3 & 4094.750 & & & \\
      \bottomrule
    \end{tabular}
    \qquad
    \begin{tabular}{lr}
      \toprule
      Raíz MSE & 6.61816 \\
      Media dependiente & 69.75000 \\
      Coef Var & 9.48840 \\
      R-cuadrado & 0.9786 \\
      R-Sq Ajust & 0.9679 \\
      \bottomrule
    \end{tabular}

    \medskip
    \begin{tabular}{lcrrrr}
      \toprule
      Variable & g.l. & Estimación & Error estándar & $t$ & $p$-valor \\
      \midrule
      \emph{Intercept} & B & 0.20000 & 7.98892 & 0.03 & 0.9823 \\
      $X_1$ & B & 10.70000 & 1.11867 & 9.56 & 0.0108 \\
      $X_2$ & 0 & 0 & -- & -- & -- \\
      \bottomrule
    \end{tabular}
  \end{center}
  \begin{itemize}
    \item Como las observaciones de $(X_1,X_2)$ están exactamente sobre una recta, hay infinitos
      planos de regresión que serían igual de buenos: la solución no es única.
    \item La columna de $X_2$ en la matriz de diseño es combinación lineal de las otras dos, por
      lo que $\XtX$ no es invertible (\cref{teo:t5-mc}) y no pueden estimarse las pendientes de
      todas las variables implicadas.
    \item SAS elige ajustar el modelo solo con $X_1$: la salida coincide con la de la regresión
      simple de $Y$ sobre $X_1$, con $\est{Y}=0.2+10.7X_1$ y $R^2_a=1-43.8/(4094.75/3)=0.9679$.
  \end{itemize}
\end{ejemplo}

\paragraph{Situación intermedia.} Normalmente, la correlación entre las $X$ no es ni 0 ni 1.
\begin{itemize}
  \item Cuando esa correlación no es muy grande, puede estimarse un buen modelo, útil para hacer
    predicciones dentro de su \emph{scope}.
  \item Cuando la correlación supone un problema grave de multicolinealidad, se observa lo
    siguiente:
    \begin{itemize}
      \item correlaciones muy altas entre pares de predictores;
      \item estimadores de los coeficientes de regresión con error estándar muy alto;
      \item el contraste global $F$ de la tabla ANOVA es significativo y $R^2$ es alto, pero los
        contrastes $t$ asociados a las variables independientes individuales no lo son (no se
        rechaza $H_0\colon\beta_k=0$);
      \item las sumas de cuadrados de tipo I son muy diferentes de las de tipo III;
      \item hay grandes cambios en los coeficientes de regresión estimados cuando se añaden o se
        eliminan variables del modelo.
    \end{itemize}
\end{itemize}

\paragraph{Factor de inflación de la varianza.} Para detectar la multicolinealidad se utiliza el
\textbf{VIF} (\emph{Variance Inflation Factor}), relacionado con la varianza estimada de los
coeficientes de regresión, que se ve inflada cuando hay multicolinealidad.

\begin{definicion}[VIF y tolerancia][def:t6-vif]
  El \textbf{factor de inflación de la varianza} del predictor $X_k$ es
  \[
    VIF_k = \frac{1}{1-R_k^2},
  \]
  donde $R_k^2$ es el coeficiente de determinación del modelo de regresión de $X_k$ sobre el
  resto de predictores. La \textbf{tolerancia} es
  \[
    T_k = 1-R_k^2 = \frac{1}{VIF_k}.
  \]
\end{definicion}

Si, por ejemplo, $R_k^2\ge0.9$, $X_k$ se predice muy bien a partir del resto de variables
independientes. Esto corresponde a $VIF_k\ge10$, que indica un problema muy serio de
multicolinealidad. Equivalentemente, tolerancias menores que 0.1 son indicadoras de
multicolinealidad.

% ==============================================================================
%  6.4  Criterios de selección de un modelo
% ==============================================================================
\section{Criterios de selección de un modelo}
\label{sec:t6-seleccion}

Al construir un modelo de regresión múltiple hay que decidir:
\begin{itemize}
  \item \textbf{cuántas variables explicativas} utilizar. Generalmente son más convenientes los
    conjuntos pequeños, aunque incluyendo muchas variables independientes se consigue explicar
    una mayor parte de la variabilidad de la respuesta;
  \item dado un número de variables a incluir, \textbf{cuáles} de entre todas las $X$ disponibles
    deben formar parte del modelo.
\end{itemize}
Hay muchos criterios para comparar modelos entre sí, basados en distintas cuestiones. Uno de ellos
es la \textbf{bondad del ajuste}, que se mide con $R^2$ cuando los modelos tienen el mismo número
de predictores y con $R^2_a$ cuando no.

\subsection{Criterios basados en la bondad del ajuste}

\paragraph{Mismo número de predictores: criterio de $R^2$.} Se define
\[
  R^2_p = \frac{SSR_p}{SST} = 1-\frac{SSE_p}{SST},
\]
donde el subíndice $p$ indica el número de parámetros del modelo.
\begin{itemize}
  \item Son mejores los modelos con $R^2_p$ mayor.
  \item Como $SST$ no depende del modelo, elegir el modelo con máximo $R^2_p$ equivale a elegir,
    entre los modelos con $p-1$ predictores, el de menor $SSE$.
\end{itemize}

\paragraph{Distinto número de predictores: criterio de $R^2_a$.} $R^2_a$ (\cref{def:t6-R2a})
penaliza el valor de $R^2$ teniendo en cuenta el número de predictores incluidos en el modelo:
\[
  R^2_a = 1-\left(\frac{n-1}{n-p}\right)\frac{SSE}{SST} = 1-\frac{MSE}{MST}.
\]
\begin{itemize}
  \item Cuando $p$ aumenta, la cantidad que se resta tiene un denominador menor, por lo que, si la
    inclusión de las nuevas variables no se compensa suficientemente con una disminución de $SSE$,
    $R^2_a$ puede disminuir a medida que se incorporan variables.
  \item Como $MST$ no depende del modelo, elegir el modelo con máximo $R^2_a$ equivale a elegir el
    modelo con menor $MSE$.
\end{itemize}

\paragraph{Búsqueda exhaustiva.} Aplicar estos criterios a todos los modelos posibles supone una
búsqueda exhaustiva. El número de modelos que pueden construirse crece exponencialmente con el
número de predictores disponibles (con $k$ candidatos hay $2^k$ subconjuntos posibles). Si hay
pocos predictores pueden compararse todos los modelos, pero con muchos esa comparación es
inabordable, y se necesitan procedimientos semiautomáticos que permitan elegir el subconjunto de
variables independientes con el que construir los mejores modelos.

\subsection{Algoritmos automáticos de selección de variables}

Se basan en los $p$-valores asociados a los contrastes $H_0\colon\beta_k=0$ frente a
$H_1\colon\beta_k\neq0$ de cada coeficiente.

\paragraph{Selección hacia adelante (\emph{forward}).}
\begin{enumerate}
  \item Se fija un $p$-valor de entrada, \emph{p-to-enter}.
  \item Se ajustan los modelos individuales, cada uno con una sola de las variables
    independientes.
  \item Se identifica el predictor con menor $p$-valor en el contraste $H_0\colon\beta_k=0$ frente
    a $H_1\colon\beta_k\neq0$.
    \begin{enumerate}[label=3.\arabic*.]
      \item Si $p$-valor $<$ \emph{p-to-enter}, se añade la variable correspondiente, se reajusta el
        modelo con ella y se vuelve al paso 3, ahora con los modelos que resultan de añadir al
        modelo actual cada una de las variables restantes.
      \item Si $p$-valor $\ge$ \emph{p-to-enter}, no entran más variables y el proceso concluye.
    \end{enumerate}
\end{enumerate}

\paragraph{Eliminación hacia atrás (\emph{backward}).}
\begin{enumerate}
  \item Se fija un $p$-valor de salida, \emph{p-to-remove}.
  \item Se ajusta el modelo de regresión completo, con todas las variables independientes.
  \item Se identifica el predictor con mayor $p$-valor en el contraste $H_0\colon\beta_k=0$ frente
    a $H_1\colon\beta_k\neq0$.
    \begin{enumerate}[label=3.\arabic*.]
      \item Si $p$-valor $>$ \emph{p-to-remove}, se elimina la variable correspondiente, se
        reajusta el modelo sin ella y se vuelve al paso 3.
      \item Si $p$-valor $\le$ \emph{p-to-remove}, no salen más variables y el proceso concluye.
    \end{enumerate}
\end{enumerate}

\paragraph{Regresión paso a paso (\emph{stepwise}).} Combina las ideas de los métodos
\emph{forward} y \emph{backward}.
\begin{enumerate}
  \item Se fijan dos $p$-valores tales que \emph{p-to-remove} $\ge$ \emph{p-to-enter}.
  \item Se ajustan los modelos individuales, cada uno con una sola de las variables
    independientes.
  \item Se identifica el predictor con menor $p$-valor en el contraste $H_0\colon\beta_k=0$ frente
    a $H_1\colon\beta_k\neq0$.
    \begin{enumerate}[label=3.\arabic*.]
      \item Si $p$-valor $<$ \emph{p-to-enter}, se añade la variable correspondiente y se reajusta
        el modelo con ella (paso \emph{forward}).
      \item En el modelo ajustado en 3.1 se comprueba si alguna de las variables que ya estaban en
        el modelo debe salir (paso \emph{backward}):
        \begin{enumerate}[label=3.2.\arabic*.]
          \item se elige la variable con mayor $p$-valor;
          \item si $p$-valor $>$ \emph{p-to-remove}, se elimina la variable correspondiente, se
            reajusta el modelo sin ella y se vuelve a 3.2;
          \item si $p$-valor $\le$ \emph{p-to-remove}, no salen más variables y se vuelve al
            paso 3, ahora con los modelos que resultan de añadir al modelo actual cada una de las
            variables restantes.
        \end{enumerate}
      \item El proceso concluye cuando no entran ni salen más variables.
    \end{enumerate}
\end{enumerate}




% ------------------------------------------------------------------------------
%  BIBLIOGRAFÍA / MATERIAL DE REFERENCIA
% ------------------------------------------------------------------------------
% ==============================================================================
%  Bibliografía / material de referencia
% ==============================================================================
\begin{thebibliography}{9}
  \addcontentsline{toc}{chapter}{Bibliografía}
  \bibitem{diapositivas}
    Itziar Fernández Martínez.
    \textit{Diapositivas Regresión y ANOVA}. Universidad de Valladolid, 2026.
  \bibitem{kutner}
    M.~H.~Kutner, C.~J.~Nachtsheim, J.~Neter, W.~Li.
    \textit{Applied Linear Statistical Models}. McGraw-Hill.
  \bibitem{weisberg}
    S.~Weisberg. \textit{Applied Linear Regression}. Wiley.
  \bibitem{freund}
    R.~J.~Freund, W.~J.~Wilson, P.~Sa.
    \textit{Regression Analysis: Statistical Modeling of a Response Variable}. Academic Press.
\end{thebibliography}


\end{document}
